\documentclass[fs9,seriesRus]{korfarm-base}
\usepackage{korfarm-russell}

% ============================================================
% passage-note.tex (v11)
% 지문 박스 + 첫 페이지 우측 메모란 (동적 높이)
% 정책: PAGE_BREAK_POLICY.md §3, §4 참조
%
% 변경 (v11):
%  · 지문 박스 폭 확대: enlarge right by=-3.7cm  (메모 3.4cm + gap 3mm)
%  · 메모란 동적 높이: frame.north east ~ frame.south east 사이 자동 채움
%  · 메모 줄선: clip 영역으로 박스 내부만 채움 (박스 높이 모르고도 작동)
%  · 문단 들여쓰기 활성화: tcolorbox 내부 \parindent=1em, \noindent 제거
%
% 사용:
%   \kfPassageNote{제목}{지문 본문}      → 메모란 포함
%   \kfPassageNoteNT{지문 본문}          → 메모란 포함, 제목 없음
%   \kfPassageFull{제목}{본문}           → 메모란 없음 (활동 내 보기)
% ============================================================

\makeatletter

% 메모란 규격 (cm 고정)
\newlength{\kf@memoW}
\setlength{\kf@memoW}{3.4cm}
\newlength{\kf@memoGap}
\setlength{\kf@memoGap}{3mm}

% --- 메인: 제목 있는 버전 ---
\newcommand{\kf@passagenote@with}[2]{%
  \par\ifkfPassageNewPage\ifdim\pagetotal>0.4\textheight\clearpage\fi\fi\smallskip\needspace{6\baselineskip}%
  \begin{tcolorbox}[
    enhanced, breakable,
    width=\dimexpr\linewidth-4cm\relax,
    colback=kfPassageBg, colframe=kfPrimary!70,
    boxrule=0.6pt, arc=3pt,
    left=\kfBoxPad, right=\kfBoxPad, top=\dimexpr\kfBoxPad+8pt\relax, bottom=\kfBoxPad,
    title={\small\bfseries\color{kfPrimary} #1},
    coltitle=kfPrimary,
    colbacktitle=kfPrimaryLight!60,
    attach boxed title to top left={yshift=-2.4mm, xshift=8pt},
    boxed title style={colframe=kfPrimary!50, boxrule=0.5pt, arc=2pt},
    parbox=false,
    overlay unbroken and first={%
      \begin{scope}
        \draw[kfRule, line width=0.4pt, fill=kfNoteBg, rounded corners=2pt]
          ([xshift=\kf@memoGap]frame.north east) rectangle
          ([xshift=\kf@memoGap+\kf@memoW]frame.south east);
        \node[anchor=north east, font=\footnotesize\bfseries\color{kfMute}, inner sep=3pt]
          at ([xshift=\kf@memoGap+\kf@memoW, yshift=-2pt]frame.north east) {메모};
        \node[anchor=north east, fill=kfPrimary!15, text=kfPrimary,
              font=\scriptsize\bfseries, rounded corners=2pt,
              inner xsep=4pt, inner ysep=2pt]
          at ([xshift=\kf@memoGap+\kf@memoW-3pt, yshift=-19pt]frame.north east) {\kf@memocat};
        \begin{scope}
          \path[clip]
            ([xshift=\kf@memoGap]frame.north east) rectangle
            ([xshift=\kf@memoGap+\kf@memoW]frame.south east);
          \foreach \i in {1,...,40}{%
            \draw[kfRule, line width=0.25pt]
              ([xshift=\kf@memoGap+2mm, yshift=-7mm - \i*5mm]frame.north east) --
              ([xshift=\kf@memoGap+\kf@memoW-2mm, yshift=-7mm - \i*5mm]frame.north east);%
          }
        \end{scope}
      \end{scope}
    }
  ]
    \setlength{\parindent}{1em}%
    \hspace*{1em}\ignorespaces#2%
  \end{tcolorbox}\par\smallskip
}

% --- 제목 없는 버전 ---
\newcommand{\kf@passagenote@notitle}[1]{%
  \par\ifkfPassageNewPage\ifdim\pagetotal>0.4\textheight\clearpage\fi\fi\smallskip\needspace{6\baselineskip}%
  \begin{tcolorbox}[
    enhanced, breakable,
    width=\dimexpr\linewidth-4cm\relax,
    colback=kfPassageBg, colframe=kfPrimary!70,
    boxrule=0.6pt, arc=3pt,
    left=\kfBoxPad, right=\kfBoxPad, top=\dimexpr\kfBoxPad+8pt\relax, bottom=\kfBoxPad,
    parbox=false,
    overlay unbroken and first={%
      \begin{scope}
        \draw[kfRule, line width=0.4pt, fill=kfNoteBg, rounded corners=2pt]
          ([xshift=\kf@memoGap]frame.north east) rectangle
          ([xshift=\kf@memoGap+\kf@memoW]frame.south east);
        \node[anchor=north east, font=\footnotesize\bfseries\color{kfMute}, inner sep=3pt]
          at ([xshift=\kf@memoGap+\kf@memoW, yshift=-2pt]frame.north east) {메모};
        \node[anchor=north east, fill=kfPrimary!15, text=kfPrimary,
              font=\scriptsize\bfseries, rounded corners=2pt,
              inner xsep=4pt, inner ysep=2pt]
          at ([xshift=\kf@memoGap+\kf@memoW-3pt, yshift=-19pt]frame.north east) {\kf@memocat};
        \begin{scope}
          \path[clip]
            ([xshift=\kf@memoGap]frame.north east) rectangle
            ([xshift=\kf@memoGap+\kf@memoW]frame.south east);
          \foreach \i in {1,...,40}{%
            \draw[kfRule, line width=0.25pt]
              ([xshift=\kf@memoGap+2mm, yshift=-7mm - \i*5mm]frame.north east) --
              ([xshift=\kf@memoGap+\kf@memoW-2mm, yshift=-7mm - \i*5mm]frame.north east);%
          }
        \end{scope}
      \end{scope}
    }
  ]
    \setlength{\parindent}{1em}%
    \hspace*{1em}\ignorespaces#1%
  \end{tcolorbox}\par\smallskip
}

\newcommand{\kfPassageNote}[2]{%
  \def\kf@tmptitle{#1}%
  \ifx\kf@tmptitle\@empty
    \kf@passagenote@notitle{#2}%
  \else
    \kf@passagenote@with{#1}{#2}%
  \fi
}
\newcommand{\kfPassageNoteNT}[1]{\kf@passagenote@notitle{#1}}
\makeatother

% ============================================================
% 풀폭 지문 박스 (메모란 없음) — 활동 내 보기 박스
% PAGE_BREAK_POLICY.md §3
% ============================================================
\makeatletter
\newcommand{\kf@passagefull@with}[2]{%
  \par\ifkfPassageNewPage\ifdim\pagetotal>0.4\textheight\clearpage\fi\fi\smallskip\needspace{6\baselineskip}%
  \begin{tcolorbox}[
    enhanced, breakable,
    colback=kfPassageBg, colframe=kfPrimary!70,
    boxrule=0.6pt, arc=3pt,
    left=\kfBoxPad, right=\kfBoxPad, top=\dimexpr\kfBoxPad+8pt\relax, bottom=\kfBoxPad,
    title={\small\bfseries\color{kfPrimary} #1},
    coltitle=kfPrimary,
    colbacktitle=kfPrimaryLight!60,
    attach boxed title to top left={yshift=-2.4mm, xshift=8pt},
    boxed title style={colframe=kfPrimary!50, boxrule=0.5pt, arc=2pt},
    parbox=false
  ]
    \setlength{\parindent}{1em}%
    \hspace*{1em}\ignorespaces#2%
  \end{tcolorbox}\par\smallskip
}
\newcommand{\kf@passagefull@notitle}[1]{%
  \par\ifkfPassageNewPage\ifdim\pagetotal>0.4\textheight\clearpage\fi\fi\smallskip\needspace{6\baselineskip}%
  \begin{tcolorbox}[
    enhanced, breakable,
    colback=kfPassageBg, colframe=kfPrimary!70,
    boxrule=0.6pt, arc=3pt,
    left=\kfBoxPad, right=\kfBoxPad, top=\dimexpr\kfBoxPad+8pt\relax, bottom=\kfBoxPad,
    parbox=false
  ]
    \setlength{\parindent}{1em}%
    \hspace*{1em}\ignorespaces#1%
  \end{tcolorbox}\par\smallskip
}
\newcommand{\kfPassageFull}[2]{%
  \def\kf@tmptitle{#1}%
  \ifx\kf@tmptitle\@empty
    \kf@passagefull@notitle{#2}%
  \else
    \kf@passagefull@with{#1}{#2}%
  \fi
}
\newcommand{\kfPassageFullNT}[1]{\kf@passagefull@notitle{#1}}
\makeatother

% ============================================================
% 작은 문장 박스 (문장 독해용) — PAGE_BREAK_POLICY.md §2.5
% ============================================================
\newcommand{\kfSentenceBox}[2]{%
  \par\smallskip\needspace{3\baselineskip}%
  \noindent\begin{tcolorbox}[
    enhanced, breakable=false,
    colback=kfPassageBg, colframe=kfMute,
    boxrule=0.4pt, arc=2pt,
    left=8pt, right=6pt, top=4pt, bottom=4pt,
    title={\footnotesize\bfseries\color{kfPrimary} #1},
    coltitle=kfPrimary, colbacktitle=kfPaper,
    attach boxed title to top left={yshift=-1.6mm, xshift=4pt},
    boxed title style={colframe=kfMute, boxrule=0.3pt, arc=1pt}
  ]%
    \setstretch{1.3}#2%
  \end{tcolorbox}\par\smallskip%
}

% ============================================================
% 지문 본문 아래 안내/정리 박스 (v16 신규)
% 사용: \kfPassageHint{한 줄 정리}{본문 안내 텍스트}
% ============================================================
\newcommand{\kfPassageHint}[2]{%
  \par\smallskip\needspace{4\baselineskip}%
  \begin{tcolorbox}[
    enhanced, breakable=false,
    colback=kfPrimary!5, colframe=kfPrimary!40,
    boxrule=0.4pt, arc=2pt,
    left=8pt, right=6pt, top=4pt, bottom=4pt,
    title={\footnotesize\bfseries\color{kfPrimary} #1},
    coltitle=kfPrimary, colbacktitle=kfPaper,
    attach boxed title to top left={yshift=-1.5mm, xshift=4pt},
    boxed title style={colframe=kfPrimary!40, boxrule=0.3pt, arc=1pt}
  ]%
    \small\setstretch{1.3}#2%
  \end{tcolorbox}\par\smallskip%
}

% ============================================================
% 환경 버전 (호환성) — 메모란 포함
% ============================================================
\makeatletter
\newenvironment{kfPassageNoteEnv}[1]%
  {\par\needspace{6\baselineskip}%
   \def\kf@psrc{#1}%
   \begin{tcolorbox}[
     enhanced, breakable,
     width=\dimexpr\linewidth-4cm\relax,
     colback=kfPassageBg, colframe=kfPrimary!70,
     boxrule=0.6pt, arc=3pt,
     left=\kfBoxPad, right=\kfBoxPad, top=\dimexpr\kfBoxPad+8pt\relax, bottom=\kfBoxPad,
     title={\small\bfseries\color{kfPrimary} \kf@psrc},
     coltitle=kfPrimary, colbacktitle=kfPrimaryLight!60,
     attach boxed title to top left={yshift=-2.4mm, xshift=8pt},
     boxed title style={colframe=kfPrimary!50, boxrule=0.5pt, arc=2pt},
     parbox=false,
     overlay unbroken and first={%
       \begin{scope}
         \draw[kfRule, line width=0.4pt, fill=kfNoteBg, rounded corners=2pt]
           ([xshift=\kf@memoGap]frame.north east) rectangle
           ([xshift=\kf@memoGap+\kf@memoW]frame.south east);
         \node[anchor=north east, font=\footnotesize\bfseries\color{kfMute}, inner sep=3pt]
           at ([xshift=\kf@memoGap+\kf@memoW, yshift=-2pt]frame.north east) {메모};
         \begin{scope}
           \path[clip]
             ([xshift=\kf@memoGap]frame.north east) rectangle
             ([xshift=\kf@memoGap+\kf@memoW]frame.south east);
           \foreach \i in {1,...,40}{%
             \draw[kfRule, line width=0.25pt]
               ([xshift=\kf@memoGap+2mm, yshift=-7mm - \i*5mm]frame.north east) --
               ([xshift=\kf@memoGap+\kf@memoW-2mm, yshift=-7mm - \i*5mm]frame.north east);%
           }
         \end{scope}
       \end{scope}
     }
   ]%
   \setlength{\parindent}{1em}%
  }%
  {\end{tcolorbox}\par\smallskip}
\makeatother

% ============================================================
% condition-box.tex
% <조건> 박스 컴포넌트 (tcolorbox)
% 사용:
%   \begin{kfCondition}
%     \item 첫째 조건
%     \item 둘째 조건
%   \end{kfCondition}
% ============================================================

\newenvironment{kfCondition}%
  {\par\smallskip\needspace{4\baselineskip}%
   \begin{tcolorbox}[
     enhanced, breakable=false,
     colback=kfPrimaryLight!40, colframe=kfPrimary,
     boxrule=0.5pt, arc=2pt,
     left=\kfBoxPad, right=\kfBoxPad, top=\kfBoxPad, bottom=\kfBoxPad,
     title={\small\bfseries\color{kfPrimary} \,조건\,},
     coltitle=kfPrimary, colbacktitle=kfPaper,
     attach boxed title to top left={yshift=-2mm, xshift=6pt},
     boxed title style={colframe=kfPrimary, boxrule=0.4pt, arc=1pt}
   ]%
   \begin{enumerate}[leftmargin=16pt, itemsep=1pt, topsep=1pt,
                    label=\textbf{\arabic*.}]}%
  {\end{enumerate}\end{tcolorbox}\par\smallskip}

% ============================================================
% evidence-box.tex
% <보기> 박스 컴포넌트
% 사용:
%   \begin{kfEvidence}
%     보기 본문 ...
%   \end{kfEvidence}
% ============================================================

\newenvironment{kfEvidence}%
  {\par\smallskip\needspace{4\baselineskip}%
   \begin{tcolorbox}[
     enhanced, breakable=false,
     colback=kfPaper, colframe=kfMute,
     boxrule=0.5pt, arc=2pt,
     left=\kfBoxPad, right=\kfBoxPad, top=\kfBoxPad, bottom=\kfBoxPad,
     title={\small\bfseries\color{kfInk} \,보기\,},
     coltitle=kfInk, colbacktitle=kfPrimaryLight!50,
     attach boxed title to top left={yshift=-2mm, xshift=6pt},
     boxed title style={colframe=kfMute, boxrule=0.4pt, arc=1pt}
   ]%
   \leavevmode\mbox{}\ignorespaces%
  }%
  {\end{tcolorbox}\par\smallskip}

% ============================================================
% choice-list.tex
% 5선지 (① ② ③ ④ ⑤) 자동 들여쓰기 컴포넌트
% 사용:
%   \begin{kfChoices}
%     \kfChoice 첫째 선택지
%     \kfChoice 둘째 선택지
%     \kfChoice 셋째 선택지
%     \kfChoice 넷째 선택지
%     \kfChoice 다섯째 선택지
%   \end{kfChoices}
% ============================================================

\newcounter{kfChoiceCnt}
\newcommand{\kfChoiceMark}[1]{%
  \ifcase#1\relax%
  \or ①\or ②\or ③\or ④\or ⑤\or ⑥\or ⑦\or ⑧\or ⑨%
  \fi
}

% 환경: 자동 번호
\newenvironment{kfChoices}%
  {\setcounter{kfChoiceCnt}{0}%
   \par\smallskip%
   \begin{list}{}{%
     \setlength{\leftmargin}{20pt}%
     \setlength{\itemindent}{0pt}%
     \setlength{\labelwidth}{18pt}%
     \setlength{\labelsep}{6pt}%
     \setlength{\itemsep}{2pt}%
     \setlength{\topsep}{1pt}%
     \setlength{\parsep}{0pt}%
     \renewcommand{\makelabel}[1]{##1\hss}%
   }}%
  {\end{list}\par\smallskip}

\newcommand{\kfChoice}{%
  \stepcounter{kfChoiceCnt}%
  \item[\textbf{\kfChoiceMark{\value{kfChoiceCnt}}}]\ignorespaces%
}

% --- 한 줄 5선지 (짧은 선택지용) ---
\newcommand{\kfChoicesInline}[5]{%
  \par\smallskip%
  \noindent\textbf{①}~#1\quad\textbf{②}~#2\quad\textbf{③}~#3\quad\textbf{④}~#4\quad\textbf{⑤}~#5%
  \par\smallskip%
}

% ============================================================
% answer-note.tex
% 답안 작성란 (노트 스타일, 줄친 영역)
% 사용:
%   \kfAnswerLines{5}        % 5줄 답안란
%   \kfAnswerNote{서술형 답안란}{8} % 라벨+8줄
% ============================================================

% 단일 줄 (필요 시 인라인)
\newcommand{\kfAnswerLine}{%
  \par\noindent\textcolor{kfRule}{\rule{\linewidth}{0.4pt}}\par
}

% N줄 답안란 (간단)
\newcommand{\kfAnswerLines}[1]{%
  \par\smallskip%
  \begin{minipage}{\linewidth}%
  \foreach \i in {1,...,#1}{%
    \textcolor{kfRule}{\rule{\linewidth}{0.4pt}}\\[6pt]%
  }%
  \end{minipage}%
  \par\smallskip%
}

% 라벨이 있는 노트형 답안란
\newcommand{\kfAnswerNote}[2]{%
  \par\smallskip\needspace{#2\baselineskip}%
  \begin{tcolorbox}[
    enhanced, breakable=false,
    colback=kfPaper, colframe=kfRule,
    boxrule=0.4pt, arc=2pt,
    left=\kfBoxPad, right=\kfBoxPad, top=\kfBoxPad, bottom=\kfBoxPad,
    title={\footnotesize\bfseries\color{kfMute} #1},
    coltitle=kfMute, colbacktitle=kfPaper,
    attach boxed title to top left={yshift=-1.5mm, xshift=6pt},
    boxed title style={colframe=kfRule, boxrule=0.3pt, arc=1pt}
  ]%
  \foreach \i in {1,...,#2}{%
    \textcolor{kfRule}{\rule{\linewidth}{0.3pt}}\\[6pt]%
  }%
  \end{tcolorbox}\par\smallskip%
}

% 짧은 빈칸 (단답형)
\newcommand{\kfBlank}[1][3cm]{\underline{\hspace{#1}}}
% 넓은 빈칸 (서술 한 줄)
\newcommand{\kfBlankWide}[1][6cm]{\underline{\hspace{#1}}}
% 괄호 빈칸
\newcommand{\kfBlankParen}{(\,\hspace{1cm}\,)}
% O/X 빈칸
\newcommand{\kfBlankOX}{(\,\hspace{1.5cm}\,)}

% ============================================================
% question-block.tex
% 문제 블록 (번호 배지 + 발문 + 보기/조건 + 선택지 + 답안란)
% 모든 구성을 하나의 minipage로 묶어 단/페이지 단절 금지
%
% 사용:
%   \begin{kfQBlock}{1}{객관식}
%     \kfQStem{다음 글에 대한 설명으로 가장 적절한 것은?}
%     \begin{kfEvidence} ... \end{kfEvidence}    % 선택
%     \begin{kfCondition} ... \end{kfCondition}  % 선택
%     \begin{kfChoices}
%       \kfChoice ...
%     \end{kfChoices}
%     \kfAnswerLines{2}                            % 선택
%   \end{kfQBlock}
% ============================================================

% 무단절 minipage로 묶고, 발문/보기/조건/선택지/답안란을 자유롭게 배치
% 사용자 피드백 #4: [객관식]/[서술형] 등 텍스트 라벨 제거 — 번호 배지만 표시
% #2(유형 라벨)는 호환성 인자로만 받고 출력하지 않음.
\newenvironment{kfQBlock}[2]%
  {\par\smallskip\needspace{6\baselineskip}%
   \noindent\begin{minipage}{\linewidth}%
   \samepage%
   \par\noindent
   \kfQNum{#1}\hspace{4pt}%
   \begin{minipage}[t]{\dimexpr\linewidth-30pt\relax}%
  }%
  {\end{minipage}\end{minipage}\par\smallskip}

% 문제 발문
\newcommand{\kfQStem}[1]{%
  \par\noindent#1\par\vspace{2pt}%
}

% 짧은 소문항 라벨 (1), (2), (3)
\newcounter{kfSubQ}
\newcommand{\kfSubQReset}{\setcounter{kfSubQ}{0}}
\newcommand{\kfSubQ}{%
  \stepcounter{kfSubQ}%
  \par\noindent\textbf{(\arabic{kfSubQ})}~\ignorespaces%
}

% ============================================================
% activity-table.tex
% 활동: 정리표 / 내용 확인표 (마크다운 표 → tabularx 변환)
% 사용:
%   % 일반 tabular (직접 컬럼 명시)
%   \begin{kfActTable}{|l|X|X|}
%     \kfTHead{항목} & \kfTHead{설명} & \kfTHead{학생 작성} \\\hline
%     ① & ... & \kfTBlank \\\hline
%   \end{kfActTable}
%
%   % adjustbox 자동 축소 (정해진 폭으로 강제)
%   \kfActTableWrap{
%     ... (\begin{tabular}...\end{tabular} 코드) ...
%   }
% ============================================================

% 사용 시 본문 마지막 행은 \\\hline 으로 끝나야 함.
% v17.1: 흰색 mask로 Canva 배경 본문 침범 차단
\newenvironment{kfActTable}[1]%
  {\par\smallskip\needspace{4\baselineskip}%
   \renewcommand{\arraystretch}{1.3}%
   \arrayrulecolor{kfRule}%
   \begin{tcolorbox}[blanker, colback=white, left=2pt, right=2pt, top=2pt, bottom=2pt]%
   \noindent\begin{adjustbox}{max width=\linewidth}%
   \begin{tabular}{#1}%
   \hline}%
  {\end{tabular}%
   \end{adjustbox}%
   \end{tcolorbox}%
   \arrayrulecolor{black}%
   \par\smallskip}

% 헤더 행 헬퍼
\newcommand{\kfTHead}[1]{\textbf{\color{kfPrimary} #1}}
\newcommand{\kfTHeadRow}[1]{\rowcolor{kfPrimaryLight!50}\textbf{\color{kfPrimary} #1}}

% 빈 셀 (학생 작성용)
\newcommand{\kfTBlank}{\rule[-1.6em]{0pt}{3.4em}}
\newcommand{\kfTBlankSm}{\rule[-1.0em]{0pt}{2.4em}}

% adjustbox 강제 축소 래퍼 — Python에서 변환된 raw tabular 블록을 감쌀 때 사용
\newcommand{\kfActTableWrap}[1]{%
  \par\smallskip\needspace{4\baselineskip}%
  \begin{tcolorbox}[blanker, colback=white, left=2pt, right=2pt, top=2pt, bottom=2pt]%
  \noindent\begin{adjustbox}{max width=\linewidth, center}%
  #1%
  \end{adjustbox}%
  \end{tcolorbox}\par\smallskip%
}

% ============================================================
% activity-vocab.tex
% 어휘 학습 표 (3열: 번호 - 뜻 - 빈칸)
% 사용:
%   \begin{kfVocab}
%     \kfVocabItem{1}{눈으로 보고 마음으로 깨달음}
%     \kfVocabItem{2}{사물의 이치를 따져 헤아림}
%   \end{kfVocab}
% ============================================================

% adjustbox로 폭 강제 + 일반 tabular + p{} 열 사용
\newenvironment{kfVocab}%
  {\par\smallskip\needspace{4\baselineskip}%
   \renewcommand{\arraystretch}{1.35}%
   \arrayrulecolor{kfRule}%
   \noindent\begin{adjustbox}{max width=\linewidth, center}%
   \begin{tabular}{|>{\centering\arraybackslash}m{1.0cm}|m{8.5cm}|>{\centering\arraybackslash}m{4.0cm}|}%
   \hline
   \rowcolor{kfPrimaryLight!50}%
   \multicolumn{1}{|c|}{\textbf{\color{kfPrimary}번호}} &
   \multicolumn{1}{c|}{\textbf{\color{kfPrimary}뜻풀이}} &
   \multicolumn{1}{c|}{\textbf{\color{kfPrimary}어휘 (정답 작성)}} \\\hline
   \ignorespaces}%
  {\end{tabular}%
   \end{adjustbox}%
   \arrayrulecolor{black}%
   \par\smallskip}

\newcommand{\kfVocabItem}[2]{%
  \textbf{#1} & #2 & \rule[-1.0em]{0pt}{2.4em}\\\hline%
}

% ============================================================
% activity-blank.tex
% 빈칸 채우기 활동
% 사용:
%   \begin{kfBlankList}
%     \kfBlankItem 첫 번째 문장에서 (\kfBlank)은(는) ...
%     \kfBlankItem 둘째 문장에서 ...
%   \end{kfBlankList}
% ============================================================

\newenvironment{kfBlankList}%
  {\par\smallskip%
   \begin{enumerate}[leftmargin=20pt, itemsep=4pt, topsep=2pt,
                    label=\textbf{\arabic*.}, labelsep=6pt]}%
  {\end{enumerate}\par\smallskip}

\newcommand{\kfBlankItem}{\item}

% 텍스트 안에 박스형 빈칸 (단어 1개 채우기)
\newcommand{\kfBlankBox}[1][2.4cm]{%
  \tikz[baseline=-0.5ex]{%
    \draw[kfRule, line width=0.4pt] (0,0) -- ++(#1,0);%
  }%
}

% ============================================================
% activity-ox.tex (v15)
% O/X 표기 활동 — 그래픽 디자인
% 사용:
%   \begin{kfOXList}
%     \kfOXItem 글쓴이는 자연을 예찬하고 있다.\kfOX
%     \kfOXItem 1연과 4연은 수미상관 구조이다.\kfOX
%   \end{kfOXList}
%
% \kfOX = 우측에 [O] [X] 두 박스 (학생이 하나 동그라미)
% ============================================================

\newenvironment{kfOXList}%
  {\par\smallskip%
   \begin{enumerate}[leftmargin=20pt, itemsep=5pt, topsep=2pt,
                    label=\textbf{\arabic*.}, labelsep=6pt]}%
  {\end{enumerate}\par\smallskip}

\newcommand{\kfOXItem}{\item}

% v16: O/X 그래픽 박스 키움 (18×14 → 24×20pt) + 영역색 강조
\newcommand{\kfOX}{%
  \hfill\nolinebreak
  \tikz[baseline=-0.9ex]{%
    \node[draw=kfPrimary, fill=kfPrimary!5, line width=0.6pt, rounded corners=3pt,
          minimum width=24pt, minimum height=20pt, inner sep=0pt,
          font=\normalsize\bfseries\color{kfPrimary}] (ob) {O};%
    \node[draw=kfMute, fill=kfMute!5, line width=0.6pt, rounded corners=3pt,
          minimum width=24pt, minimum height=20pt, inner sep=0pt,
          font=\normalsize\bfseries\color{kfMute},
          right=4pt of ob] (xb) {X};%
  }%
}

% ============================================================
% activity-connect.tex
% 좌우 선 긋기 활동 (2단 양식 + 중앙 여백)
% 사용:
%   \begin{kfConnect}
%     \kfPair{사르트르}{실존은 본질에 앞선다}
%     \kfPair{니체}{초인 사상}
%     \kfPair{하이데거}{현존재}
%   \end{kfConnect}
% ============================================================

\newenvironment{kfConnect}%
  {\par\smallskip%
   \renewcommand{\arraystretch}{1.6}%
   \begin{tabular}{@{}>{\raggedleft\arraybackslash}m{4.5cm}
                     @{\hspace{2.5cm}}
                     >{\raggedright\arraybackslash}m{6.5cm}@{}}}%
  {\end{tabular}\par\smallskip}

\newcommand{\kfPair}[2]{%
  \tikz[baseline]{\node[draw=kfRule, fill=kfPrimaryLight!30, rounded corners=2pt,
        inner sep=4pt, font=\small]{#1};} \tikz[baseline]{\node[circle, draw=kfMute, fill=kfPaper, inner sep=1.5pt]{};}
  &
  \tikz[baseline]{\node[circle, draw=kfMute, fill=kfPaper, inner sep=1.5pt]{};} \tikz[baseline]{\node[draw=kfRule, fill=kfNoteBg, rounded corners=2pt,
        inner sep=4pt, font=\small]{#2};} \\
}

% ============================================================
% activity-writing.tex
% 글쓰기 활동 (큰 노트 영역)
% 사용:
%   \kfWriting{독후감}{12}     % 라벨 + 12줄 큰 노트
%   \kfWritingPrompt{프롬프트 텍스트}
% ============================================================

\newcommand{\kfWritingPrompt}[1]{%
  \par\smallskip%
  \begin{tcolorbox}[
    enhanced, breakable=false,
    colback=kfPrimaryLight!30, colframe=kfPrimary,
    boxrule=0.4pt, arc=3pt,
    left=\kfBoxPad, right=\kfBoxPad, top=\kfBoxPad, bottom=\kfBoxPad,
  ]%
  \small\textbf{\color{kfPrimary}글쓰기 안내}\\[2pt]%
  #1
  \end{tcolorbox}\par\smallskip%
}

\newcommand{\kfWriting}[2]{%
  \par\smallskip\needspace{\numexpr#2+2\relax\baselineskip}%
  \begin{tcolorbox}[
    enhanced, breakable=false,
    colback=kfPaper, colframe=kfRule,
    boxrule=0.4pt, arc=2pt,
    left=\kfBoxPad, right=\kfBoxPad, top=\kfBoxPad, bottom=\kfBoxPad,
    title={\footnotesize\bfseries\color{kfMute} #1},
    coltitle=kfMute, colbacktitle=kfPaper,
    attach boxed title to top left={yshift=-1.5mm, xshift=6pt},
    boxed title style={colframe=kfRule, boxrule=0.3pt, arc=1pt}
  ]%
  \foreach \i in {1,...,#2}{%
    \textcolor{kfRule}{\rule{\linewidth}{0.3pt}}\\[7pt]%
  }%
  \end{tcolorbox}\par\smallskip%
}

% ============================================================
% activity-structure.tex
% 구조도 (트리 + 빈칸)
% 사용:
%   \begin{kfStructure}
%     \kfNode{0}{서론}{문제 제기}
%     \kfNode{1}{본론}{(\,\quad\,)}
%     \kfNode{1}{본론}{근거 2}
%     \kfNode{0}{결론}{주장 강화}
%   \end{kfStructure}
% ============================================================

\newenvironment{kfStructure}%
  {\par\smallskip\needspace{6\baselineskip}%
   \begin{tcolorbox}[
     enhanced, breakable=false,
     colback=kfPaper, colframe=kfRule,
     boxrule=0.4pt, arc=2pt,
     left=\kfBoxPad, right=\kfBoxPad, top=\kfBoxPad, bottom=\kfBoxPad,
     title={\footnotesize\bfseries\color{kfMute} 글의 구조},
     coltitle=kfMute, colbacktitle=kfPrimaryLight!40,
     attach boxed title to top left={yshift=-1.5mm, xshift=6pt},
     boxed title style={colframe=kfRule, boxrule=0.3pt, arc=1pt}
   ]%
   \begin{tabbing}
     \hspace{1.4cm}\=\hspace{1.4cm}\=\hspace{8cm}\kill}%
  {\end{tabbing}\end{tcolorbox}\par\smallskip}

% 노드: #1=들여쓰기 단계 (0~3), #2=라벨, #3=내용
\newcommand{\kfNode}[3]{%
  \ifcase#1\relax%
    \>\>\textbf{\color{kfPrimary} ▶ #2}\quad #3\\
  \or
    \>\>\quad\textbf{\color{kfMute} ◦ #2}\quad #3\\
  \or
    \>\>\quad\quad\textbf{\color{kfMute} - #2}\quad #3\\
  \or
    \>\>\quad\quad\quad\textbf{\color{kfMute} \textperiodcentered\ #2}\quad #3\\
  \fi
}

% 빈칸 노드
\newcommand{\kfNodeBlank}[2]{\kfNode{#1}{#2}{(\hspace{2.5cm})}}

% ============================================================
% activity-sentence-reading.tex
% 문장 독해 활동 — 문장 박스 1개 + 그에 묶인 문제 N개를 한 minipage로
% 사용:
%   \begin{kfSentenceUnit}{1}{"바람은 내 귀에 속삭이며..."}
%     \kfSentQ{(1)}{서술어 '속삭이며'의 주어는?}
%     \begin{kfChoices}
%       \kfChoice 바람은
%       ...
%     \end{kfChoices}
%   \end{kfSentenceUnit}
% ============================================================

% kfSentenceBox 매크로는 passage-note.tex에 정의됨

\newenvironment{kfSentenceUnit}[2]%
  {\par\smallskip\needspace{6\baselineskip}%
   \noindent\begin{minipage}{\linewidth}%
   \samepage%
   \kfSentenceBox{문장 #1}{#2}%
   \par\smallskip%
  }%
  {\end{minipage}\par\smallskip}

% 같은 문장에 묶인 소문항 (문장 안내 + 발문)
\newcommand{\kfSentQ}[2]{%
  \par\noindent\textbf{#1}~#2\par\smallskip%
}

% 헤더 없이 문장만 있는 묶음 (sentence_ref가 없는 일반 문항)
\newenvironment{kfSentencelessUnit}%
  {\par\smallskip\noindent\begin{minipage}{\linewidth}\samepage}%
  {\end{minipage}\par\smallskip}

% ============================================================
% chapter-cover.tex (v6)
% 챕터 진입 페이지 (표지) — 하단에 영역 목차 추가
% 사용:
%   \kfChapterCover{1}{근대성과 자아}{Chapter 1}{이번 챕터에서는 ...}
%   \begin{kfChapterAreaList}
%     \kfChapterAreaItem{어휘}{kfAreaVocab}{핵심 어휘 12개}
%     ...
%   \end{kfChapterAreaList}
%
%   영역 목차가 없을 때는 \kfChapterCoverEnd 호출
% ============================================================

\newcommand{\kfChapterCover}[4]{%
  \clearpage%
  \thispagestyle{kfBlank}%
  \kfMarkChapter{Chapter #1. #2}%
  \kfChapterCoverBg%
  % v18.5: 챕터 표지 우상단에 로고
  \begin{tikzpicture}[remember picture, overlay]
    \node[anchor=north east, inner sep=0pt] at ($(current page.north east)+(-18mm,-18mm)$) {\kfLogoMd};
  \end{tikzpicture}%
  \vspace*{20mm}%
  \noindent
  \begin{tikzpicture}[remember picture, overlay]
    % 좌측 액센트 바
    \fill[kfPrimary] (current page.west) ++(12mm,25mm) rectangle ++(5mm, -50mm);
    % 챕터 번호 배지 (이중 원, 약간 작게)
    \node[anchor=north west, inner sep=0pt] at ($(current page.north west)+(24mm,-45mm)$) {%
      \begin{tikzpicture}
        \draw[kfPrimary, line width=1pt] (0,0) circle (10mm);
        \draw[kfPrimary, line width=0.4pt] (0,0) circle (8mm);
        \node[font=\normalsize\bfseries\color{kfPrimary}] at (0,2mm) {CHAPTER};
        \node[font=\LARGE\bfseries\color{kfPrimary}] at (0,-3mm) {#1};
      \end{tikzpicture}%
    };
  \end{tikzpicture}%
  \vspace{42mm}%
  \par\noindent\hspace{32mm}%
  \begin{minipage}{0.65\linewidth}%
    {\footnotesize\color{kfMute} #3}\\[4pt]
    {\LARGE\bfseries\color{kfPrimary} #2}\\[8pt]
    {\color{kfRule}\rule{50mm}{0.5pt}}\\[6pt]
    {\small\color{kfInk} #4}
  \end{minipage}%
  \par\vspace{14mm}%
}

% v6 / v18.5 / v18.6: 챕터 표지의 영역 목차 — 흰색 반투명 박스
% v18.6: 배경 가로선/도형과 안 겹치도록 TikZ overlay로 우하단 절대 위치
\makeatletter
% 영역 항목들을 모아 둘 박스
\newsavebox{\kf@arealistbox}
\newenvironment{kfChapterAreaList}{%
  \begin{lrbox}{\kf@arealistbox}%
  \begin{minipage}{0.62\linewidth}%
    \begin{tcolorbox}[%
      enhanced, colback=white, colframe=white,
      opacityback=0.92, opacityframe=0,
      boxrule=0pt, arc=4pt,
      width=\linewidth,
      top=8pt, bottom=8pt, left=10pt, right=10pt,
    ]%
      {\footnotesize\bfseries\color{kfMute} 이 챕터의 영역}\par\vspace{4pt}%
      {\color{kfRule}\hrule height 0.3pt}\par\vspace{4pt}%
      \begin{itemize}[leftmargin=14pt, itemsep=2pt, topsep=0pt, label={\color{kfPrimary!70}\textbullet}]%
}{%
      \end{itemize}%
    \end{tcolorbox}%
  \end{minipage}%
  \end{lrbox}%
  % v18.7: 배경 상단 가로선 ~ 하단 가로선 사이의 중간 영역에 배치
  % 가로선이 내용을 가리지 않도록 박스 중심이 두 선 사이 중점에 위치
  \begin{tikzpicture}[remember picture, overlay]%
    \node[anchor=center, inner sep=0pt]
      at ($(current page.south)+(0mm,90mm)$)
      {\usebox{\kf@arealistbox}};%
  \end{tikzpicture}%
  \vfill%
  \noindent\hfill\kfSeriesBadge\hspace{12mm}%
  \clearpage%
}
\makeatother

% 영역 항목: \kfChapterAreaItem{영역명}{kfArea색}{부제}
\makeatletter
\newcommand{\kfChapterAreaItem}[3]{%
  \item {\small\bfseries\color{#2} #1}%
  \def\kf@cci@sub{#3}%
  \ifx\kf@cci@sub\@empty\else
    {\small\color{kfMute}\, \textemdash\, #3}%
  \fi
}
\makeatother

% 영역 목록 없이 cover만 (legacy)
\newcommand{\kfChapterCoverEnd}{%
  \vfill%
  \noindent\hfill\kfSeriesBadge\hspace{12mm}%
  \clearpage%
}

% ============================================================
% area-header.tex (v16)
% 영역 시작 표제 — 시리즈별 차별화 라벨 + 큰 영역명 + 부제
% PAGE_BREAK_POLICY.md §1.4
%
% v16 (디자인 고도화 Phase 1):
%   - 시리즈별 라벨 박스 추가:
%     소쉬르 = AREA START (영역색 채움 박스)
%     프레게 = SECTION OPEN (영역색 테두리)
%     러셀   = RUSSELL (영역색 라이트 박스)
%     비트   = WITTGENSTEIN (회색 라이트 박스)
%   - 라벨 위치: 영역명 위 좌측
% ============================================================

\makeatletter

% 시리즈 라벨 텍스트
\newcommand{\kf@serieslabel}{%
  \ifcase\kf@series\relax
    AREA START%
  \or
    AREA START%
  \or
    SECTION OPEN%
  \or
    RUSSELL%
  \or
    WITTGENSTEIN%
  \fi
}

% 시리즈 라벨 박스 스타일 (#1 = 영역색)
\newcommand{\kf@labelbox}[1]{%
  \ifcase\kf@series\relax
    % 0/1 = 소쉬르: 영역색 채움
    \tikz[baseline=-2pt]{\node[fill=#1, text=white,
      rounded corners=3pt, inner xsep=6pt, inner ysep=2pt,
      font=\scriptsize\bfseries]{\kf@serieslabel};}%
  \or
    % 1 = 소쉬르: 영역색 채움
    \tikz[baseline=-2pt]{\node[fill=#1, text=white,
      rounded corners=3pt, inner xsep=6pt, inner ysep=2pt,
      font=\scriptsize\bfseries]{\kf@serieslabel};}%
  \or
    % 2 = 프레게: 영역색 테두리
    \tikz[baseline=-2pt]{\node[draw=#1, line width=0.6pt, text=#1,
      rounded corners=3pt, inner xsep=6pt, inner ysep=2pt,
      font=\scriptsize\bfseries]{\kf@serieslabel};}%
  \or
    % 3 = 러셀: 영역색 라이트 박스
    \tikz[baseline=-2pt]{\node[fill=#1!12, text=#1,
      rounded corners=2pt, inner xsep=5pt, inner ysep=1.5pt,
      font=\scriptsize\bfseries]{\kf@serieslabel};}%
  \or
    % 4 = 비트: 회색 미니 박스
    \tikz[baseline=-2pt]{\node[fill=kfMute!10, text=kfMute,
      rounded corners=2pt, inner xsep=5pt, inner ysep=1.5pt,
      font=\scriptsize\bfseries]{\kf@serieslabel};}%
  \fi
}

% v17: 시리즈+영역별 Canva 배경 자동 매핑
% \kf@hasbg=1로 set되면 \kfAreaStart에서 라벨 박스 출력 생략
\def\kf@areaVocab{어휘}\def\kf@areaGrammar{문법}\def\kf@areaConcept{개념}
\def\kf@areaLit{문학}\def\kf@areaNonlit{비문학}
\newif\ifkf@bgon
% 파일 있으면 배경 적용 + boolean true. 없으면 nothing.
\newcommand{\kfBgSet}[1]{%
  \IfFileExists{#1.pdf}{\kfPageBg{#1}\kf@bgontrue}{}%
}
\newcommand{\kf@trybg}[1]{%
  % v18.4: 영역 배경 PDF 비활성화 — 큰 도형이 본문 영역을 침범해
  % "영역 헤더 후 빈 페이지" 문제를 일으킴. 라벨 박스로만 영역 표시.
  \kf@bgonfalse%
}

% Note: 러셀(rus_*)/비트(wit_*) 배경 PDF 미존재 시 \kfPageBg가 에러 → \IfFileExists로 가드
% (현재는 파일이 모두 존재한다고 가정. 부분 제공 시 cls 수정 필요)

\newcommand{\kfAreaStart}[3]{%
  \clearpage%
  \thispagestyle{fancy}%
  \kfMarkArea{#1}%
  \kf@trybg{#1}%
  \vspace*{1mm}%
  % 배경 없으면 라텍스 라벨 박스 출력 (배경 있으면 일러스트가 라벨 역할)
  \ifkf@bgon
    \vspace*{4mm}%
  \else
    \noindent\kf@labelbox{#2}\par\vspace{4pt}%
  \fi
  % 영역명 + 부제
  \noindent
  \begin{tikzpicture}
    \node[anchor=west, inner sep=0pt] (bar) at (0,0) {\color{#2}\rule{3pt}{22pt}};
    \node[anchor=base west, inner sep=0pt, xshift=8pt, yshift=2pt] (lbl) at (bar.east |- 0,0) {%
      {\LARGE\bfseries\color{#2} #1}%
    };
    \def\kf@areasub{#3}%
    \ifx\kf@areasub\@empty\else
      \node[anchor=base west, inner sep=0pt, xshift=8pt, font=\normalsize\color{kfMute}]
        at (lbl.base east) {\,\textperiodcentered\, #3};%
    \fi
  \end{tikzpicture}%
  \par\vspace{2pt}
  {\color{#2!70}\hrule height 1pt}%
  \par\vspace{1pt}%
  {\color{kfRule}\hrule height 0.3pt}%
  \par\vspace{4pt}%
  \needspace{6\baselineskip}\nopagebreak[4]%
  \global\kf@areaJustOpenedtrue% v18.5: 첫 컨텐츠 needspace 무시
}
\makeatother

% 시리즈/영역 단축 매크로
\newcommand{\kfStartVocab}[1]{\kfAreaStart{어휘}{kfAreaVocab}{#1}}
\newcommand{\kfStartGrammar}[1]{\kfAreaStart{문법}{kfAreaGrammar}{#1}}
\newcommand{\kfStartConcept}[1]{\kfAreaStart{개념}{kfAreaConcept}{#1}}
\newcommand{\kfStartLit}[1]{\kfAreaStart{문학}{kfAreaLiterature}{#1}}
\newcommand{\kfStartNonlit}[1]{\kfAreaStart{비문학}{kfAreaNonlit}{#1}}
\newcommand{\kfStartWeekly}[1]{\kfAreaStart{실력 확인}{kfAreaWeekly}{#1}}

% ============================================================
% section-header.tex
% 섹션 시작 라벨 헤더 — [지문] / [활동] / [문제] 라벨
% 좌측 액센트 바 + 라벨 + 부제 (영역 액센트 색)
%
% 사용:
%   \kfSecHeader{kfAreaLiterature}{지문}{빼앗긴 들에도 봄은 오는가 / 이상화}
%   \kfSecHeader{kfAreaConcept}{활동}{내용 확인}
%   \kfSecHeader{kfAreaNonlit}{문제}{}
%
% 헤더 직후 페이지 분할 금지: \needspace{4\baselineskip} + \nopagebreak
% ============================================================

\providecommand{\kfSecHeader}[3]{%
  \par\needspace{10\baselineskip}\filbreak%
  \vspace{3pt}%
  \noindent
  \begin{tikzpicture}
    \node[anchor=west, inner sep=0pt] (bar) at (0,0) {\color{#1}\rule{3pt}{14pt}};
    \node[anchor=west, inner sep=0pt, xshift=6pt, font=\normalsize\bfseries\color{#1}]
       (lbl) at (bar.east) {[\,#2\,]};
    \node[anchor=west, inner sep=0pt, xshift=4pt, font=\small\color{kfMute}]
       at (lbl.east) {#3};
  \end{tikzpicture}%
  \par\vspace{1pt}%
  {\color{#1!50}\hrule height 0.4pt}%
  \par\vspace{2pt}\nopagebreak%
}

% 영역별 단축 매크로
\newcommand{\kfSecPassageVocab}[1]{\kfSecHeader{kfAreaVocab}{지문}{#1}}
\newcommand{\kfSecPassageGrammar}[1]{\kfSecHeader{kfAreaGrammar}{지문}{#1}}
\newcommand{\kfSecPassageConcept}[1]{\kfSecHeader{kfAreaConcept}{지문}{#1}}
\newcommand{\kfSecPassageLit}[1]{\kfSecHeader{kfAreaLiterature}{지문}{#1}}
\newcommand{\kfSecPassageNonlit}[1]{\kfSecHeader{kfAreaNonlit}{지문}{#1}}
\newcommand{\kfSecPassageWeekly}[1]{\kfSecHeader{kfAreaWeekly}{지문}{#1}}

\newcommand{\kfSecActVocab}[1]{\kfSecHeader{kfAreaVocab}{활동}{#1}}
\newcommand{\kfSecActGrammar}[1]{\kfSecHeader{kfAreaGrammar}{활동}{#1}}
\newcommand{\kfSecActConcept}[1]{\kfSecHeader{kfAreaConcept}{활동}{#1}}
\newcommand{\kfSecActLit}[1]{\kfSecHeader{kfAreaLiterature}{활동}{#1}}
\newcommand{\kfSecActNonlit}[1]{\kfSecHeader{kfAreaNonlit}{활동}{#1}}
\newcommand{\kfSecActWeekly}[1]{\kfSecHeader{kfAreaWeekly}{활동}{#1}}

\newcommand{\kfSecQVocab}[1]{\kfSecHeader{kfAreaVocab}{문제}{#1}}
\newcommand{\kfSecQGrammar}[1]{\kfSecHeader{kfAreaGrammar}{문제}{#1}}
\newcommand{\kfSecQConcept}[1]{\kfSecHeader{kfAreaConcept}{문제}{#1}}
\newcommand{\kfSecQLit}[1]{\kfSecHeader{kfAreaLiterature}{문제}{#1}}
\newcommand{\kfSecQNonlit}[1]{\kfSecHeader{kfAreaNonlit}{문제}{#1}}
\newcommand{\kfSecQWeekly}[1]{\kfSecHeader{kfAreaWeekly}{문제}{#1}}


\begin{document}

\thispagestyle{kfBlank}
\kfBookCoverBg
\vspace*{\kfBookCoverVspace}
\begin{center}
{\kfLogoLg}\\[14pt]
{\fontsize{72}{84}\selectfont\bfseries\kfBookCoverColor 러셀}\\[18pt]
{\fontsize{28}{32}\selectfont\kfBookCoverSubColor \textsc{Level} \textbf{2}}\\[22pt]
{\large\kfBookCoverSubColor 제 1 권 \quad\textbar\quad 챕터 1\textendash{} 4}
\end{center}
\vfill
\hfill\kfSeriesBadge\hspace{15mm}
\clearpage
\kfChapterCover{1}{러셀2 (챕터1)}{Chapter 1}{이번 챕터의 학습 내용을 시작합니다.}
\begin{kfChapterAreaList}
\kfChapterAreaItem{문법}{kfAreaGrammar}{단어에게 이름표를! 품사 분류의 원리와 9품사}
\kfChapterAreaItem{개념}{kfAreaConcept}{누가, 어떤 상황에서, 무엇을 느꼈을까?}
\kfChapterAreaItem{문학}{kfAreaLiterature}{「서시」 — 윤동주}
\kfChapterAreaItem{비문학}{kfAreaNonlit}{[인문] 양심을 통한 최대 행복 실현, 밀의 공리주의}
\end{kfChapterAreaList}

\kfStartGrammar{단어에게 이름표를! 품사 분류의 원리와 9품사}


\kfPassageNoteNT{%
　세상에 존재하는 수많은 단어는 저마다 고유한 성격을 가지고 있다. 문법학자들은 이 단어들을 일정한 기준에 따라 갈래를 나누고 이름을 붙였는데, 이것이 바로 품사이다. 품사는 단어의 성질을 밝혀주는 이름표와 같아서, 품사를 이해하는 것은 문법 학습의 가장 기본이 되는 첫걸음이다. 단어를 품사로 분류할 때는 크게 형태, 기능, 의미라는 세 가지 기준을 종합적으로 사용한다.

　가장 먼저 형태를 기준으로 단어의 ‘변화’ 여부를 살필 수 있다. ‘하늘’, ‘매우’처럼 문장 속에서 형태가 고정되어 변하지 않는 단어를 불변어라 하고, ‘먹다’가 ‘먹고, 먹으니’처럼 상황에 따라 모습이 바뀌는 단어를 가변어라고 한다. 국어에서는 동사, 형용사와 서술격 조사 ‘이다’만이 가변어에 속하며, 이를 제외한 대부분의 단어는 불변어이다.

　다음으로 기능에 따라 단어가 문장 안에서 어떤 역할을 수행하는지 묶을 수 있다. 문장의 뼈대 역할을 하는 체언, 주로 체언에 붙어 문법적 관계를 나타내는 관계언, 주어의 움직임이나 상태를 서술하는 용언, 다른 말을 꾸며주는 수식언, 그리고 문장의 다른 성분과 관계없이 독립적으로 쓰이는 독립언, 이렇게 다섯 가지로 기능적 구분을 할 수 있다.

　마지막으로 이 기능적 갈래를 각 단어가 지닌 고유한 의미에 따라 더 잘게 나누면 비로소 아홉 개의 품사가 완성된다. 먼저 체언에는 ‘학교’, ‘우정’처럼 사람이나 사물의 이름을 나타내는 명사, ‘우리’, ‘저것’과 같이 그 이름을 대신하는 대명사, 그리고 ‘다섯’, ‘셋째’처럼 수량이나 순서를 가리키는 수사가 있다. 관계언에는 ‘이/가’, ‘마저’, ‘와/과’처럼 다른 말과의 문법적 관계를 표시해 주는 조사가 유일하게 속한다. 용언은 ‘뛰다’, ‘생각하다’와 같이 대상의 움직임을 나타내는 동사와 ‘착하다’, ‘아름답다’처럼 상태나 성질을 나타내는 형용사로 구분된다. 수식언은 꾸밈을 받는 대상에 따라, ‘새’, ‘모든’, ‘한’과 같이 체언을 꾸미는 관형사와 ‘정말’, ‘멀리’처럼 주로 용언을 꾸미는 부사로 나뉜다. 끝으로 독립언에는 ‘어머나’, ‘네’, ‘아’처럼 놀람, 부름, 대답 등을 표현하는 감탄사가 있다.

　결국 품사란, 형태와 기능 그리고 의미라는 체계적인 기준을 통해 단어의 성격을 종합적으로 규정한 문법적 신분증이라고 할 수 있다.
\par\medskip\begin{itemize}[leftmargin=18pt, itemsep=2pt, topsep=2pt, label=\textbullet]
\item 형태 — 불변어(하늘) / 가변어(먹다·예쁘다·이다)
\item 기능 — 체언·관계언·용언·수식언·독립언
\item 9품사 — 명·대·수 / 조 / 동·형 / 관·부 / 감
\end{itemize}
}

\kfActivityBg \par\kfMaybeNeedspace{32\baselineskip}\noindent{\kfTipMark\,\,}{\small\color{kfInk} 지문의 핵심 내용을 떠올리며 빈칸에 알맞은 말을 채워 보세요.}\par\nopagebreak[4]\smallskip\nopagebreak[4]

\kfActivityBg \kfSecHeader{kfAreaGrammar}{활동}{분석 훈련}
\par\kfMaybeNeedspace{24\baselineskip}\noindent{\kfTipMark\,\,}{\small\color{kfInk} 다음 문장에 쓰인 단어들을 체언과 관계언으로 나누어 각각 모두 찾아 쓰시오.}\par\nopagebreak[4]\smallskip\nopagebreak[4]
\begin{enumerate}[leftmargin=22pt, itemsep=6pt, topsep=4pt, label=\textbf{\arabic*.}]
\item 다음 문장에서 문장의 뼈대가 되는 체언을 모두 찾아 쓰시오.
\begin{kfEvidence}저 학생은 우리 학교의 자랑이자 희망이다.\end{kfEvidence}
\item 다음 문장에서 다른 말에 붙어 관계를 맺어주는 조사를 모두 찾아 쓰시오.
\begin{kfEvidence}너는 사과와 빵 중에서 무엇을 먹겠니?\end{kfEvidence}
\item 다음 문장에서 주어의 움직임이나 상태를 설명하는 용언을 찾아 기본형으로 쓰시오.
\begin{kfEvidence}가을 하늘이 참 푸르고 아름답게 빛난다.\end{kfEvidence}
\item 다음 문장에서 다른 말을 꾸며주는 수식언을 모두 찾아 쓰시오.
\begin{kfEvidence}새 옷을 입으니 기분이 아주 좋았다.\end{kfEvidence}
\item 다음 문장에서 놀람이나 느낌을 나타내는 감탄사를 찾아 쓰시오.
\begin{kfEvidence}어머나, 벌써 시간이 이렇게 되었네!\end{kfEvidence}
\item 다음 문장에서 형태가 변하는 가변어를 모두 찾아 기본형으로 쓰시오.
\begin{kfEvidence}예쁜 꽃이 피고 향기가 진하다.\end{kfEvidence}
\item 다음 문장에서 형태가 변하지 않는 불변어를 모두 찾아 쓰시오.
\begin{kfEvidence}와, 저 새는 정말 빨리 나는구나!\end{kfEvidence}
\item 다음 문장에서 체언(명사, 대명사, 수사)을 꾸며주는 관형사를 찾아 쓰시오.
\begin{kfEvidence}모든 학생은 각자의 자리에서 한 문제를 풀었다.\end{kfEvidence}
\item 다음 문장에서 용언(동사, 형용사)이나 다른 부사를 꾸며주는 부사를 찾아 쓰시오.
\begin{kfEvidence}기차가 매우 빨리 달려서 금방 도착했다.\end{kfEvidence}
\item 다음 문장에서 움직임을 나타내는 동사와 상태를 나타내는 형용사를 구분하여 찾아보시오.
\begin{kfEvidence}착한 아이가 의자에 앉아 책을 읽는다.\end{kfEvidence}
\end{enumerate}

\kfSecHeader{kfAreaGrammar}{문제}{}

\par\medskip
\kfBeginMC
\begin{kfQBlock}{01}{객관식}
\kfQStem{다음 중 품사에 대한 설명으로 적절하지 \uline{\underline{않은}} 것은?}
\begin{kfChoices}
\kfChoice 단어를 기능, 의미, 형태에 따라 나눈 것이다.
\kfChoice 단어의 문법적 성질을 알려주는 이름표이다.
\kfChoice 모든 품사는 문장에서 형태가 변하지 않는다.
\kfChoice 우리말의 품사는 총 9개로 나눌 수 있다.
\kfChoice 품사를 이해하는 것은 문법 학습의 기초이다.
\end{kfChoices}
\end{kfQBlock}
\begin{kfQBlock}{02}{객관식}
\kfQStem{다음 단어들 중 형태에 따른 분류가 나머지와 \uline{다른} 하나는?}
\begin{kfChoices}
\kfChoice 하늘
\kfChoice 그리고
\kfChoice 첫째
\kfChoice 예쁘다
\kfChoice 에서
\end{kfChoices}
\end{kfQBlock}
\begin{kfQBlock}{03}{객관식}
\kfQStem{<보기>의 설명에 해당하는 품사들로만 짝지어진 것은?}
\begin{kfEvidence}문장에서 다른 말을 꾸며주는 기능을 한다.\end{kfEvidence}
\begin{kfChoices}
\kfChoice 명사, 대명사
\kfChoice 동사, 형용사
\kfChoice 조사, 감탄사
\kfChoice 관형사, 부사
\kfChoice 수사, 명사
\end{kfChoices}
\end{kfQBlock}
\begin{kfQBlock}{04}{객관식}
\kfQStem{다음 문장의 밑줄 친 단어 중 품사가 \uline{다른} 하나는?}
\begin{kfChoices}
\kfChoice \uline{새} 신을 신었다.
\kfChoice \uline{저} 아이는 내 동생이다.
\kfChoice \uline{옛} 추억이 떠오른다.
\kfChoice 나는 \uline{아주} 행복하다.
\kfChoice \uline{모든} 국민은 법 앞에 평등하다.
\end{kfChoices}
\end{kfQBlock}
\begin{kfQBlock}{05}{객관식}
\kfQStem{다음 단어 중 기능에 따른 분류가 나머지와 \uline{다른} 하나는?}
\begin{kfChoices}
\kfChoice 학교
\kfChoice 그것
\kfChoice 셋째
\kfChoice 에게
\kfChoice 사랑
\end{kfChoices}
\end{kfQBlock}
\begin{kfQBlock}{06}{객관식}
\kfQStem{밑줄 친 단어의 품사가 바르게 연결되지 \uline{\underline{않은}} 것은?}
\begin{kfChoices}
\kfChoice \uline{다섯} 사람이 한자리에 모두 모였다. - 수사
\kfChoice \uline{와}, 오늘 옷이 정말 멋지다! - 감탄사
\kfChoice 가을 하늘\uline{이} 유난히 푸르고 맑다. - 조사
\kfChoice 그는 \uline{몹시} 지쳐 보였다. - 관형사
\kfChoice 나는 따뜻한 밥을 \uline{먹는다}. - 동사
\end{kfChoices}
\end{kfQBlock}
\begin{kfQBlock}{07}{객관식}
\kfQStem{다음 중 용언에 해당하는 단어들로만 묶인 것은?}
\begin{kfChoices}
\kfChoice 책상, 연필
\kfChoice 나, 우리
\kfChoice 하나, 첫째
\kfChoice 걷다, 슬프다
\kfChoice 매우, 빨리
\end{kfChoices}
\end{kfQBlock}
\begin{kfQBlock}{08}{객관식}
\kfQStem{문장 속 역할에 대한 설명이 바르지 \uline{\underline{않은}} 것은?}
\begin{kfChoices}
\kfChoice 명사: 주로 주어나 목적어 역할을 한다.
\kfChoice 조사: 홀로 쓰일 수 없고 다른 말에 붙어 쓰인다.
\kfChoice 형용사: 대상의 움직임을 나타낸다.
\kfChoice 관형사: 체언 앞에 놓여 체언을 꾸민다.
\kfChoice 감탄사: 문장의 다른 성분과 관계없이 독립적으로 쓰인다.
\end{kfChoices}
\end{kfQBlock}
\begin{kfQBlock}{09}{객관식}
\kfQStem{다음 중 불변어가 \uline{\underline{아닌}} 것은?}
\begin{kfChoices}
\kfChoice 학생이다
\kfChoice 나무
\kfChoice 그리고
\kfChoice 까지
\kfChoice 새
\end{kfChoices}
\end{kfQBlock}
\begin{kfQBlock}{10}{객관식}
\kfQStem{다음 문장에 사용되지 \uline{\underline{않은}} 품사는?}
\begin{kfEvidence}와, 저 새는 정말 높이 나는구나!\end{kfEvidence}
\begin{kfChoices}
\kfChoice 형용사
\kfChoice 관형사
\kfChoice 부사
\kfChoice 동사
\kfChoice 감탄사
\end{kfChoices}
\end{kfQBlock}
\begin{kfQBlock}{11}{객관식}
\kfQStem{다음 중 품사의 종류가 \uline{다른} 하나는?}
\begin{kfChoices}
\kfChoice 우리
\kfChoice 이것
\kfChoice 다섯
\kfChoice 무엇
\kfChoice 그
\end{kfChoices}
\end{kfQBlock}
\begin{kfQBlock}{12}{객관식}
\kfQStem{<보기>의 ㉠, ㉡에 들어갈 품사가 바르게 짝지어진 것은?}
\begin{kfEvidence}㉠은 체언을 꾸며주고, ㉡은 주로 용언을 꾸며준다.\end{kfEvidence}
\begin{kfChoices}
\kfChoice ㉠ 부사, ㉡ 관형사
\kfChoice ㉠ 관형사, ㉡ 부사
\kfChoice ㉠ 명사, ㉡ 동사
\kfChoice ㉠ 동사, ㉡ 형용사
\kfChoice ㉠ 조사, ㉡ 감탄사
\end{kfChoices}
\end{kfQBlock}
\begin{kfQBlock}{13}{객관식}
\kfQStem{밑줄 친 단어 중 품사가 동사인 것은?}
\begin{kfChoices}
\kfChoice 하늘이 매우 \uline{푸르다}.
\kfChoice 나는 학생\uline{이다}.
\kfChoice 아기가 방에서 잠을 \uline{잔다}.
\kfChoice 내 동생은 참 \uline{착하다}.
\kfChoice 산이 \uline{높다}.
\end{kfChoices}
\end{kfQBlock}
\begin{kfQBlock}{14}{객관식}
\kfQStem{품사를 기능에 따라 나눌 때, 짝지어진 것 중 성격이 \uline{다른} 하나는?}
\begin{kfChoices}
\kfChoice 수식언 - 형용사
\kfChoice 관계언 - 조사
\kfChoice 독립언 - 감탄사
\kfChoice 용언 - 동사
\kfChoice 체언 - 명사
\end{kfChoices}
\end{kfQBlock}
\begin{kfQBlock}{15}{객관식}
\kfQStem{다음 문장에서 조사를 제외한 단어의 개수는?}
\begin{kfEvidence}저 아이가 새 책을 매우 빨리 읽었다.\end{kfEvidence}
\begin{kfChoices}
\kfChoice 4개
\kfChoice 5개
\kfChoice 6개
\kfChoice 7개
\kfChoice 8개
\end{kfChoices}
\end{kfQBlock}
\begin{kfQBlock}{16}{단답형}
\kfQStem{단어를 성질에 따라 나눈 갈래를 무엇이라고 하는가?}
\kfAnswerLines{1}
\end{kfQBlock}
\begin{kfQBlock}{17}{단답형}
\kfQStem{품사를 분류하는 세 가지 기준은 무엇인가?}
\kfAnswerLines{1}
\end{kfQBlock}
\begin{kfQBlock}{18}{단답형}
\kfQStem{문장에서 형태가 변하는 단어를 무엇이라고 하는가?}
\kfAnswerLines{1}
\end{kfQBlock}
\begin{kfQBlock}{19}{단답형}
\kfQStem{문장에서 형태가 변하지 않는 단어를 무엇이라고 하는가?}
\kfAnswerLines{1}
\end{kfQBlock}
\begin{kfQBlock}{20}{단답형}
\kfQStem{문장의 뼈대를 이루는 명사, 대명사, 수사를 묶어 기능에 따라 무엇이라고 하는가?}
\kfAnswerLines{1}
\end{kfQBlock}
\begin{kfQBlock}{21}{단답형}
\kfQStem{주로 체언 뒤에 붙어 문법적 관계를 나타내는 조사가 속한 기능적 분류는 무엇인가?}
\kfAnswerLines{1}
\end{kfQBlock}
\begin{kfQBlock}{22}{단답형}
\kfQStem{주어의 움직임이나 상태를 서술하는 동사, 형용사가 속한 기능적 분류는 무엇인가?}
\kfAnswerLines{1}
\end{kfQBlock}
\begin{kfQBlock}{23}{단답형}
\kfQStem{다른 말을 꾸며주는 관형사, 부사가 속한 기능적 분류는 무엇인가?}
\kfAnswerLines{1}
\end{kfQBlock}
\begin{kfQBlock}{24}{단답형}
\kfQStem{문장에서 독립적으로 쓰이는 감탄사가 속한 기능적 분류는 무엇인가?}
\kfAnswerLines{1}
\end{kfQBlock}
\begin{kfQBlock}{25}{단답형}
\kfQStem{‘뛰다’, ‘웃다’처럼 움직임을 나타내는 품사는 무엇인가?}
\kfAnswerLines{1}
\end{kfQBlock}
\begin{kfQBlock}{26}{서술형}
\kfQStem{품사를 분류하는 세 가지 기준인 '형태', '기능', '의미'가 각각 무엇을 기준으로 단어를 나누는 것인지 간략하게 설명하시오.}
\kfAnswerNote{답안}{4}
\end{kfQBlock}
\kfEndMC


\kfStartConcept{누가, 어떤 상황에서, 무엇을 느꼈을까?}


\kfPassageNoteNT{%
　시는 함축적인 언어로 정서와 의미를 전달하는 문학이다. 작품의 의미를 깊이 있게 이해하려면, 시를 이루는 요소들을 차례대로 분석하는 과정이 필요하다.

　 시를 읽을 때 가장 먼저 살펴보아야 할 것은 '누가 말하고 있는가'이다. 시에서 자신의 생각과 감정을 전달하는 이를 '시적 화자' 또는 '서정적 자아'라고 한다. 화자는 시를 쓴 시인 자신일 수도 있고, 시인이 만들어 낸 가상의 인물일 수도 있다. 예를 들어, 윤동주 「서시」에서 "죽는 날까지 하늘을 우러러 한 점 부끄럼이 없기를"이라고 다짐하는 화자는 시인 윤동주 본인의 모습으로 이해할 수 있다. 이와 달리 김소월 「엄마야 누나야」에서 "엄마야 누나야 강변 살자"라고 말하는 화자는 평화로운 삶을 꿈꾸는 아이의 모습으로 형상화된 가상의 인물이다. 이처럼 시인은 자신이 전달하고자 하는 주제 의식을 가장 효과적으로 드러내기 위해, 의도적으로 시인 자신을 내세우거나 특정한 인물을 화자로 설정한다.

　 화자를 파악한 후에는 그 화자가 처한 '시적 상황'을 이해해야 한다. '시적 상황'이란 화자가 놓인 시간적, 공간적 배경이나 구체적인 처지를 말한다. 예를 들어 「서시」의 화자는 '별이 바람에 스치는 밤'이라는 시간적 배경 속에서 '죽는 날까지'의 삶을 다짐하며 자신의 과거를 돌아본다. 이러한 상황은 자신의 삶을 엄격하게 돌아보는 '성찰'의 순간이라 할 수 있다.

　 '시적 상황'을 이해했다면, 다음은 화자가 느끼는 '정서', 즉 기쁨, 슬픔, 그리움과 같은 다양한 마음 상태를 파악해야 한다. 화자의 정서는 시에 사용된 단어나 표현에 주목할 때 잘 드러난다. 「서시」의 화자는 "잎새에 이는 바람에도 나는 괴로워했다"라고 고백한다. 아주 사소한 흔들림을 뜻하는 '잎새에 이는 바람'에도 '괴로워했다'는 표현은, 화자가 순결한 삶을 지향하면서 느끼는 '고뇌'의 정서를 드러낸다. 이러한 정서를 바탕으로 시인이 궁극적으로 말하고자 하는 '주제'가 드러난다.「서시」의 '고뇌'와 '부끄러움'이라는 정서는 결국 '순수한 삶에 대한 소망과 의지'라는 주제로 연결된다.

　 이처럼 시적 화자, 시적 상황, 정서, 주제를 파악하는 것은, 시인이 전달하고자 한 주제를 정확히 이해하고 작품을 보다 풍부하게 감상하는 데 효과적인 방법이다.
\par\medskip\begin{itemize}[leftmargin=18pt, itemsep=2pt, topsep=2pt, label=\textbullet]
\item 화자=시인 — 「서시」의 윤동주 자신
\item 화자=가상 인물 — 「엄마야 누나야」의 어린 아이
\item 정서·주제 — 「서시」의 '고뇌·부끄러움' → '순수한 삶의 소망과 의지'
\end{itemize}
}

\kfSecHeader{kfAreaConcept}{문제}{}

\par\medskip
\kfBeginMC
\begin{kfQBlock}{01}{객관식}
\kfQStem{시적 화자에 대한 설명으로 적절하지 \uline{\underline{않은}} 것은?}
\begin{kfChoices}
\kfChoice 시인 자신일 수 있다.
\kfChoice 가상의 인물일 수 있다.
\kfChoice 서정적 자아라고도 한다.
\kfChoice 시의 주제 의식과는 관련이 없다.
\kfChoice 자신의 생각과 감정을 시에서 전달한다.
\end{kfChoices}
\end{kfQBlock}
\begin{kfQBlock}{02}{객관식}
\kfQStem{화자의 정서를 파악하는 방법으로 적절하지 \uline{\underline{않은}} 것은?}
\begin{kfChoices}
\kfChoice 시에 사용된 단어와 표현의 결을 직접 살펴본다.
\kfChoice 시의 창작 연도가 작가의 출생 연도와 일치하는지 본다.
\kfChoice 화자가 바라보는 대상이나 풍경의 인상을 살펴본다.
\kfChoice 시 전체에 흐르는 분위기와 어조를 함께 읽어 낸다.
\kfChoice 화자가 처한 시적 상황과 처지를 함께 고려한다.
\end{kfChoices}
\end{kfQBlock}
\begin{kfQBlock}{03}{객관식}
\kfQStem{위 글의 내용과 일치하지 \uline{않는} 것은?}
\begin{kfChoices}
\kfChoice 시는 함축적인 언어를 사용한다.
\kfChoice 「서시」의 화자는 시인 자신으로 볼 수 있다.
\kfChoice 「엄마야 누나야」의 화자는 가상의 인물이다.
\kfChoice 시적 상황은 화자가 느끼는 감정이나 기분을 말한다.
\kfChoice 시적 요소를 분석하면 작품을 풍부하게 감상할 수 있다.
\end{kfChoices}
\end{kfQBlock}
\begin{kfQBlock}{04}{객관식}
\kfQStem{다음 <보기>의 시적 화자에 대한 설명으로 적절하지 \uline{\underline{않은}} 것은?}
\begin{kfEvidence}열무 삼십 단을 이고 \\\relax
시장에 간 우리 엄마 \\\relax
안 오시네, 해는 시든지 오래 \\\relax
나는 찬밥처럼 방에 담겨 \\\relax
아무리 천천히 숙제를 해도 \\\relax
엄마 안 오시네 배춧잎 같은 발소리 타박타박 \\\relax
안 들리네, 어둡고 무서워 \\\relax
금간 창 틈으로 고요히 빗소리 \\\relax
빈방에 혼자 엎드려 훌쩍거리던

아주 먼 옛날 \\\relax
지금도 내 눈시울을 뜨겁게 하는 \\\relax
그 시절, 내 유년의 윗목. \\\relax
— 기형도 「엄마 걱정」\end{kfEvidence}
\begin{kfChoices}
\kfChoice 1연의 화자는 '아이'이다.
\kfChoice 2연의 화자는 '어른'이다.
\kfChoice 화자는 '찬밥'을 먹고 있다.
\kfChoice 1연의 화자는 엄마를 기다리고 있다.
\kfChoice 2연의 화자는 유년 시절을 회상한다.
\end{kfChoices}
\end{kfQBlock}
\begin{kfQBlock}{05}{객관식}
\kfQStem{다음 <보기>의 화자가 처한 '시적 상황'에 대한 이해로 가장 적절한 것은?}
\begin{kfEvidence}엄마야 누나야 강변 살자, \\\relax
뜰에는 반짝이는 금모래빛, \\\relax
뒷문 밖에는 갈잎의 노래, \\\relax
엄마야 누나야 강변 살자. \\\relax
— 김소월 「엄마야 누나야」\end{kfEvidence}
\begin{kfChoices}
\kfChoice 화자가 원하는 이상적인 삶의 모습을 그리고 있다.
\kfChoice 누나와 함께 엄마를 기다리고 있다.
\kfChoice 현재 살고 있는 강변의 풍경을 묘사하고 있다.
\kfChoice 가족과 헤어져 홀로 남겨진 슬픔을 노래하고 있다.
\kfChoice 엄마와 함께 강변에 살고 있다.
\end{kfChoices}
\end{kfQBlock}
\begin{kfQBlock}{06}{객관식}
\kfQStem{이 시의 화자가 처한 '시적 상황'으로 가장 적절한 것은?}
\begin{kfEvidence}내를 건너서 숲으로 \\\relax
고개를 넘어서 마을로 \\\relax
어제도 가고 오늘도 갈 \\\relax
나의 길, 새로운 길 \\\relax
민들레가 피고 까치가 날고 \\\relax
아가씨가 지나고 바람이 일고 \\\relax
나의 길은 언제나 새로운 길 \\\relax
오늘도... 내일도... \\\relax
내를 건너서 숲으로 \\\relax
고개를 넘어서 마을로 \\\relax
— 윤동주 「새로운 길」\end{kfEvidence}
\begin{kfChoices}
\kfChoice 숲 속에서 길을 잃고 마을을 찾아 헤매고 있다.
\kfChoice 처음 가보는 길에서 낯선 풍경을 마주하고 있다.
\kfChoice 매일 걷는 익숙한 길을 새롭게 인식하며 걷고 있다.
\kfChoice 고개를 넘다가 지쳐서 잠시 걸음을 멈추고 쉬고 있다.
\kfChoice 어제 걸었던 길을 떠올리며 과거의 일을 회상하고 있다.
\end{kfChoices}
\end{kfQBlock}
\begin{kfQBlock}{07}{객관식}
\kfQStem{이 시의 화자가 ‘나의 길’에 대해 느끼는 정서로 적절하지 \uline{\underline{않은}} 것은?}
\begin{kfEvidence}내를 건너서 숲으로 \\\relax
고개를 넘어서 마을로 \\\relax
어제도 가고 오늘도 갈 \\\relax
나의 길, 새로운 길 \\\relax
민들레가 피고 까치가 날고 \\\relax
아가씨가 지나고 바람이 일고 \\\relax
나의 길은 언제나 새로운 길 \\\relax
오늘도... 내일도... \\\relax
내를 건너서 숲으로 \\\relax
고개를 넘어서 마을로 \\\relax
— 윤동주 「새로운 길」\end{kfEvidence}
\begin{kfChoices}
\kfChoice 매일 가는 길에서 느끼는 새로움에 대한 기대감
\kfChoice 자신의 잘못을 깊이 뉘우치는 부끄러움
\kfChoice 익숙한 풍경 속에서도 새로움을 발견하는 설렘
\kfChoice 자신의 길을 묵묵히 걸어가려는 의지
\kfChoice ‘오늘도… 내일도…’ 이어질 길에 대한 긍정적 다짐
\end{kfChoices}
\end{kfQBlock}
\begin{kfQBlock}{08}{단답형}
\kfQStem{함축적인 언어를 사용하여 정서와 의미를 전달하는 문학 양식은 무엇인가?}
\kfAnswerLines{1}
\end{kfQBlock}
\begin{kfQBlock}{09}{단답형}
\kfQStem{시에서 자신의 생각과 감정을 전달하는 이를 무엇이라고 하는가?}
\kfAnswerLines{1}
\end{kfQBlock}
\begin{kfQBlock}{10}{단답형}
\kfQStem{화자가 놓인 시간적, 공간적 배경이나 구체적인 처지를 무엇이라고 하는가?}
\kfAnswerLines{1}
\end{kfQBlock}
\begin{kfQBlock}{11}{단답형}
\kfQStem{시적 상황 속에서 화자가 느끼는 감정을 무엇이라고 하는가?}
\kfAnswerLines{1}
\end{kfQBlock}
\begin{kfQBlock}{12}{단답형}
\kfQStem{화자의 정서를 바탕으로 시인이 궁극적으로 말하고자 하는 바를 무엇이라고 하는가?}
\kfAnswerLines{1}
\end{kfQBlock}
\begin{kfQBlock}{13}{단답형}
\kfQStem{윤동주의 시 「서시」에서 '잎새에 이는 바람'에도 괴로워하는 화자가 느끼는 정서는 무엇인가?}
\kfAnswerLines{1}
\end{kfQBlock}
\begin{kfQBlock}{14}{단답형}
\kfQStem{김소월의 시 「엄마야 누나야」에서 평화로운 삶을 꿈꾸는 화자는 어떤 모습으로 형상화되었는가?}
\kfAnswerLines{1}
\end{kfQBlock}
\begin{kfQBlock}{15}{서술형}
\kfQStem{시적 화자의 두 가지 유형을 위 글에 제시된 작품을 예로 들어 서술하시오.}
\kfAnswerNote{답안}{4}
\end{kfQBlock}
\kfEndMC


\kfStartLit{「서시」 — 윤동주}


\kfPassageNote{}{%
죽는 날까지 하늘을 우러러 \\[4pt]
한 점 부끄럼이 없기를, \\[4pt]
잎새에 이는 바람에도 \\[4pt]
나는 괴로워했다. \\[14pt]
별을 노래하는 마음으로 \\[4pt]
모든 죽어가는 것을 사랑해야지 \\[4pt]
그리고 나한테 주어진 길을 \\[4pt]
걸어가야겠다. \\[14pt]
오늘 밤에도 별이 바람에 스치운다.
}


\kfPassageNote{문장 독해 — 발췌 문장}{%
1. 죽는 날까지 하늘을 우러러 한 점 부끄럼이 없기를, 잎새에 이는 바람에도 나는 괴로워했다.

2. \kfFillBlank[30mm]{} 별을 노래하는 마음으로 모든 죽어가는 것을 사랑해야지.

3. 그리고 나한테 주어진 길을 걸어가야겠다.
}


\kfPassageNote{해설 학습 — 본문}{%
[단원 1] 문학은 그 시대의 아픔과 희망을 담는 그릇과 같다. 일제강점기라는 어두운 시대에 활동한 윤동주 시인의 「서시」는 절망적인 현실 속에서도 굴복하지 않는 젊은 지식인의 마음을 보여주는 대표작이다. 시인의 유고 시집 『하늘과 바람과 별과 시』의 서문 역할을 하는 이 시에는, 그의 문학 세계를 이해하는 핵심 정서인 '부끄러움'이 잘 나타나 있다.

윤동주가 느끼는 부끄러움은 단순히 개인의 잘못에서 비롯된 감정이 아니다. 이는 나라를 잃은 시대에 적극적으로 저항하지 못하고 시를 쓰는 일에 머무는 자신을 끊임없이 돌아보는 고결한 성찰의 결과물이다. '잎새에 이는 바람에도 괴로워했다'는 구절은 이러한 작고 섬세한 양심을 상징적으로 보여준다. 이러한 부끄러움은 좌절에 머무르지 않고, '모든 죽어가는 것'으로 표현된 고통받는 민족에 대한 깊은 사랑과 연민으로 이어진다.

결국 시인은 이 사랑을 실천하기 위해 자신에게 '주어진 길'을 운명으로 받아들인다. 여기서 길이란, 민족의 아픔을 대신 짊어지려는 희생적 태도를 의미한다. 따라서 윤동주의 부끄러움은 무력한 지식인의 고뇌를 넘어, 자신의 희생을 통해 어두운 시대를 밝히려는 위대한 결심으로 승화되는 것이다.

[단원 2] 시인은 직접적인 표현보다는 상징을 통해 자신의 마음을 전달하는 경우가 많다. 윤동주 「서시」도 여러 상징적인 단어들을 사용하여 시의 깊은 뜻을 효과적으로 드러낸다.

시에 나오는 '하늘'은 화자가 자신의 모든 행동과 마음을 비추어보는 절대적인 기준이다. 부끄러움 없는 삶을 살고 싶을 때 우러러보는 도덕적인 거울과 같은 존재이다.

'바람'은 시의 흐름에 따라 두 가지 의미를 지닌다. 1연에 나오는 '잎새에 이는 바람'은 깨끗하게 살고 싶은 화자의 마음을 흔드는 내면의 자아 성찰과 괴로움을 의미한다. 반면 마지막 연의 '별이 바람에 스치운다'에서 '바람'은 시인을 억압하는 일제강점기의 차가운 현실, 즉 외부에서 오는 시련을 상징한다.

어두운 밤하늘에서 빛나는 '별'은 암담한 현실 속에서도 사라지지 않는 희망을 의미한다. 동시에 시인이 간직하고 싶은 순수한 양심과 그가 꿈꾸는 이상적인 세계를 나타낸다.

'오늘 밤'은 화자가 처해 있는 일제강점기의 절망적인 상황 그 자체를 뜻한다.

그리고 '길'은 이러한 고통스러운 현실 속에서도 앞으로 나아가야 할 시인의 사명과 삶의 자세를 보여준다.

이처럼 윤동주 시인은 단순한 자연물을 통해 자신의 복잡한 내면세계와 시대의 아픔을 깊이 있게 표현했다.

[단원 3] 문학 작품은 시간의 흐름을 효과적으로 활용하여 주제 의식을 깊이 있게 드러낸다. 「서시」는 과거, 미래, 현재를 넘나드는 입체적인 시간 구성을 통해 화자의 내면을 생생하게 보여주는 작품이다.

1연은 과거의 삶을 스스로 돌아보는 시간이다. 화자는 지금까지 한 점 부끄럼 없는 삶을 살기 위해 아주 작은 바람에도 괴로워할 만큼 자신을 진지하게 살피며 살아왔음을 고백한다. 

2연은 미래를 향한 굳은 다짐을 보여준다. 화자는 순수한 꿈을 간직한 채 고통받는 모든 것을 사랑하고, 자신에게 주어진 특별한 길을 조용히 걸어가겠다는 의지를 다진다. 

3연은 다시 어둡고 힘든 현재로 돌아와 화자가 놓인 상황을 상징적으로 보여준다. 여전히 현실을 뜻하는 '밤'은 어둡고 차가우며, 지켜야 할 소중한 꿈인 '별'은 바깥의 어려움을 뜻하는 '바람'에 위태롭게 흔들리고 있다.

이처럼 과거를 돌아본 것을 바탕으로 미래를 다짐하고, 그 힘으로 힘든 현재를 견뎌내는 짜임새 있는 구성은 화자의 깊은 생각과 의지를 효과적으로 보여준다.
}


\kfActivityBg \par\needspace{15\baselineskip}
\kfSecHeader{kfAreaLiterature}{활동}{문장 독해}
\par\kfMaybeNeedspace{32\baselineskip}\noindent{\kfTipMark\,\,}{\small\color{kfInk} 다음은 작품에서 발췌한 문장입니다. 각 문장을 읽고 물음에 답하시오.}\par\nopagebreak[4]\smallskip\nopagebreak[4]
\kfBeginMC
\begin{kfSentenceUnit}{1}{}
\kfSentQ{1.}{다음 문장에서 서술어 '우러러', '없기를', '이는', '괴로워했다'의 주어로 적절하게 짝지어진 것은?}
\begin{kfChoices}
\kfChoice 우러러(하늘), 없기를(나), 이는(잎새), 괴로워했다(바람)
\kfChoice 우러러(나), 없기를(부끄럼), 이는(바람), 괴로워했다(나)
\kfChoice 우러러(나), 없기를(하늘), 이는(잎새), 괴로워했다(나)
\kfChoice 우러러(죽는 날), 없기를(점), 이는(바람), 괴로워했다(잎새)
\end{kfChoices}
\end{kfSentenceUnit}

\begin{kfSentenceUnit}{2}{}
\kfSentQ{2.}{다음 시구에서 '그리고 나에게 주어진 길'이 의미하는 바로 가장 적절한 것은?}
\begin{kfChoices}
\kfChoice 과거의 잘못을 후회하는 삶
\kfChoice 현실을 도피하여 안주하는 삶
\kfChoice 주어진 운명을 받아들이는 희생적인 삶
\kfChoice 개인의 성공과 영달을 추구하는 삶
\end{kfChoices}
\end{kfSentenceUnit}

\kfEndMC

\kfActivityBg \kfSecHeader{kfAreaLiterature}{활동}{해설 문제}
\par\kfMaybeNeedspace{18\baselineskip}\noindent{\kfTipMark\,\,}{\small\color{kfInk} 해설 본문을 읽고, 단원별 단답형 물음에 답하시오.}\par\nopagebreak[4]\smallskip\nopagebreak[4]
\begin{enumerate}[leftmargin=22pt, itemsep=6pt, topsep=4pt, label=\textbf{\arabic*.}]
\item 윤동주 시인처럼 시대의 아픔을 노래하고 불의에 맞서는 시인을 무엇이라고 부르는가?
\item 「서시」는 윤동주 시인의 어떤 시집에서 서문 역할을 하는가?
\item 이 글에서 설명하는 윤동주 문학 세계의 핵심 정서는 무엇인가?
\item 윤동주가 느끼는 부끄러움은 무엇에 대한 성찰의 결과물인가?
\item 시인의 섬세한 양심을 상징적으로 보여주는 구절은 무엇인가?
\item 시인의 부끄러움은 어떤 감정으로 이어지는가?
\item 시인이 사랑하고자 다짐한 '모든 죽어가는 것'은 누구를 가리키는가?
\item 시에서 '주어진 길'은 어떤 태도를 의미하는가?
\item 화자가 삶의 기준으로 삼는, 절대적 순수함을 상징하는 시어는 무엇인가?
\item 1연에서는 내면의 흔들림을, 마지막 연에서는 외부의 시련을 의미하는 시어는 무엇인가?
\item 암울한 현실 속에서도 잃지 않으려는 희망과 순수한 이상을 상징하는 시어는 무엇인가?
\item 화자가 처한 절망적인 시대를 상징하는 시간적 배경은 무엇인가?
\item 화자가 묵묵히 걸어가야 할 숙명적인 삶의 자세를 비유하는 시어는 무엇인가?
\item 마지막 연의 '바람'은 구체적으로 어떤 시대 현실을 상징하는가?
\item 이 시에서 '비둘기'가 '평화'를 나타내는 것처럼, 눈에 보이지 않는 추상적인 생각이나 의미를 눈에 보이는 구체적인 사물로 대신 나타내는 표현 방법을 무엇이라고 하는가?
\item 「서시」는 어떤 시간 구성을 통해 화자의 내면을 보여주고 있는가?
\item 시에 나오는 1연은 어떤 시간에 해당하며, 화자가 무엇을 하는 시간인가?
\item 2연에서 화자가 미래를 향해 다짐하는 의지 두 가지는 무엇인가?
\item 3연은 어떤 시간으로 돌아오며, 무엇을 상징적으로 보여주고 있는가?
\item 3연의 현재 상황에서, 지켜야 할 꿈인 '별'은 무엇에 의해 위태롭게 흔들리고 있는가?
\end{enumerate}

\kfSecHeader{kfAreaLiterature}{문제}{}

\par\medskip
\kfBeginMC
\begin{kfQBlock}{01}{객관식}
\kfQStem{이 시의 화자가 ‘부끄러움’을 느끼게 되는 까닭으로 가장 알맞은 것은?}
\begin{kfChoices}
\kfChoice 순수한 꿈과 어두운 현실의 차이를 느끼며 자기 자신을 깊이 돌아보기 때문이다.
\kfChoice 다른 사람과 자신을 비교하며 부족한 점을 깨닫고 스스로를 미워하기 때문이다.
\kfChoice 과거에 저질렀던 잘못을 자꾸 떠올리며 마음의 불편함을 느끼기 때문이다.
\kfChoice 시대의 아픔을 외면한 채 자신의 안위만을 생각하는 이기심 때문이었다.
\kfChoice 정해진 운명을 따르지 않고 저항하려는 마음 때문에 현실과 갈등하기 때문이다.
\end{kfChoices}
\end{kfQBlock}
\begin{kfQBlock}{02}{객관식}
\kfQStem{화자가 “잎새에 이는 바람에도 / 나는 괴로워했다.”고 말한 이유로 가장 알맞은 것은?}
\begin{kfChoices}
\kfChoice 몸이 너무 약해서 작은 바람에도 병이 들었기 때문이다.
\kfChoice 아주 작은 부끄러움도 받아들이지 못하는 순수한 양심을 지녔기 때문이다.
\kfChoice 가난 때문에 바람을 잘 막아 줄 좋은 집에서 살지 못했기 때문이다.
\kfChoice 친구들이 자신을 놀리고 괴롭히는 것을 바람에 비유했기 때문이다.
\kfChoice 자신이 저지른 큰 잘못이 들킬까 봐 항상 불안했기 때문이다.
\end{kfChoices}
\end{kfQBlock}
\begin{kfQBlock}{03}{객관식}
\kfQStem{㉠\textasciitilde{}㉤ 중, 화자가 처한 암울한 시대 현실을 상징하는 시어로 가장 알맞은 것은?}
\begin{kfChoices}
\kfChoice ㉤
\kfChoice ㉡
\kfChoice ㉢
\kfChoice ㉣
\kfChoice ㉠
\end{kfChoices}
\end{kfQBlock}
\begin{kfQBlock}{04}{객관식}
\kfQStem{화자가 ‘모든 죽어가는 것’을 사랑하겠다고 다짐한 이유로 가장 알맞은 것은?}
\begin{kfChoices}
\kfChoice 약한 존재들과 힘을 합쳐 자신의 나약함을 이겨내고 강한 마음을 얻기 위해
\kfChoice 신에게 받은 사랑을 자신보다 약한 존재들에게 나눠주려는 종교적인 믿음 때문에
\kfChoice 자신의 마음이 얼마나 깨끗하고 순수한지 증명할 좋은 방법이라고 생각해서
\kfChoice 고통받는 사람들의 아픔에 공감하며 그 슬픔을 함께 나누려는 따뜻한 마음에서
\kfChoice 죽음 앞에서는 모든 생명이 똑같다는 깊은 진리를 혼자 깨달았기 때문에
\end{kfChoices}
\end{kfQBlock}
\begin{kfQBlock}{05}{객관식}
\kfQStem{1연의 ‘바람’과 3연의 ‘바람’에 대한 설명으로 적절하지 \uline{\underline{않은}} 것은?}
\begin{kfChoices}
\kfChoice 1연의 바람은 화자의 한없이 예민한 양심과 내면 깊은 곳의 고뇌를 형상화한다.
\kfChoice 두 ‘바람’은 자연의 섭리를 나타내어 화자의 자연 친화적 태도를 보여 준다.
\kfChoice 1연은 화자 내면을, 3연은 외부 세계를 나타낸다는 점에서 의미층이 구별된다.
\kfChoice 두 ‘바람’은 모두 이 시의 전체 주제를 강조하는 핵심 시어로 작용한다.
\kfChoice 3연의 바람은 밤하늘의 별을 스치며 지나가는 외부의 가혹한 현실을 상징한다.
\end{kfChoices}
\end{kfQBlock}
\begin{kfQBlock}{06}{객관식}
\kfQStem{이 시의 표현 방식에 대한 설명으로 적절하지 \uline{\underline{않은}} 것은?}
\begin{kfChoices}
\kfChoice 대상을 우스꽝스럽게 비꼬는 풍자적 표현이 두드러지게 사용된다.
\kfChoice 추상적 가치를 구체적인 사물에 빗대어 형상화한다.
\kfChoice ‘과거 → 미래 → 현재’의 시간 흐름에 따라 시상이 전개된다.
\kfChoice 의지적·고백적 어조로 화자의 다짐을 드러낸다.
\kfChoice ‘하늘·별·바람·길’ 등 상징적인 시어를 사용하여 함축적 의미를 전달한다.
\end{kfChoices}
\end{kfQBlock}
\begin{kfQBlock}{07}{객관식}
\kfQStem{<보기>를 참고하여 이 시를 감상한 내용으로 적절하지 \uline{\underline{않은}} 것은?}
\begin{kfEvidence}윤동주는 일제 강점기에 활동했던 시인이다. 그는 일본 유학 중 독립운동 혐의로 체포되어 감옥에서 생을 마감했다. 그의 시에는 이러한 시대적 아픔 속에서 고뇌하는 지식인의 모습과 자기희생을 통해 조국의 구원을 바라는 마음이 잘 나타나 있다.\end{kfEvidence}
\begin{kfChoices}
\kfChoice '오늘 밤'은 시인이 겪었던 일제 강점기의 암울한 현실을 의미하겠군.
\kfChoice '바람'은 시인에게 닥쳤던 시련과 억압을 상징한다고 볼 수 있겠군.
\kfChoice '나한테 주어진 길'은 결국 감옥에서의 죽음으로 이어진 희생의 길이었겠군.
\kfChoice '한 점 부끄럼이 없기를' 바라는 마음은 조국 독립에 대한 무관심을 보여주는군.
\kfChoice '모든 죽어가는 것'은 일제에 의해 고통받던 우리 민족을 가리키는 것이겠군.
\end{kfChoices}
\end{kfQBlock}
\begin{kfQBlock}{08}{서술형}
\kfQStem{이 시에서 화자가 간직하고 싶은 소중한 '희망'을 상징하는 시어와, 화자를 힘들게 하는 '시련'을 상징하는 시어를 <조건>에 맞게 쓰시오.}
\kfAnswerNote{답안}{4}
\end{kfQBlock}
\begin{kfQBlock}{09}{서술형}
\kfQStem{2연에서 화자는 앞으로 어떻게 살겠다고 다짐한다. 화자가 마음속으로 다짐한 두 가지를 시에 나온 표현을 사용하여 <조건>에 맞게 쓰시오.}
\kfAnswerNote{답안}{4}
\end{kfQBlock}
\kfEndMC


\kfStartNonlit{[인문] 양심을 통한 최대 행복 실현, 밀의 공리주의}


\kfPassageNote{[인문] 양심을 통한 최대 행복 실현, 밀의 공리주의}{%
　공리주의는 공리를 통해 최대 행복을 실현하는 것을 중시하는 이론이다. 여기서 공리란 이익과 유용함을 뜻하며, 공리주의에서 행복이란 고통을 피하고 쾌락을 추구하는 것을 의미한다. 이때 행복은 개인의 쾌락만이 아니라 개인의 행위와 관련된 사회 구성원의 쾌락도 함께 고려하는 것이다.

　밀 이전의 공리주의는 양적 쾌락주의의 입장을 가졌다. 이는 모든 쾌락이 측정 가능하고 어디서 나오든 상관없이 똑같은 성질을 가지므로 단지 양에서만 차이가 난다고 보는 것이다. 동물적 욕망에서 나오는 감각적 쾌락과 인간의 지성, 도덕 감정, 상상력 등에서 나오는 정신적 쾌락이 본질적으로 동일하다고 본 것이다.

　하지만 이런 입장은 여러 비판을 받았다. 상대적으로 쉽게 쾌락을 누릴 수 있는 동물이 가장 행복한 존재가 될 수 있기 때문에 '천박한 돼지의 철학'이라는 비판을 받았다. 또한 최대 행복의 추구가 인간의 이기심과 충돌할 수 있어서 실현하기 어렵다는 비판도 있었다. 이런 문제점들을 해결하기 위해 밀은 공리주의 이론을 발전시켰다.

　밀은 질적 쾌락주의를 주장했다. 그는 쾌락이 본래부터 질적 차이가 있다고 봤다. 감각적이고 육체적인 쾌락은 저급 쾌락이고, 정신적 쾌락은 고급 쾌락이다. 고급 쾌락은 저급 쾌락보다 더 바람직하고 가치 있는 우월성을 지닌다고 생각했다. 동물과 달리 인간은 고급 쾌락의 추구를 통해 인간의 품위를 높일 수 있다고 봤다.

　밀 이전의 공리주의는 최대 행복 추구와 이기심이 충돌할 때 법률, 여론 등과 같은 외적 제재가 개인의 이기적 본성을 제어할 수 있다고 생각했다. 하지만 밀은 이것이 근본적인 해결책이 아니라고 봤다. 외적 제재가 최대 행복의 원리에 맞는 행동을 하게 할 수는 있지만, 자발적으로 그런 행동을 하도록 이끄는 힘은 아니라고 생각했기 때문이다. 그래서 밀은 내적 제재인 ㉠\uline{양심}을 강조했다. 양심은 우리 마음 안에서 형성되는 일종의 도덕적 의무감으로, 이를 어기면 내면에 고통을 준다. 양심은 구성원들과 일체감을 이루고자 하는 타고난 사회적 감정에 바탕을 두고, 교육과 외적 제재 등의 후천적인 경험을 통해 기를 수 있다. 이를 통해 비로소 인간은 자기 이익만 추구하는 성향을 극복하고 최대 행복의 원리에 따르는 삶을 실현할 수 있다고 봤다.

　밀은 외적 제재와 내적 제재를 통해 최대 행복의 원리를 실현하여 사회 구성원의 복지를 높일 수 있다고 보았고, 그런 점에서 공리주의가 인간 윤리의 타당한 기준이 될 수 있다고 강조했다.
}


\kfPassageNote{문장 독해 — 본문 문장}{%
1. 공리주의는 공리를 통해 최대 행복을 실현하는 것을 중시하는 이론이다.

2. 여기서 공리란 이익과 유용함을 뜻하며, 공리주의에서 행복이란 고통을 피하고 쾌락을 추구하는 것을 의미한다.

3. 이때 행복은 개인의 쾌락만이 아니라 개인의 행위와 관련된 사회 구성원의 쾌락도 함께 고려하는 것이다.

4. 밀 이전의 공리주의는 양적 쾌락주의의 입장을 가졌다.

5. 이는 모든 쾌락이 측정 가능하고 어디서 나오든 상관없이 똑같은 성질을 가지므로 단지 양에서만 차이가 난다고 보는 것이다.

6. 동물적 욕망에서 나오는 감각적 쾌락과 인간의 지성, 도덕 감정, 상상력 등에서 나오는 정신적 쾌락이 본질적으로 동일하다고 본 것이다.

7. 하지만 이런 입장은 여러 비판을 받았다.

8. 상대적으로 쉽게 쾌락을 누릴 수 있는 동물이 가장 행복한 존재가 될 수 있기 때문에 '천박한 돼지의 철학'이라는 비판을 받았다.

9. 또한 최대 행복의 추구가 인간의 이기심과 충돌할 수 있어서 실현하기 어렵다는 비판도 있었다.

10. 이런 문제점들을 해결하기 위해 밀은 공리주의 이론을 발전시켰다.

11. 밀은 질적 쾌락주의를 주장했다.

12. 그는 쾌락이 본래부터 질적 차이가 있다고 봤다.

13. 감각적이고 육체적인 쾌락은 저급 쾌락이고, 정신적 쾌락은 고급 쾌락이다.

14. 고급 쾌락은 저급 쾌락보다 더 바람직하고 가치 있는 우월성을 지닌다고 생각했다.

15. 동물과 달리 인간은 고급 쾌락의 추구를 통해 인간의 품위를 높일 수 있다고 봤다.

16. 밀 이전의 공리주의는 최대 행복 추구와 이기심이 충돌할 때 법률, 여론 등과 같은 외적 제재가 개인의 이기적 본성을 제어할 수 있다고 생각했다.

17. 하지만 밀은 이것이 근본적인 해결책이 아니라고 봤다.

18. 외적 제재가 최대 행복의 원리에 맞는 행동을 하게 할 수는 있지만, 자발적으로 그런 행동을 하도록 이끄는 힘은 아니라고 생각했기 때문이다.

19. 그래서 밀은 내적 제재인 양심을 강조했다.

20. 양심은 우리 마음 안에서 형성되는 일종의 도덕적 의무감으로, 이를 어기면 내면에 고통을 준다.

21. 양심은 구성원들과 일체감을 이루고자 하는 타고난 사회적 감정에 바탕을 두고, 교육과 외적 제재 등의 후천적인 경험을 통해 기를 수 있다.

22. 이를 통해 비로소 인간은 자기 이익만 추구하는 성향을 극복하고 최대 행복의 원리에 따르는 삶을 실현할 수 있다고 봤다.

23. 밀은 외적 제재와 내적 제재를 통해 최대 행복의 원리를 실현하여 사회 구성원의 복지를 높일 수 있다고 보았고, 그런 점에서 공리주의가 인간 윤리의 타당한 기준이 될 수 있다고 강조했다.
}


\kfActivityBg \par\kfMaybeNeedspace{32\baselineskip}\noindent{\kfTipMark\,\,}{\small\color{kfInk} 다음 표의 '의미'를 읽고, 그에 해당하는 '어휘'를 지문에서 찾아 적으시오.}\par\nopagebreak[4]\smallskip\nopagebreak[4]
\begin{kfVocab}
\kfVocabItem{1}{최대 다수의 최대 행복을 윤리의 기준으로 삼는 사상}
\kfVocabItem{2}{이익과 유용함}
\kfVocabItem{3}{즐겁고 만족스러운 감정}
\kfVocabItem{4}{모든 쾌락이 질적으로 동일하고 양으로만 측정 가능하다고 보는 입장}
\kfVocabItem{5}{시각, 청각 등 감각 기관을 통해 느끼는 즐거움}
\kfVocabItem{6}{지성, 도덕 감정, 상상력 등에서 나오는 즐거움}
\kfVocabItem{7}{자기 자신의 이익만을 생각하고 추구하는 마음}
\kfVocabItem{8}{쾌락이 질적으로 다른 종류가 있다고 보는 입장}
\kfVocabItem{9}{다른 것보다 더 뛰어나고 나은 성질}
\kfVocabItem{10}{사람이 갖추어야 할 위엄이나 기품}
\kfVocabItem{11}{법률, 여론 등 외부에서 이기심을 통제하는 장치}
\kfVocabItem{12}{양심처럼 자기 마음속에서 우러나오는 도덕적 통제}
\kfVocabItem{13}{어떤 행위에 대하여 옳고 그름을 판단하는 내면의 의식}
\end{kfVocab}

\kfActivityBg \par\needspace{15\baselineskip}
\kfSecHeader{kfAreaNonlit}{활동}{문장 독해}
\par\kfMaybeNeedspace{32\baselineskip}\noindent{\kfTipMark\,\,}{\small\color{kfInk} 다음은 본문의 번호 매긴 문장입니다. 각 문장을 읽고 물음에 답하시오.}\par\nopagebreak[4]\smallskip\nopagebreak[4]
\kfBeginMC
\begin{kfSentenceUnit}{1}{}
\kfSentQ{1.}{공리주의가 중시하는 것은 무엇인가?}
\begin{kfChoices}
\kfChoice 개인의 이익만 추구하는 것
\kfChoice 법률에 의지해 도덕을 지키는 것
\kfChoice 고통 없이 쾌락만 영원히 지속되는 것
\kfChoice 공리를 통해 최대 행복을 실현하는 것
\end{kfChoices}
\end{kfSentenceUnit}

\begin{kfSentenceUnit}{2}{}
\kfSentQ{2.}{이 글에서 '공리'는 어떤 의미인가?}
\begin{kfChoices}
\kfChoice 희생과 봉사
\kfChoice 법률과 제도
\kfChoice 이익과 유용함
\kfChoice 도덕적 의무감
\end{kfChoices}
\end{kfSentenceUnit}

\begin{kfSentenceUnit}{3}{}
\kfSentQ{3.}{공리주의에서 행복이란 무엇을 피하고 무엇을 추구하는 것을 의미하는가?}
\begin{kfChoices}
\kfChoice 고통을 피하고 명예를 추구하는 것
\kfChoice 쾌락을 피하고 고통을 추구하는 것
\kfChoice 고통을 피하고 쾌락을 추구하는 것
\kfChoice 이기심을 피하고 법률을 추구하는 것
\end{kfChoices}
\end{kfSentenceUnit}

\begin{kfSentenceUnit}{4}{}
\kfSentQ{4.}{공리주의에서 행복은 누구의 쾌락까지 고려해야 하는가?}
\begin{kfChoices}
\kfChoice 나 자신
\kfChoice 국가 지도자
\kfChoice 가족 구성원
\kfChoice 사회 구성원
\end{kfChoices}
\end{kfSentenceUnit}

\begin{kfSentenceUnit}{5}{}
\kfSentQ{5.}{밀 이전의 공리주의는 쾌락에 대해 어떤 입장이었는가?}
\begin{kfChoices}
\kfChoice 양적 쾌락주의
\kfChoice 질적 쾌락주의
\kfChoice 이기적 쾌락주의
\kfChoice 감각적 쾌락주의
\end{kfChoices}
\end{kfSentenceUnit}

\begin{kfSentenceUnit}{6}{}
\kfSentQ{6.}{모든 쾌락은 양에서만 차이가 날 뿐 질적으로는 동일하다고 본 밀 이전의 공리주의 입장은 무엇인가?}
\begin{kfChoices}
\kfChoice 양적 쾌락주의
\kfChoice 질적 쾌락주의
\kfChoice 이기적 쾌락주의
\kfChoice 감각적 쾌락주의
\end{kfChoices}
\end{kfSentenceUnit}

\begin{kfSentenceUnit}{7}{}
\kfSentQ{7.}{양적 쾌락주의가 쾌락의 차이를 판단하는 유일한 기준은 무엇인가?}
\begin{kfChoices}
\kfChoice 양
\kfChoice 질
\kfChoice 기간
\kfChoice 출처
\end{kfChoices}
\end{kfSentenceUnit}

\begin{kfSentenceUnit}{8}{}
\kfSentQ{8.}{양적 쾌락주의는 감각적 쾌락과 정신적 쾌락이 본질적으로 어떻다고 보았는가?}
\begin{kfChoices}
\kfChoice 다르다
\kfChoice 무관하다
\kfChoice 동일하다
\kfChoice 상반된다
\end{kfChoices}
\end{kfSentenceUnit}

\begin{kfSentenceUnit}{9}{}
\kfSentQ{9.}{'이런 입장'이 가리키는 것은 무엇인가?}
\begin{kfChoices}
\kfChoice 질적 쾌락주의
\kfChoice 양적 쾌락주의
\kfChoice 이기적 쾌락주의
\kfChoice 감각적 쾌락주의
\end{kfChoices}
\end{kfSentenceUnit}

\begin{kfSentenceUnit}{10}{}
\kfSentQ{10.}{'이런 입장(양적 쾌락주의)'이 받은 비판으로 가장 적절한 것은?}
\begin{kfChoices}
\kfChoice 천박한 돼지의 철학
\kfChoice 현실을 무시한 이상주의
\kfChoice 감정을 배제한 합리주의
\kfChoice 고상한 소크라테스의 철학
\end{kfChoices}
\end{kfSentenceUnit}

\begin{kfSentenceUnit}{11}{}
\kfSentQ{11.}{양적 쾌락주의가 받은 비판의 별칭은 무엇인가?}
\begin{kfChoices}
\kfChoice 소크라테스의 철학
\kfChoice 배부른 돼지의 철학
\kfChoice 천박한 돼지의 철학
\kfChoice 고상한 인간의 철학
\end{kfChoices}
\end{kfSentenceUnit}

\begin{kfSentenceUnit}{12}{}
\kfSentQ{12.}{최대 행복의 추구가 무엇과 충돌할 수 있다는 비판이 있었는가?}
\begin{kfChoices}
\kfChoice 인간의 이타심
\kfChoice 인간의 이기심
\kfChoice 인간의 동정심
\kfChoice 인간의 자비심
\end{kfChoices}
\end{kfSentenceUnit}

\begin{kfSentenceUnit}{13}{}
\kfSentQ{13.}{밀이 공리주의 이론을 발전시킨 이유는 무엇인가?}
\begin{kfChoices}
\kfChoice 새로운 법률을 제정하기 위해
\kfChoice 종교적 교리를 전파하기 위해
\kfChoice 개인의 이기심을 옹호하기 위해
\kfChoice 기존 이론의 문제점을 해결하기 위해
\end{kfChoices}
\end{kfSentenceUnit}

\begin{kfSentenceUnit}{14}{}
\kfSentQ{14.}{밀이 새롭게 주장한 공리주의 이론은 무엇인가?}
\begin{kfChoices}
\kfChoice 양적 쾌락주의
\kfChoice 질적 쾌락주의
\kfChoice 이기적 쾌락주의
\kfChoice 감각적 쾌락주의
\end{kfChoices}
\end{kfSentenceUnit}

\begin{kfSentenceUnit}{15}{}
\kfSentQ{15.}{밀이 쾌락에 질적인 차이가 있다고 주장한 이론은 무엇인가?}
\begin{kfChoices}
\kfChoice 질적 쾌락주의
\kfChoice 양적 쾌락주의
\kfChoice 감각 쾌락주의
\kfChoice 정신 쾌락주의
\end{kfChoices}
\end{kfSentenceUnit}

\begin{kfSentenceUnit}{16}{}
\kfSentQ{16.}{밀은 쾌락에 어떤 차이가 있다고 보았는가?}
\begin{kfChoices}
\kfChoice 질적 차이
\kfChoice 양적 차이
\kfChoice 시간적 차이
\kfChoice 공간적 차이
\end{kfChoices}
\end{kfSentenceUnit}

\begin{kfSentenceUnit}{17}{}
\kfSentQ{17.}{밀이 말한 두 가지 종류의 쾌락은 무엇인가?}
\begin{kfChoices}
\kfChoice 순간 쾌락, 영원 쾌락
\kfChoice 육체 쾌락, 물질 쾌락
\kfChoice 저급 쾌락, 고급 쾌락
\kfChoice 개인 쾌락, 사회 쾌락
\end{kfChoices}
\end{kfSentenceUnit}

\begin{kfSentenceUnit}{18}{}
\kfSentQ{18.}{밀에 따르면, 정신적 쾌락과 같이 더 바람직하고 가치 있는 쾌락은 무엇인가?}
\begin{kfChoices}
\kfChoice 저급 쾌락
\kfChoice 순간 쾌락
\kfChoice 육체 쾌락
\kfChoice 고급 쾌락
\end{kfChoices}
\end{kfSentenceUnit}

\begin{kfSentenceUnit}{19}{}
\kfSentQ{19.}{밀은 고급 쾌락이 저급 쾌락보다 어떤 면에서 뛰어나다고 생각했는가?}
\begin{kfChoices}
\kfChoice 우월성
\kfChoice 양적 우위
\kfChoice 간편성
\kfChoice 대중성
\end{kfChoices}
\end{kfSentenceUnit}

\begin{kfSentenceUnit}{20}{}
\kfSentQ{20.}{밀은 인간이 고급 쾌락을 추구함으로써 무엇을 높일 수 있다고 보았는가?}
\begin{kfChoices}
\kfChoice 품위
\kfChoice 재산
\kfChoice 권력
\kfChoice 사치
\end{kfChoices}
\end{kfSentenceUnit}

\begin{kfSentenceUnit}{21}{}
\kfSentQ{21.}{최대 행복 추구와 이기심이 충돌할 때, 밀 이전의 공리주의가 제시한 해결책은 무엇인가?}
\begin{kfChoices}
\kfChoice 외적 제재
\kfChoice 내적 제재
\kfChoice 종교 신앙
\kfChoice 도덕 교육
\end{kfChoices}
\end{kfSentenceUnit}

\begin{kfSentenceUnit}{22}{}
\kfSentQ{22.}{법률이나 여론처럼 외부에서 개인의 이기심을 제어하는 제재를 무엇이라고 하는가?}
\begin{kfChoices}
\kfChoice 사적 제재
\kfChoice 공적 제재
\kfChoice 내적 제재
\kfChoice 외적 제재
\end{kfChoices}
\end{kfSentenceUnit}

\begin{kfSentenceUnit}{23}{}
\kfSentQ{23.}{밀은 외적 제재를 어떤 해결책이라고 평가했는가?}
\begin{kfChoices}
\kfChoice 유일한 해결책이다
\kfChoice 불필요한 해결책이다
\kfChoice 가장 완벽한 해결책이다
\kfChoice 근본적인 해결책이 아니다
\end{kfChoices}
\end{kfSentenceUnit}

\begin{kfSentenceUnit}{24}{}
\kfSentQ{24.}{밀은 외적 제재가 사람들을 자발적으로 행동하게 만드는 힘이라고 생각했는가?}
\begin{kfChoices}
\kfChoice 힘이다
\kfChoice 힘이 아니다
\kfChoice 알 수 없다
\kfChoice 확신할 수 없다
\end{kfChoices}
\end{kfSentenceUnit}

\begin{kfSentenceUnit}{25}{}
\kfSentQ{25.}{밀이 외적 제재의 대안으로 강조한 '내적 제재'는 무엇인가?}
\begin{kfChoices}
\kfChoice 여론
\kfChoice 양심
\kfChoice 처벌
\kfChoice 법률
\end{kfChoices}
\end{kfSentenceUnit}

\begin{kfSentenceUnit}{26}{}
\kfSentQ{26.}{밀이 이기심을 극복할 근본적인 해결책으로 강조한 내적 제재는 무엇인가?}
\begin{kfChoices}
\kfChoice 여론
\kfChoice 양심
\kfChoice 처벌
\kfChoice 법률
\end{kfChoices}
\end{kfSentenceUnit}

\begin{kfSentenceUnit}{27}{}
\kfSentQ{27.}{밀에 따르면 '양심'이란 우리 마음 안에서 형성되는 무엇인가?}
\begin{kfChoices}
\kfChoice 이기적 욕망
\kfChoice 사회적 명예
\kfChoice 본능적 공포
\kfChoice 도덕적 의무감
\end{kfChoices}
\end{kfSentenceUnit}

\begin{kfSentenceUnit}{28}{}
\kfSentQ{28.}{양심이 바탕을 두고 있는 타고난 감정은 무엇인가?}
\begin{kfChoices}
\kfChoice 질투의 감정
\kfChoice 공포의 감정
\kfChoice 배타적 감정
\kfChoice 사회적 감정
\end{kfChoices}
\end{kfSentenceUnit}

\begin{kfSentenceUnit}{29}{}
\kfSentQ{29.}{밀에 따르면, 양심의 바탕이 되는 타고난 감정은 무엇인가?}
\begin{kfChoices}
\kfChoice 개인적 감정
\kfChoice 이기적 감정
\kfChoice 배타적 감정
\kfChoice 사회적 감정
\end{kfChoices}
\end{kfSentenceUnit}

\begin{kfSentenceUnit}{30}{}
\kfSentQ{30.}{밀은 양심을 통해 인간이 어떤 성향을 극복할 수 있다고 보았는가?}
\begin{kfChoices}
\kfChoice 고통을 피하려는 성향
\kfChoice 질서를 지키려는 성향
\kfChoice 남을 돕고자 하는 성향
\kfChoice 자기 이익만 추구하는 성향
\end{kfChoices}
\end{kfSentenceUnit}

\begin{kfSentenceUnit}{31}{}
\kfSentQ{31.}{밀이 생각한, 최대 행복의 원리를 실현하는 두 가지 방법은 무엇인가?}
\begin{kfChoices}
\kfChoice 외적 제재와 내적 제재
\kfChoice 법률 강화와 처벌 확대
\kfChoice 물질적 보상과 명예 부여
\kfChoice 개인의 희생과 사회적 강요
\end{kfChoices}
\end{kfSentenceUnit}

\kfEndMC

\kfActivityBg \par\kfMaybeNeedspace{32\baselineskip}\noindent{\kfTipMark\,\,}{\small\color{kfInk} 글의 전체적인 논리 흐름을 파악하며 빈칸에 알맞은 내용을 채워 보세요.}\par\nopagebreak[4]\smallskip\nopagebreak[4]
\begin{kfStructure}
\kfNode{0}{}{}
\end{kfStructure}

\kfSecHeader{kfAreaNonlit}{문제}{}

\par\medskip
\kfBeginMC
\begin{kfQBlock}{01}{객관식}
\kfQStem{윗글에 대한 설명으로 적절하지 \uline{\underline{않은}} 것은?}
\begin{kfChoices}
\kfChoice 공리주의와 다른 윤리 사상 사이의 차이점을 항목별로 비교 분석한다.
\kfChoice 밀 이전 공리주의가 지닌 문제점을 짚어 준다.
\kfChoice 밀의 ‘질적 쾌락주의’와 ‘내적 제재’가 어떤 발전이었는지 설명한다.
\kfChoice 밀에 의해 공리주의 이론이 심화되었음을 보여 준다.
\kfChoice 공리주의의 기본 개념을 먼저 소개한다.
\end{kfChoices}
\end{kfQBlock}
\begin{kfQBlock}{02}{객관식}
\kfQStem{윗글의 내용과 일치하지 \uline{않는} 것은?}
\begin{kfChoices}
\kfChoice 공리주의는 사회 전체의 행복을 고려하는 이론이다.
\kfChoice 밀은 모든 쾌락이 본질적으로 동일하다고 생각했다.
\kfChoice '천박한 돼지의 철학'은 양적 쾌락주의에 대한 비판이다.
\kfChoice 밀은 외적 제재만으로는 이기심 문제를 근본적으로 해결할 수 없다고 보았다.
\kfChoice 밀에 따르면 양심은 교육을 통해 길러질 수 있다.
\end{kfChoices}
\end{kfQBlock}
\begin{kfQBlock}{03}{객관식}
\kfQStem{밀과 밀 이전 공리주의의 입장을 비교한 것으로 적절하지 \uline{\underline{않은}} 것은?}
\begin{kfChoices}
\kfChoice 쾌락: 밀은 질적 차이를, 이전에는 양적 차이만을 인정했다.
\kfChoice 인간관: 밀은 인간의 품위를, 이전에는 감각적 쾌락을 중시했다.
\kfChoice 이기심 제어: 밀은 내적 제재를, 이전에는 외적 제재를 강조했다.
\kfChoice 행복: 밀은 개인의 행복만을, 이전에는 사회 전체의 행복을 추구했다.
\kfChoice 가치 판단: 밀은 고급 쾌락을, 이전에는 양적으로 많은 쾌락을 더 가치있게 봤다.
\end{kfChoices}
\end{kfQBlock}
\begin{kfQBlock}{04}{객관식}
\kfQStem{㉠ ‘양심’에 대한 밀의 설명에서 추론할 수 있는 내용으로 적절하지 \uline{\underline{않은}} 것은?}
\begin{kfChoices}
\kfChoice 양심은 개인의 이익과 사회 전체의 이익을 조화시키는 역할을 한다.
\kfChoice 사람마다 타고나는 것이므로 어떤 후천적 노력으로도 바꿀 수 없다.
\kfChoice 최대 행복의 원리에 따라 살도록 행위 주체를 이끈다.
\kfChoice 외적 제재가 아닌 내면에서 작동하는 ‘내적 제재’에 해당한다.
\kfChoice 자기 이익만 추구하는 성향을 극복하는 데 도움을 준다.
\end{kfChoices}
\end{kfQBlock}
\begin{kfQBlock}{05}{단답형}
\kfQStem{공리주의에서 행복이란 무엇을 피하고 무엇을 추구하는 것을 의미하는가?}
\kfAnswerLines{1}
\end{kfQBlock}
\begin{kfQBlock}{06}{단답형}
\kfQStem{모든 쾌락은 양에서만 차이가 날 뿐 질적으로는 동일하다고 본 밀 이전의 공리주의 입장은 무엇인가?}
\kfAnswerLines{1}
\end{kfQBlock}
\begin{kfQBlock}{07}{단답형}
\kfQStem{양적 쾌락주의가 받은 비판의 별칭은 무엇인가?}
\kfAnswerLines{1}
\end{kfQBlock}
\begin{kfQBlock}{08}{단답형}
\kfQStem{밀이 쾌락에 질적인 차이가 있다고 주장한 이론은 무엇인가?}
\kfAnswerLines{1}
\end{kfQBlock}
\begin{kfQBlock}{09}{단답형}
\kfQStem{밀에 따르면, 정신적 쾌락과 같이 더 바람직하고 가치 있는 쾌락은 무엇인가?}
\kfAnswerLines{1}
\end{kfQBlock}
\begin{kfQBlock}{10}{서술형}
\kfQStem{밀이 기존의 공리주의가 받은 비판을 극복하기 위해 제시한 두 가지 핵심 주장을 <조건>에 맞게 서술하시오.}
\kfAnswerNote{답안}{4}
\end{kfQBlock}
\kfEndMC



\clearpage

\kfChapterCover{2}{러셀2 (챕터2)}{Chapter 2}{이번 챕터의 학습 내용을 시작합니다.}
\begin{kfChapterAreaList}
\kfChapterAreaItem{문법}{kfAreaGrammar}{문장의 핵심 재료인 체언, 문법적 관계를 맺는 관계언}
\kfChapterAreaItem{개념}{kfAreaConcept}{소설, 지금 누가 이야기하고 있을까?}
\kfChapterAreaItem{문학}{kfAreaLiterature}{「동백꽃」 — 김유정}
\kfChapterAreaItem{비문학}{kfAreaNonlit}{[사회] 미성년자 보호를 위한 계약 취소 제도}
\end{kfChapterAreaList}

\kfStartGrammar{문장의 핵심 재료인 체언, 문법적 관계를 맺는 관계언}


\kfPassageNoteNT{%
　우리가 사용하는 수많은 단어는 문장 속에서 맡은 역할과 성질에 따라 몇 개의 그룹으로 나눌 수 있는데, 이를 ‘품사’라고 한다. 품사 분류는 단어의 기능, 의미, 형태라는 세 가지 기준에 따라 이루어진다. 그중에서도 문장에 들어갈 핵심 내용물을 담고 있는 단어 그룹이 있다. 바로 ‘몸 체(體)’ 자를 쓰는 체언이다. 체언은 문장이라는 요리의 맛과 정체를 결정하는 ‘핵심 재료’와 같다고 할 수 있다. 이러한 체언에는 구체적인 대상의 이름을 나타내는 명사, 그 이름을 대신하는 대명사, 그리고 수량이나 순서를 가리키는 수사가 속한다.

　명사는 세상 모든 존재에 붙여진 이름이다. ‘이순신’, ‘서울’처럼 세상에 단 하나뿐인 대상을 가리키는 고유 명사가 있고, ‘사람’, ‘도시’처럼 공통된 속성을 지닌 대상 전체를 아우르는 보통 명사가 있다. 또한 ‘책상’처럼 홀로 쓰일 수 있는 자립 명사와 달리, ‘것’, ‘수’, ‘뿐’처럼 반드시 앞에 꾸며주는 말이 있어야만 의미를 갖는 의존 명사도 존재한다.

　대명사는 명사를 대신하여 가리키는 말로, 반복을 피하고 문장을 간결하게 만든다. ‘나, 너, 우리, 그’처럼 사람을 가리키는 인칭 대명사와 ‘이것, 저기, 거기’처럼 사물이나 장소를 가리키는 지시 대명사로 나뉜다.

　수사는 사물의 수량이나 순서를 나타내는 말로, ‘하나, 둘’처럼 양을 세는 양수사와 ‘첫째, 둘째’처럼 순서를 나타내는 서수사가 있다.

　그런데 요리의 재료들을 그냥 늘어놓는다고 음식이 완성되지 않듯이, 이 체언들도 혼자서는 문장 속에서 다른 단어와 완전한 관계를 맺기 어렵다. 이때 등장하는 것이 바로 관계언에 속하는 조사이다. 조사는 체언 뒤에 껌딱지처럼 달라붙어 그 단어가 문장 속에서 어떤 자격을 갖는지 표시해 주거나 특별한 의미를 더해주는 역할을 한다.

　예를 들어 ‘학생’이라는 명사에 붙어 주어의 자격을 주는 ‘-이/가’는 격 조사이고, ‘너만 와라’에서처럼 특별한 의미를 더하는 ‘-만’은 보조사이며, ‘너와 나’처럼 둘을 이어주는 ‘-와/과’는 접속 조사이다. 이처럼 조사는 주로 체언이라는 핵심 재료에 붙어 문법적 관계를 명시해주는 중요한 양념이자 접착제 역할을 담당한다.
\par\medskip\begin{itemize}[leftmargin=18pt, itemsep=2pt, topsep=2pt, label=\textbullet]
\item 체언 — 고유 명사(이순신)·보통 명사(도시)·의존 명사(것·수·뿐)
\item 대명사·수사 — '나·이것(대명)' / '하나·첫째(수사)'
\item 조사 — 격 조사(이/가)·보조사(만·도)·접속 조사(와/과)
\end{itemize}
}

\kfActivityBg \par\kfMaybeNeedspace{32\baselineskip}\noindent{\kfTipMark\,\,}{\small\color{kfInk} 지문의 핵심 내용을 떠올리며 빈칸에 알맞은 말을 채워 보세요.}\par\nopagebreak[4]\smallskip\nopagebreak[4]

\kfActivityBg \kfSecHeader{kfAreaGrammar}{활동}{분석 훈련}
\par\kfMaybeNeedspace{24\baselineskip}\noindent{\kfTipMark\,\,}{\small\color{kfInk} 다음 문장을 읽고, 설명에 해당하는 단어를 모두 찾아 써 보세요.}\par\nopagebreak[4]\smallskip\nopagebreak[4]
\begin{enumerate}[leftmargin=22pt, itemsep=6pt, topsep=4pt, label=\textbf{\arabic*.}]
\item 다음 문장에서 사물이나 사람의 이름을 나타내는 명사를 모두 찾아 쓰시오.
\begin{kfEvidence}철수가 학교에서 친구와 축구를 했다.\end{kfEvidence}
\item 다음 문장에서 명사를 대신하여 쓰인 대명사를 모두 찾아 쓰시오.
\begin{kfEvidence}이것은 바로 우리가 찾던 그것이다.\end{kfEvidence}
\item 다음 문장에서 수량이나 순서를 나타내는 수사를 모두 찾아 쓰시오.
\begin{kfEvidence}사과 하나와 배 둘을 샀고, 달리기에서는 첫째로 들어왔다.\end{kfEvidence}
\item 다음 문장에서 체언 뒤에 붙어 문법적 관계를 나타내거나 의미를 더해주는 조사를 모두 찾아 쓰시오.
\begin{kfEvidence}나도 너처럼 할 수만 있다면 좋겠다.\end{kfEvidence}
\item 다음 문장에서 문장의 뼈대 역할을 하는 체언을 모두 찾아 쓰시오.
\begin{kfEvidence}우리 셋이 힘을 합쳐 저 문제를 풀었다.\end{kfEvidence}
\item 다음 문장에서 이름을 나타내는 명사를 모두 찾아 쓰시오.
\begin{kfEvidence}소년은 푸른 바다에서 꿈과 희망을 보았다.\end{kfEvidence}
\item 다음 문장에서 명사를 대신하는 대명사를 모두 찾아 쓰시오.
\begin{kfEvidence}그는 누구에게 어디로 가는지 묻지 않았다.\end{kfEvidence}
\item 다음 문장에서 수량을 나타내는 수사를 찾아 쓰시오.
\begin{kfEvidence}다섯 사람이 길을 걷고 있다.\end{kfEvidence}
\item 다음 문장에서 문법적 관계를 표시하는 조사를 모두 찾아 쓰시오.
\begin{kfEvidence}아빠가 동생에게 과자를 주셨다.\end{kfEvidence}
\item 다음 문장에서 문장의 주체가 되는 체언을 모두 찾아 쓰시오.
\begin{kfEvidence}이것은 나의 조카가 가장 좋아하는 책이다.\end{kfEvidence}
\end{enumerate}

\kfSecHeader{kfAreaGrammar}{문제}{}

\par\medskip
\kfBeginMC
\begin{kfQBlock}{01}{객관식}
\kfQStem{다음 중 품사가 \uline{다른} 하나는?}
\begin{kfChoices}
\kfChoice 행복
\kfChoice 우정
\kfChoice 우리
\kfChoice 학교
\kfChoice 연필
\end{kfChoices}
\end{kfQBlock}
\begin{kfQBlock}{02}{객관식}
\kfQStem{체언에 대한 설명으로 적절하지 \uline{\underline{않은}} 것은?}
\begin{kfChoices}
\kfChoice 어미와 결합해 다양한 형태로 활용되며 서술어 역할을 한다.
\kfChoice 명사·대명사·수사가 이에 속한다.
\kfChoice 주로 조사와 결합하여 주어·목적어 등 문장 성분이 된다.
\kfChoice 형태가 변하지 않는 ‘불변어’에 속한다.
\kfChoice 문장의 뼈대를 이루는 핵심 품사이다.
\end{kfChoices}
\end{kfQBlock}
\begin{kfQBlock}{03}{객관식}
\kfQStem{다음 밑줄 친 단어의 품사가 명사가 \uline{\underline{아닌}} 것은?}
\begin{kfChoices}
\kfChoice 나는 \uline{사과}를 좋아한다.
\kfChoice \uline{이것}은 내 책이다.
\kfChoice \uline{사랑}은 위대하다.
\kfChoice \uline{철수}는 내 친구다.
\kfChoice \uline{컴퓨터}가 고장 났다.
\end{kfChoices}
\end{kfQBlock}
\begin{kfQBlock}{04}{객관식}
\kfQStem{다음 중 수사가 포함된 문장은?}
\begin{kfChoices}
\kfChoice 나는 너를 좋아한다.
\kfChoice 우리는 친구이다.
\kfChoice 사과가 다섯 개 있다.
\kfChoice 저 산이 매우 높다.
\kfChoice 어서 집에 가자.
\end{kfChoices}
\end{kfQBlock}
\begin{kfQBlock}{05}{객관식}
\kfQStem{밑줄 친 단어와 품사가 같은 것은?}
\begin{kfEvidence}\uline{그}는 학생이다.\end{kfEvidence}
\begin{kfChoices}
\kfChoice \uline{새} 옷을 입었다.
\kfChoice \uline{나}는 밥을 먹었다.
\kfChoice \uline{빨리} 달린다.
\kfChoice \uline{와}, 정말 예쁘다!
\kfChoice \uline{하늘}이 맑다.
\end{kfChoices}
\end{kfQBlock}
\begin{kfQBlock}{06}{객관식}
\kfQStem{다음 중 홀로 쓰일 수 \uline{\underline{없는}} 품사는?}
\begin{kfChoices}
\kfChoice 명사
\kfChoice 대명사
\kfChoice 수사
\kfChoice 조사
\kfChoice 동사
\end{kfChoices}
\end{kfQBlock}
\begin{kfQBlock}{07}{객관식}
\kfQStem{다음 문장에서 조사는 모두 몇 개인가?}
\begin{kfEvidence}너는 학교에서 친구랑 공부를 했다.\end{kfEvidence}
\begin{kfChoices}
\kfChoice 2개
\kfChoice 3개
\kfChoice 4개
\kfChoice 5개
\kfChoice 6개
\end{kfChoices}
\end{kfQBlock}
\begin{kfQBlock}{08}{객관식}
\kfQStem{다음 <보기>의 ㉠\textasciitilde{}㉢에 해당하는 품사를 바르게 짝지은 것은?}
\begin{kfEvidence}㉠\uline{이것}은 ㉡\uline{나}의 ㉢\uline{하나}뿐인 보물이다.\end{kfEvidence}
\begin{kfChoices}
\kfChoice ㉠:명사, ㉡:대명사, ㉢:수사
\kfChoice ㉠:대명사, ㉡:대명사, ㉢:수사
\kfChoice ㉠:대명사, ㉡:명사, ㉢:조사
\kfChoice ㉠:명사, ㉡:대명사, ㉢:조사
\kfChoice ㉠:대명사, ㉡:대명사, ㉢:명사
\end{kfChoices}
\end{kfQBlock}
\begin{kfQBlock}{09}{객관식}
\kfQStem{다음 중 대명사가 \uline{\underline{없는}} 문장은?}
\begin{kfChoices}
\kfChoice 저기 우리 집이 보인다.
\kfChoice 그녀는 무엇을 찾고 있니?
\kfChoice 아빠가 동생을 불렀다.
\kfChoice 거기에는 아무도 없었다.
\kfChoice 너는 학생이 아니다.
\end{kfChoices}
\end{kfQBlock}
\begin{kfQBlock}{10}{객관식}
\kfQStem{다음 밑줄 친 단어의 품사가 \uline{다른} 하나는?}
\begin{kfChoices}
\kfChoice \uline{첫째} 아들
\kfChoice \uline{둘} 더하기 둘은 넷
\kfChoice \uline{셋} 셀 동안 나와
\kfChoice \uline{넷}이서 놀러 갔다
\kfChoice 사과 \uline{다섯}
\end{kfChoices}
\end{kfQBlock}
\begin{kfQBlock}{11}{객관식}
\kfQStem{다음 중 체언이 \uline{\underline{아닌}} 것은?}
\begin{kfChoices}
\kfChoice 컴퓨터
\kfChoice 그들
\kfChoice 백(百)
\kfChoice 빨리
\kfChoice 우정
\end{kfChoices}
\end{kfQBlock}
\begin{kfQBlock}{12}{객관식}
\kfQStem{다음 <보기>의 설명에 해당하는 품사들의 묶음은?}
\begin{kfEvidence}문장의 뼈대를 이루며, 주로 주어, 목적어 등으로 사용된다.\end{kfEvidence}
\begin{kfChoices}
\kfChoice 명사, 동사
\kfChoice 명사, 대명사, 수사
\kfChoice 동사, 형용사
\kfChoice 관형사, 부사
\kfChoice 조사, 감탄사
\end{kfChoices}
\end{kfQBlock}
\begin{kfQBlock}{13}{객관식}
\kfQStem{다음 문장에서 밑줄 친 단어의 품사는?}
\begin{kfEvidence}나는 밥\uline{도} 먹고 빵\uline{도} 먹었다.\end{kfEvidence}
\begin{kfChoices}
\kfChoice 명사
\kfChoice 대명사
\kfChoice 수사
\kfChoice 조사
\kfChoice 부사
\end{kfChoices}
\end{kfQBlock}
\begin{kfQBlock}{14}{객관식}
\kfQStem{품사의 분류가 바르게 된 것은?}
\begin{kfChoices}
\kfChoice 그 - 명사
\kfChoice 일곱 - 대명사
\kfChoice 에게 - 조사
\kfChoice 사랑 - 수사
\kfChoice 여기 - 명사
\end{kfChoices}
\end{kfQBlock}
\begin{kfQBlock}{15}{객관식}
\kfQStem{다음 중 관계언에 해당하는 단어는?}
\begin{kfChoices}
\kfChoice 첫째
\kfChoice 어디
\kfChoice 나무
\kfChoice 까지
\kfChoice 생각
\end{kfChoices}
\end{kfQBlock}
\begin{kfQBlock}{16}{단답형}
\kfQStem{‘책상’, ‘사랑’, ‘대한민국’처럼 사람이나 사물의 이름을 나타내는 품사는 무엇인가?}
\kfAnswerLines{1}
\end{kfQBlock}
\begin{kfQBlock}{17}{단답형}
\kfQStem{‘나’, ‘너’, ‘이것’, ‘저기’처럼 명사를 대신하여 가리키는 품사는 무엇인가?}
\kfAnswerLines{1}
\end{kfQBlock}
\begin{kfQBlock}{18}{단답형}
\kfQStem{‘하나’, ‘둘’, ‘첫째’처럼 수량이나 순서를 나타내는 품사는 무엇인가?}
\kfAnswerLines{1}
\end{kfQBlock}
\begin{kfQBlock}{19}{단답형}
\kfQStem{문장에서 주로 몸통 역할을 하는 명사, 대명사, 수사를 통틀어 무엇이라고 하는가?}
\kfAnswerLines{1}
\end{kfQBlock}
\begin{kfQBlock}{20}{단답형}
\kfQStem{체언 뒤에 붙어 문법적 관계를 맺어주거나 특별한 뜻을 더해주는 품사는 무엇인가?}
\kfAnswerLines{1}
\end{kfQBlock}
\begin{kfQBlock}{21}{단답형}
\kfQStem{주로 체언 뒤에 붙어 문법적 관계를 표시하는 조사가 기능상 속하는 품사의 갈래는 무엇인가?}
\kfAnswerLines{1}
\end{kfQBlock}
\begin{kfQBlock}{22}{단답형}
\kfQStem{‘너와 나’에서 ‘너’와 ‘나’의 품사는 무엇인가?}
\kfAnswerLines{1}
\end{kfQBlock}
\begin{kfQBlock}{23}{단답형}
\kfQStem{‘하늘이 푸르다’에서 ‘이’의 품사는 무엇인가?}
\kfAnswerLines{1}
\end{kfQBlock}
\begin{kfQBlock}{24}{단답형}
\kfQStem{문장 '나는 책을 읽는다'에서 ‘책’의 품사는 무엇인가?}
\kfAnswerLines{1}
\end{kfQBlock}
\begin{kfQBlock}{25}{단답형}
\kfQStem{문장 ‘사과가 열이 있다.’에서 ‘열’의 품사는 무엇인가?}
\kfAnswerLines{1}
\end{kfQBlock}
\begin{kfQBlock}{26}{서술형}
\kfQStem{다음 문장에서 체언에 해당하는 단어를 모두 찾고, 각각의 품사(명사, 대명사, 수사)를 쓰시오.}
\begin{kfEvidence}우리 둘이 저 산에 올라가 보자.\end{kfEvidence}
\kfAnswerNote{답안}{4}
\end{kfQBlock}
\kfEndMC


\kfStartConcept{소설, 지금 누가 이야기하고 있을까?}


\kfPassageNoteNT{%
　소설은 작가가 창조한 허구의 이야기이다. 작가는 이 이야기를 독자에게 직접 말하지 않고, 자신 대신 이야기를 전달하는 '서술자'를 내세운다. 이 서술자는 독자에게 사건을 전달하며 소설의 전체적인 분위기를 이끌어간다.

　이때 서술자가 어디에 서서, 무엇을 얼마만큼 알고 이야기하느냐에 따라 독자가 받아들이는 정보와 느낌은 완전히 달라진다. 이처럼 서술자가 사건을 바라보는 위치나 관점을 '시점'이라고 한다. 시점은 크게 서술자가 작품 안에 있는 1인칭 시점과 작품 밖에 있는 3인칭 시점으로 나뉜다.

　1인칭 시점은 작품 속 등장인물인 '나'가 직접 서술하는 방식이다. 이 중 '나'가 주인공이 아닌 주변 인물인 경우가 1인칭 관찰자 시점이다. 주요섭의 「사랑손님과 어머니」에서 서술자인 '나'는 어린아이 '옥희'이다. '옥희'는 주인공인 '어머니'와 '사랑손님'의 모습을 자신이 관찰한 것만 독자에게 전달한다. 따라서 독자는 '옥희'의 순수한 눈을 통해 두 주인공의 내면을 짐작해야 한다.

　반면, 서술자인 '나'가 이야기의 주인공이라면 1인칭 주인공 시점이 된다. 김유정의 「동백꽃」이 대표적이다. 이 작품의 서술자인 '나'는 점순이와의 갈등을 겪는 주인공으로서, 자신이 겪은 일과 속마음을 직접 고백한다. 이처럼 1인칭 시점은 '나'의 목소리를 직접 듣기 때문에 독자에게 친밀감과 신뢰감을 주지만, 독자는 오직 '나'가 보거나 아는 것만 알 수 있다는 한계가 있다.

　3인칭 시점은 서술자가 작품 밖에 존재하며 인물을 이름이나 '그', '그녀'와 같은 3인칭 대명사로 지칭하는 방식이다. 이 중 3인칭 관찰자 시점은 서술자가 인물의 내면을 전혀 드러내지 않고, 오직 밖으로 드러나는 행동과 대화만을 객관적으로 묘사하는 방식이다. 황순원의 「학」이 대표적이다. 이 작품의 서술자는 주인공 성삼과 덕재의 속마음을 직접 말해주지 않고, 두 사람의 대화와 행동만을 담담하게 전달한다. 이 시점은 인물의 속마음을 독자가 직접 추측하게 만들어 극적인 긴장감을 높이는 효과가 있다.

　같은 3인칭 시점이라도 서술자가 신처럼 모든 것을 아는 위치에 있다면 전지적 시점이 된다. 이 시점에서는 서술자가 인물의 행동은 물론 내면 심리까지 모두 꿰뚫어 본다. 현진건의 「운수 좋은 날」이 대표적이다. 이 작품의 서술자는 주인공 김 첨지가 돈을 벌며 느끼는 기쁨과 아내에 대한 불안감이 뒤섞인 내면 심리를 모두 서술한다. 이는 인물을 깊이 있게 이해하도록 돕지만, 때로 독자의 상상력을 제한할 수도 있다.

　이처럼 소설의 시점은 각각 뚜렷한 특징과 효과를 지닌다. 작가는 자신이 전달하고자 하는 주제 의식을 가장 잘 드러낼 수 있는 시점을 의도적으로 선택한다. 그러므로 독자는 시점을 파악함으로써 작가의 의도에 더 가까이 다가갈 수 있으며, 작품을 더욱 깊이 있게 감상할 수 있다.
\par\medskip\begin{itemize}[leftmargin=18pt, itemsep=2pt, topsep=2pt, label=\textbullet]
\item 1인칭 주인공 — 김유정 「동백꽃」의 '나'
\item 1인칭 관찰자 — 주요섭 「사랑손님과 어머니」의 옥희
\item 3인칭 관찰자 — 황순원 「학」 / 전지적 — 「운수 좋은 날」
\end{itemize}
}

\kfSecHeader{kfAreaConcept}{문제}{}

\par\medskip
\kfBeginMC
\begin{kfQBlock}{01}{객관식}
\kfQStem{소설의 서술자에 대한 설명으로 적절하지 \uline{\underline{않은}} 것은?}
\begin{kfChoices}
\kfChoice 시점과 무관하게 인물의 속마음을 자유롭게 꿰뚫어 보는 존재이다.
\kfChoice 소설 속 사건과 인물에 대한 정보를 독자에게 알려 준다.
\kfChoice 작품 안에 위치할 수도 있고, 작품 밖에서 이야기할 수도 있다.
\kfChoice 시점에 따라 알 수 있는 정보의 범위가 달라진다.
\kfChoice 작가 대신 독자에게 이야기를 전달하는 존재이다.
\end{kfChoices}
\end{kfQBlock}
\begin{kfQBlock}{02}{객관식}
\kfQStem{1인칭 시점과 3인칭 시점을 구분하는 기준에 대한 설명으로 적절하지 \uline{\underline{않은}} 것은?}
\begin{kfChoices}
\kfChoice 서술자가 작품 안에 있는지, 밖에 있는지를 기준으로 한다.
\kfChoice 작품의 길이와 갈등 구조의 복잡성에 따라 시점이 결정된다.
\kfChoice 1인칭은 작품 속 ‘나’가 직접 이야기하는 시점이다.
\kfChoice 3인칭은 작품 밖 서술자가 이야기를 풀어 가는 시점이다.
\kfChoice 서술자의 위치가 1인칭과 3인칭을 가르는 핵심이다.
\end{kfChoices}
\end{kfQBlock}
\begin{kfQBlock}{03}{객관식}
\kfQStem{1인칭 주인공 시점에 대한 설명으로 적절하지 \uline{\underline{않은}} 것은?}
\begin{kfChoices}
\kfChoice 서술자가 작품 밖에서 모든 인물의 내면을 다 알고 서술한다.
\kfChoice 주인공인 ‘나’의 내면 심리를 직접 드러낼 수 있다.
\kfChoice 독자에게 친밀감과 현장감을 주는 효과가 있다.
\kfChoice ‘나’가 알지 못하는 다른 인물의 속마음은 직접 서술하기 어렵다.
\kfChoice 작품 속 ‘나’가 곧 주인공이다.
\end{kfChoices}
\end{kfQBlock}
\begin{kfQBlock}{04}{객관식}
\kfQStem{3인칭 관찰자 시점의 효과에 대한 설명으로 적절하지 \uline{\underline{않은}} 것은?}
\begin{kfChoices}
\kfChoice 서술자가 모든 인물의 속마음을 직접 설명해 보여 준다.
\kfChoice 독자가 인물의 내면을 직접 추측해 보게 만든다.
\kfChoice 사건의 극적 긴장감을 한층 높이는 효과가 있다.
\kfChoice 독자에게 다양한 해석의 폭을 열어 주는 장점이 있다.
\kfChoice 인물의 행동을 객관적으로 보여 주는 효과가 있다.
\end{kfChoices}
\end{kfQBlock}
\begin{kfQBlock}{05}{객관식}
\kfQStem{다음 <보기>의 서술상 특징에 대한 설명으로 가장 적절한 것은?}
\begin{kfEvidence}다음 날 나는 유치원에서 돌아오면서 ‘오늘은 엄마를 좀 기쁘게 해 드리고 싶은데, 무얼 갖다 드리면 기뻐하실까？’ 하는 궁리를 했어요.

그때 문득 유치원 선생님 책상 위에 놓여 있던 꽃병이 생각났어요. 그 꽃병에는 곱고 빨간 꽃이 꽂혀 있었어요. 개나리도 아니고, 진달래도 아니고 무슨 서양 꽃 같았어요. 엄마는 꽃을 무척 좋아하시니까, 그 꽃을 갖다 드리면 기뻐하실 것 같았어요.

나는 도로 유치원으로 가서 그 꽃을 두어 송이 얼른 빼들고 도망쳐나왔어요. 집에 오니 엄마는 대문 앞에서 나를 기다리고 계셨어요.

“그 꽃 어디서 났니？ 퍽 곱구나.”

나는 갑자기 말문이 막혔어요. 유치원에서 가져왔다고 말하면 혼이 날 것 같았어요. 그래서 잠깐 망설이다가,

“응, 이 꽃！ 저, 사랑 아저씨가 엄마 갖다 주라고 주셨어.”

하고 말했어요. 그런 거짓말이 어디서 그렇게 툭 튀어나왔는지 모르겠어요.

꽃을 들고 냄새를 맡던 엄마는 내 말에 몹시 놀라는 표정이었어요. 그러고는 금세 얼굴이 꽃보다 더 빨개지면서 부끄러워하셨어요.

“옥희야, 그런 걸 받아 오면 안 돼.”

엄마는 화가 난 것 같은 표정으로 말했지만, 목소리는 떨리고 있었어요. 그러더니 엄마는 꽃 얘기를 아무에게도 하지 말라고 제게 다짐을 받으셨어요.

나는 엄마가 그 꽃을 내다 버릴 거라고 생각했어요.

그런데 엄마는 그 꽃을 꽃병에 꽂아서 풍금 위에 놓아두었어요. 며칠 뒤 꽃이 시들자 줄기는 잘라 버리고 꽃송이만 찬송가 갈피에 얌전하게 넣어 두셨어요.   \\\relax
- 주요섭 「사랑손님과 어머니」\end{kfEvidence}
\begin{kfChoices}
\kfChoice 주인공인 ‘나’가 자신의 깊은 내면적 갈등을 솔직히 고백한다.
\kfChoice 서술자가 ‘엄마’의 속마음까지 모두 꿰뚫어 보고 알려 준다.
\kfChoice 작품 밖의 서술자가 사건을 철저히 객관적인 시선으로 묘사한다.
\kfChoice 서술자가 ‘엄마’의 행동을 관찰하지만 그 의미는 이해하지 못한다.
\kfChoice 어른이 된 서술자가 과거의 일을 객관적으로 평가하며 서술한다.
\end{kfChoices}
\end{kfQBlock}
\begin{kfQBlock}{06}{객관식}
\kfQStem{다음 <보기>의 서술상 특징에 대한 설명으로 가장 적절한 것은?}
\begin{kfEvidence}나흘 전 감자 쪼간(건, 사건)만 하더라도 나는 저에게 조금도 잘못한 것은 없다. 계집애가 나물을 캐러 가면 갔지 남 울타리 엮는 데 쌩이질(한창 바쁠 때에 쓸데없는 일로 남을 귀찮게 구는 짓)을 하는 것은 다 뭐냐.

그것도 발소리를 죽여 가지고 등뒤로 살며시 와서,

"얘! 너 혼자만 일하니?"

하고 긴치 않는 수작을 하는 것이다.

어제까지도 저와 나는 이야기도 잘 않고 서로 만나도 본체 만 척하고 이렇게 점잖게 지내던 터이련만 오늘로 갑작스레 대견해졌음은 웬일인가. 항차(하물며) 망아지만 한 계집애가 남 일하는 놈 보구…….

"그럼 혼자 하지 떼루 하디?"

내가 이렇게 내배앝는 소리를 하니까,

"너 일하기 좋니?"

또는,

"한여름이나 되거든 하지 벌써 울타리를 하니?"

잔소리를 두루 늘어놓다가 남이 들을까 봐 손으로 입을 틀어막고는 그 속에서 깔깔댄다. 별로 우스울 것도 없는데 날씨가 풀리더니 이 놈의 계집애가 미쳤나 하고 의심하였다.   \\\relax
- 김유정 「동백꽃」\end{kfEvidence}
\begin{kfChoices}
\kfChoice 서술자인 '나'가 주인공의 입장에서 사건을 서술한다.
\kfChoice 서술자인 '나'가 '점순이'의 행동을 객관적으로 관찰한다.
\kfChoice 작품 밖의 서술자가 두 인물의 갈등을 묘사한다.
\kfChoice 서술자가 '점순이'의 숨겨진 의도를 정확히 파악한다.
\kfChoice 서술자가 어른이 되어 과거의 일을 회상하고 있다.
\end{kfChoices}
\end{kfQBlock}
\begin{kfQBlock}{07}{객관식}
\kfQStem{이야기의 서술자를 ‘나’에서 ‘그’로 바꾸었을 때 나타나는 변화로 적절하지 \uline{\underline{않은}} 것은?}
\begin{kfChoices}
\kfChoice 서술자와 주인공의 거리가 멀어진다.
\kfChoice 주인공의 주관적인 감상이 더욱 강화된다.
\kfChoice 주인공 외 다른 인물의 내면을 서술할 수 있게 된다.
\kfChoice 작품의 시점이 1인칭에서 3인칭으로 바뀐다.
\kfChoice 독자가 느끼는 친밀감이 감소할 수 있다.
\end{kfChoices}
\end{kfQBlock}
\begin{kfQBlock}{08}{단답형}
\kfQStem{작가가 창조한 허구의 이야기를 무엇이라고 하는가?}
\kfAnswerLines{1}
\end{kfQBlock}
\begin{kfQBlock}{09}{단답형}
\kfQStem{소설에서 작가 대신 이야기를 독자에게 전달하는 이를 무엇이라고 하는가?}
\kfAnswerLines{1}
\end{kfQBlock}
\begin{kfQBlock}{10}{단답형}
\kfQStem{서술자가 사건을 바라보는 위치나 관점을 무엇이라고 하는가?}
\kfAnswerLines{1}
\end{kfQBlock}
\begin{kfQBlock}{11}{단답형}
\kfQStem{「사랑손님과 어머니」의 '옥희'처럼, '나'가 주인공이 \underline{아닌} 주변 인물로서 주인공을 관찰하는 시점은 무엇인가?}
\kfAnswerLines{1}
\end{kfQBlock}
\begin{kfQBlock}{12}{단답형}
\kfQStem{「동백꽃」처럼, 서술자인 '나'가 이야기의 주인공이 되어 자신의 속마음을 직접 고백하는 시점은 무엇인가?}
\kfAnswerLines{1}
\end{kfQBlock}
\begin{kfQBlock}{13}{단답형}
\kfQStem{「학」처럼 서술자가 인물의 내면을 드러내지 않고 행동과 대화만을 객관적으로 묘사하는 시점은 무엇인가?}
\kfAnswerLines{1}
\end{kfQBlock}
\begin{kfQBlock}{14}{단답형}
\kfQStem{「운수 좋은 날」처럼 서술자가 신처럼 인물의 행동과 내면 심리까지 모두 꿰뚫어 보는 시점은 무엇인가?}
\kfAnswerLines{1}
\end{kfQBlock}
\begin{kfQBlock}{15}{서술형}
\kfQStem{'1인칭 주인공 시점'과 '1인칭 관찰자 시점'의 공통점과 차이점을 <조건>에 맞게 서술하시오.}
\kfAnswerNote{답안}{4}
\end{kfQBlock}
\kfEndMC


\kfStartLit{「동백꽃」 — 김유정}


\kfPassageNote{}{%
　오늘도 또 우리 수탉이 막 쫓기었다. 내가 점심을 먹고 나무를 하러 갈 생각으로 나올 때이었다. 산으로 올라서려니까 등뒤에서 푸드득 푸드득 하고 닭의 날갯짓 소리가 야단이다. 깜짝 놀라서 고개를 돌려 보니 역시나 두 놈이 또 싸움이 붙었다.

　점순네 수탉(대가리가 크고 똑 오소리같이 몸이 단단해 보이는 놈)이 덩치 작은 우리 수탉을 함부로 대하는 것이다. 그것도 그냥 대하는 것이 아니라 푸드득하고 볏을 쪼고 물러섰다가 좀 사이를 두고 푸드득하고 모가지를 쪼았다. 이렇게 멋을 부려 가며 봐주지 않고 혼쭐을 내준다. 그러면 이 못생긴 것은 쪼일 적마다 주둥이로 땅을 받으며 그 비명이 킥, 킥, 할뿐이다. 물론 미처 아물지도 않은 볏을 또 쪼이며 붉은 피는 뚝뚝 떨어진다. 이걸 가만히 내려다보자니 내 머리가 터져서 피가 흐르는 것같이 두눈에서 불이 번쩍 난다. 대뜸 지게 막대기를 메고 달려들어 점순네 닭을 후려칠까 하다가 생각을 고쳐먹고 때리는 시늉만 해서 떼어놓았다.

　이번에도 점순이가 쌈을 붙여 놨을 것이다. 나를 약 올리려고 그랬음에 틀림없을 것이다. 고놈의 계집애가 요새 들어 왜 나를 그렇게 못살게 구는지 모른다.

　나흘 전 감자 사건만 하더라도 나는 저에게 조금도 잘못한 것은 없다. 계집애가 나물을 캐러 가면 갔지 남 울타리 엮는 데 훼방을 놓는 것은 다 뭐냐. 그것도 발소리를 죽여 가지고 등뒤로 살며시 와서,

　“얘! 너 혼자만 일하니?”

　하고 쓸데없는 말을 거는 것이다.

　어제까지도 저와 나는 이야기도 잘 않고 서로 만나도 못 본 체하고 이렇게 점잖게 지내던 터인데, 오늘 갑작스레 다정하게 구는 것은 웬일인가. 하물며 망아지만 한 계집애가 남 일하는 놈 보구…….

　“그럼 혼자 하지 여럿이 하니?”

　내가 이렇게 퉁명스럽게 내뱉으니까,

　“너 일하기 좋니?”

또는,

　“한여름이나 되거든 하지 벌써 울타리를 하니?”

　잔소리를 두루 늘어놓다가 남이 들을까봐 손으로 입을 틀어막고는 그 속에서 깔깔댄다. 별로 우스울 것도 없는데 날씨가 풀리더니 이 놈의 계집애가 미쳤나 하고 의심하였다. 게다가 조금 뒤에는 제 집 쪽을 곁눈질로 돌아보더니 행주치마의 속으로 꼈던 바른손을 뽑아서 나의 턱밑으로 불쑥 내미는 것이다. 언제 구웠는지 더운 김이 나는 굵은 감자 세 개가 손에 가득 쥐여 있었다.

　“너희 집엔 이거 없지?”

　하고 자기가 대단한 것을 준다는 듯이 큰소리를 하고는 제가 준 것을 남이 알면 큰일 날 테니 여기서 얼른 먹어 버리란다. 그리고 또 하는 소리가,

　“너 봄감자가 맛있단다.”

　“난 감자 안 먹는다. 너나 먹어라.”

　나는 고개도 돌리지 않고 일하던 손으로 그 감자를 도로 어깨 너머로 쑥 밀어 버렸다. 그랬더니 그래도 가는 기색이 없고, 뿐만 아니라 쌔근쌔근하고 심상치 않게 숨소리가 점점 거칠어진다. 이건 또 뭐야 싶어서 그때에야 비로소 돌아다보니 나는 참으로 놀랐다. 우리가 이 동네에 들어온 것은 거의 삼 년째 되어오지만 여태껏 살짝 검은 편인 점순이의 얼굴이 이렇게까지 홍당무처럼 새빨개진 적이 없었다. 게다가 눈에 독을 올리고 한참 나를 그렇게 쏘아보더니 나중에는 눈물까지 어리는 것이 아니냐.

[중략 줄거리] 감자를 거절당한 뒤로 점순이는 '나'의 닭들을 괴롭히기 시작한다. '나'는 닭싸움에서 이기기 위해 수탉에게 고추장을 먹여보지만 소용이 없다. 점순의 괴롭힘이 계속되자 '나'의 분노는 점점 커져만 간다.

　거의 집에 다 내려와서 나는 호드기 소리를 듣고 발이 딱 멈추었다. 산기슭에 널려 있는 굵은 바윗돌 틈에 노란 동백꽃이 수북하게 깔려 있었다. 그 틈에 끼어 앉아서 점순이가 쓸쓸하게 호드기를 불고 있는 것이다. 그보다도 더 놀란 것은 그 앞에서 또 푸드득, 푸드득, 하고 들리는 닭의 날갯짓 소리다. 틀림없이 요년이 나의 약을 올리느라고 또 닭을 집어내다가 내가 내려올 길목에다 쌈을 시켜 놓고 저는 그 앞에 앉아서 태연하게 호드기를 불고 있음에 틀림없으리라.

　나는 너무너무 화가 나서 두 눈에서 불과 함께 눈물이 퍽 쏟아졌다. 나뭇지게도 벗어 놓을 새 없이 그대로 내동댕이치고는 지게 막대기를 뻗치고 허둥지둥 달려들었다.

　가까이 와 보니 과연 나의 짐작대로 우리 수탉이 피를 흘리고 거의 죽을 지경에 이르렀다. 닭도 닭이지만, 그러함에도 불구하고 눈 하나 깜짝 없이 그대로 앉아서 호드기만 부는 그 꼴에 더욱 화가 났다.

　나는 대뜸 달려들어서 나도 모르는 사이에 큰 수탉을 단번에 때려 엎었다. 닭은 푹 엎어진 채 다리 하나 꼼짝 못 하고 그대로 죽어 버렸다. 그리고 나는 멍하니 섰다가 점순이가 매섭게 눈을 부릅뜨고 닥치는 바람에 뒤로 벌렁 나자빠졌다.

　“이놈아! 너 왜 남의 닭을 때려죽이니?”

“그럼 어때?”

하고 일어나다가,

　“뭐 이 자식아! 누구네 집 닭인데?”

　하고 가슴을 미는 바람에 다시 벌렁 자빠졌다. 그리고 나서 가만히 생각을 하니 분하기도 하고 창피하기도 하고, 또 한편 일을 저질렀으니, 이제는 땅도 떨어지고 집도 내쫓겨야 될는지 모른다.

　나는 비슬비슬 일어나며 소맷자락으로 눈을 가리고는, 얼떨결에 엉 하고 울음을 터뜨렸다. 그러나 점순이가 앞으로 다가와서,

　“그럼 너 다음부터 안 그럴 테냐?”

　하고 물을 때에야 비로소 살길을 찾은 듯싶었다. 나는 눈물을 우선 씻고 뭘 안 그런다는 것인지 뜻도 모르면서,

“그래!”

　하고 무턱대고 대답하였다.

　“다음부터 또 그래 봐라, 내 자꾸 못살게 굴 테니.”

　“그래 그래 이젠 안 그럴게!”

　“닭 죽은 건 걱정 마라, 내 안 이를 테니.”

　그리고 뭣에 떠밀렸는지 나의 어깨를 짚은 채 그대로 퍽 쓰러진다. 그 바람에 나의 몸뚱이도 겹쳐서 쓰러지며, 활짝 핀 노란 동백꽃 속으로 폭 파묻혀 버렸다.

　코끝을 톡 쏘는 듯한, 그리고 향긋한 그 냄새에 나는 땅이 꺼지는 듯이 온 정신이 고만 아찔하였다.

- 김유정 「동백꽃」
}


\kfPassageNote{문장 독해 — 발췌 문장}{%
1. 이걸 가만히 내려다보자니 내 머리가 터져서 피가 흐르는 것같이 두눈에서 불이 번쩍 난다.

2. 너희 집엔 이거 없지?

3. 나는 고개도 돌리지 않고 일하던 손으로 그 감자를 도로 어깨 너머로 쑥 밀어 버렸다. 그랬더니 그래도 가는 기색이 없고, 뿐만 아니라 쌔근쌔근하고 심상치 않게 숨소리가 점점 거칠어진다. 이건 또 뭐야 싶어서 그때에야 비로소 돌아다보니 나는 참으로 놀랐다.

4. 틀림없이 요년이 나의 약을 올리느라고 또 닭을 집어내다가 내가 내려올 길목에다 쌈을 시켜 놓고 저는 그 앞에 앉아서 태연하게 호드기를 불고 있음에 틀림없으리라.

5. 닭도 닭이지만, 그러함에도 불구하고 눈 하나 깜짝 없이 그대로 앉아서 호드기만 부는 그 꼴에 더욱 화가 났다.

6. 다음부터 또 그래 봐라
}


\kfPassageNote{해설 학습 — 본문}{%
[단원 1] 김유정의 「동백꽃」은 1930년대 강원도의 한 시골 마을을 배경으로, 사춘기 소년과 소녀의 순수한 사랑을 다룬다. 이 소설의 가장 큰 재미는 어수룩하고 눈치 없는 소년인 ‘나’의 시점에서 이야기가 진행된다는 점이다. 문학 작품에서 인물의 성격은 갈등을 만들고 이야기를 발전시키는 중요한 요소인데, 「동백꽃」의 두 주인공은 서로 다른 성격으로 인해 오해하며 주요 갈등을 만들어 낸다.

서술자인 ‘나’는 순박하고 눈치가 없으며 자존심이 강하다. 그래서 마름의 딸인 점순이가 자신을 무시한다고 생각한다. 반면 ‘점순이’는 ‘나’보다 정신적으로 성숙하며 자신의 감정을 표현하는 데 매우 적극적이다. 하지만 좋아하는 마음을 솔직하게 말하지 못하고, 일부러 닭싸움을 붙이는 등 서툴고 짓궂은 방식으로 관심을 표현한다.

독자는 점순이가 호감의 표시로 감자를 건네는 것을 금방 알 수 있지만, ‘나’는 이를 자신을 무시하는 행동으로 받아들여 차갑게 거절한다. 이처럼 좋아하는 마음을 반대로 표현하는 점순이의 행동과 그것을 전혀 이해하지 못하는 ‘나’의 순박한 모습이 계속 어긋나면서 웃음을 유발한다. 이 소설은 순진한 소년과 활달한 소녀의 갈등과 화해 과정을 통해 독자들에게 따뜻하고 해학적인 감동을 선사한다.

[단원 2] 소설에서 소재는 단순한 배경이 아니라 인물의 감정과 관계 변화를 나타내는 역할을 한다. 「동백꽃」에서 감자, 닭싸움, 동백꽃은 두 주인공의 마음과 관계 변화를 보여주는 상징적인 소재이다.

첫 번째 소재인 '감자'는 점순이가 ‘나’에게 보여주는 순수한 관심과 좋아하는 마음의 표현이다. 점순이는 부끄러워서 퉁명스럽게 건네지만, 사실은 ‘나’와 친해지고 싶은 마음이 담겨 있다. 하지만 눈치 없는 ‘나’는 이 마음을 전혀 알아채지 못하고 자존심 때문에 거절하며, 이 사건으로 둘의 갈등이 시작된다.

두 번째 소재인 ‘닭싸움’은 감자 사건으로 시작된 갈등의 표현이다. 마음이 상한 점순이는 ‘나’의 관심을 끌기 위해 매일 닭싸움을 붙인다. 하지만 ‘나’는 점순이의 속마음을 전혀 모른 채, 그저 자기를 괴롭힌다고만 생각하며 미워하는 마음을 키운다.

마지막 소재인 '동백꽃'은 길었던 갈등이 풀리고 두 사람이 화해하는 배경이 된다. 화가 난 ‘나’가 점순이네 닭을 죽이면서 갈등은 최고조에 이른다. 하지만 점순이는 ‘나’를 용서하고, 둘은 향기로운 노란 동백꽃 속으로 함께 쓰러진다. 이 장면은 모든 오해와 다툼이 끝나고 풋풋한 사랑이 시작되는 순간을 효과적으로 보여준다.

이처럼 세 가지 소재는 갈등의 시작부터 절정, 그리고 해결까지의 과정을 효과적으로 보여준다.

[단원 3] 소설은 누가 이야기를 들려주느냐에 따라 독자가 느끼는 재미가 달라진다. 같은 사건이라도 어떤 인물의 눈으로 보느냐에 따라 전혀 다른 이야기가 될 수 있기 때문이다.

이 작품은 '1인칭 주인공 시점'으로 쓰였다. 이는 주인공인 '나'가 직접 자신의 이야기를 독자에게 들려주는 방식이다. 이 시점에서는 '나'가 자신의 생각과 감정을 자세히 설명하고, 다른 인물들의 행동을 '나'의 관점에서만 바라본다. 따라서 다른 인물의 진짜 마음은 정확히 알 수 없고, '나'가 추측한 대로만 서술된다.

소설 속 '나'는 매우 순진하고 눈치가 없는 소년이다. 특히 점순이의 마음을 전혀 이해하지 못한다. 점순이가 '나'를 좋아해서 하는 행동들을 단순히 자신을 괴롭히는 것으로만 받아들인다. 반면 독자는 ‘나’의 관찰과 묘사를 통해 점순이가 '나'를 좋아한다는 사실을 짐작할 수 있다. 이처럼 독자가 아는 것과 '나'가 아는 것 사이에 차이가 생긴다.

이러한 정보의 차이로 인해 독자는 이런 '나'의 모습을 보며 답답해하면서도 동시에 웃음을 느낀다. 이때 느끼는 안타까우면서도 따뜻한 웃음을 '해학성'이라고 한다. 해학성은 단순히 웃기기만 하는 것이 아니라, 인간의 어리석음이나 순진함을 따뜻한 시선으로 바라보며 생기는 건강한 웃음을 말한다. 이런 해학성은 우리 문학의 중요한 특징 중 하나이다.
}


\kfActivityBg \par\needspace{15\baselineskip}
\kfSecHeader{kfAreaLiterature}{활동}{문장 독해}
\par\kfMaybeNeedspace{32\baselineskip}\noindent{\kfTipMark\,\,}{\small\color{kfInk} 다음은 작품에서 발췌한 문장입니다. 각 문장을 읽고 물음에 답하시오.}\par\nopagebreak[4]\smallskip\nopagebreak[4]
\kfBeginMC
\begin{kfSentenceUnit}{1}{}
\kfSentQ{1.}{다음 문장에서 서술어 '내려다보자니', '터져서', '번쩍 난다'의 주어는 각각 무엇인가?}
\begin{kfChoices}
\kfChoice 내가, 내 머리가, 불이
\kfChoice 내가, 나의 지게 막대기가, 불이
\kfChoice 점순이가, 점순네 수탉이, 눈물이
\kfChoice 점순네 수탉이, 나의 몸뚱이가, 피가
\end{kfChoices}
\end{kfSentenceUnit}

\begin{kfSentenceUnit}{2}{}
\kfSentQ{2.}{다음 대화에서 '이거'가 가리키는 것은 무엇인가?}
\begin{kfChoices}
\kfChoice 점순이가 직접 만든 떡
\kfChoice 점순네 집에서 기른 닭
\kfChoice 점순이가 아끼는 호드기
\kfChoice 더운 김이 나는 굵은 감자 세 개
\end{kfChoices}
\end{kfSentenceUnit}

\begin{kfSentenceUnit}{3}{}
\kfSentQ{3.}{다음 문장에서 '이건'이 가리키는 상황은 무엇인가?}
\begin{kfChoices}
\kfChoice 점순이가 감자를 주며 수줍어하는 상황
\kfChoice 점순이가 몸이 아파서 괴로워하는 상황
\kfChoice 점순이가 거절당해 화가 나고 흥분한 상황
\kfChoice 점순이가 닭싸움을 시키며 즐거워하는 상황
\end{kfChoices}
\end{kfSentenceUnit}

\begin{kfSentenceUnit}{4}{}
\kfSentQ{4.}{다음 문장 속에서 닭싸움을 시키고 호드기를 부는 주체는 각각 누구인가?}
\begin{kfChoices}
\kfChoice 나, 나
\kfChoice 점순이, 점순이
\kfChoice 나, 점순네 수탉
\kfChoice 우리 수탉, 점순이
\end{kfChoices}
\end{kfSentenceUnit}

\begin{kfSentenceUnit}{5}{}
\kfSentQ{5.}{다음 문장에서 '나'를 더욱 화나게 한 '그 꼴'은 구체적으로 어떤 모습인가?}
\begin{kfChoices}
\kfChoice 감자를 거절당하고 우는 모습
\kfChoice 우리 수탉이 피를 흘리는 모습
\kfChoice 닭이 죽어가는데도 점순이가 태연하게 앉아 호드기만 불고 있는 모습
\kfChoice 점순이가 나를 약 올리기 위해 일부러 닭싸움을 붙여 놓고 즐거워하는 모습
\end{kfChoices}
\end{kfSentenceUnit}

\begin{kfSentenceUnit}{6}{}
\kfSentQ{6.}{다음 대화에서 점순이가 "다음부터 또 그래 봐라"라고 말할 때, '그래'가 가리키는 '나'의 행동은 무엇인가?

"그럼 너 다음부터 안 그럴 테냐?" / "그래!" / "다음부터 또 그래 봐라, 내 자꾸 못살게 굴 테니."}
\begin{kfChoices}
\kfChoice 욕을 대드는 행동
\kfChoice 점순이의 호의를 거절하는 행동
\kfChoice 점순네 닭을 때려죽인 행동
\kfChoice 호드기 소리를 무시하는 행동
\end{kfChoices}
\end{kfSentenceUnit}

\kfEndMC

\kfActivityBg \kfSecHeader{kfAreaLiterature}{활동}{해설 문제}
\par\kfMaybeNeedspace{18\baselineskip}\noindent{\kfTipMark\,\,}{\small\color{kfInk} 해설 본문을 읽고, 단원별 단답형 물음에 답하시오.}\par\nopagebreak[4]\smallskip\nopagebreak[4]
\kfBeginMC
\begin{enumerate}[leftmargin=22pt, itemsep=6pt, topsep=4pt, label=\textbf{\arabic*.}]
\item 이 소설의 시대적, 공간적 배경은 어디인가?
\item 이 소설의 가장 큰 재미를 만들어내는 서술 방식의 특징은 무엇인가?
\item 서술자인 ‘나’의 성격을 세 가지로 요약하면 무엇인가?
\item 점순이는 좋아하는 마음을 어떤 방식으로 표현하는가?
\item ‘나’는 점순이가 감자를 건넨 진짜 이유를 무엇이라고 오해했는가?
\item 이 소설에서 웃음을 만들어내는 핵심적인 원인은 무엇인가?
\item 두 주인공 사이에 갈등이 생기는 가장 큰 이유는 무엇인가?
\item 점순이가 '나'에게 자신의 마음을 표현하기 위해 사용한 첫 번째 소재는 무엇인가?
\item '감자'는 점순이의 어떤 마음을 상징하는가?
\item '나'가 감자를 거절한 사건은 둘 사이에 무엇이 시작되는 계기가 되는가?
\item 점순이가 '나'의 관심을 끌기 위해 계속해서 이용하는 사건은 무엇인가?
\item '닭싸움'은 두 사람의 갈등을 어떻게 만들어주는 역할을 하는가?
\item '나'와 점순이가 극적으로 화해하게 되는 공간적 배경이 되는 소재는 무엇인가?
\item 소설 마지막에 두 사람이 동백꽃 속으로 쓰러지는 장면은 무엇을 의미하는가?
\item 이 소설이 채택하고 있는 시점은 무엇인가?
\item 1인칭 주인공 시점에서는 누구의 속마음만 직접 알 수 있는가?
\item 1인칭 주인공 시점을 사용함으로써 독자는 누구의 입장에서 사건을 생생하게 체험하게 되는가?
\item 독자들은 점순이의 속마음을 어떻게 알 수 있는가?
\item '해학'은 무엇을 유발하는 표현 방식인가?
\item 독자가 '나'를 보며 안타까움과 웃음을 동시에 느끼는 이유는 무엇 때문인가?
\end{enumerate}
\kfEndMC

\kfSecHeader{kfAreaLiterature}{문제}{}

\par\medskip
\kfBeginMC
\begin{kfQBlock}{01}{객관식}
\kfQStem{이 글에 대한 설명으로 가장 적절한 것은?}
\begin{kfChoices}
\kfChoice 역사적 사건을 배경으로 인물의 고난을 그리고 있다.
\kfChoice 배경에 대한 섬세한 묘사를 통해 신비로운 분위기를 조성하고 있다.
\kfChoice 서술자의 눈에 비친 사건을 중심으로 이야기가 전개되고 있다.
\kfChoice 과거와 현재를 오가며 인물의 복잡한 내면을 보여주고 있다.
\kfChoice 비현실적인 사건을 통해 현실의 문제점을 비판하고 있다.
\end{kfChoices}
\end{kfQBlock}
\begin{kfQBlock}{02}{객관식}
\kfQStem{점순이가 '나'의 닭을 괴롭히는 진짜 이유로 가장 알맞은 것은?}
\begin{kfChoices}
\kfChoice ‘나’의 닭이 자기 집 곡식을 먹었기 때문에
\kfChoice 자신의 마음을 몰라주는 '나'에게 관심을 받고 싶어서
\kfChoice 어떤 특별한 악의가 있는 것이 아니라, 단순히 심심해서
\kfChoice 닭을 이용해 돈을 벌려는 계산적인 마음이 있었기 때문에
\kfChoice ‘나’의 집안에 대해 좋지 않은 감정을 가지고 있었기 때문에
\end{kfChoices}
\end{kfQBlock}
\begin{kfQBlock}{03}{객관식}
\kfQStem{이 소설의 서술상 특징이 주는 효과에 대한 설명으로 적절하지 \uline{\underline{않은}} 것은?}
\begin{kfChoices}
\kfChoice 어수룩한 서술자 ‘나’의 시점을 통해 작품 전체의 해학적 분위기를 만든다.
\kfChoice 전지적 서술자가 점순이의 내면 심리를 직접 설명해 인물 이해를 돕는다.
\kfChoice 서술자의 인식과 실제 상황의 차이로 독자에게 재미를 준다.
\kfChoice 1인칭 주인공 시점이 인물의 순박함을 효과적으로 드러낸다.
\kfChoice ‘나’가 점순이의 마음을 오해하면서 발생하는 상황적 웃음을 유발한다.
\end{kfChoices}
\end{kfQBlock}
\begin{kfQBlock}{04}{객관식}
\kfQStem{'감자'와 '닭'의 기능에 대한 설명으로 가장 적절한 것은?}
\begin{kfChoices}
\kfChoice '감자'는 갈등의 시작, '닭'은 갈등의 심화를 의미한다.
\kfChoice '감자'는 물질적 가치, '닭'은 생명의 소중함을 상징한다.
\kfChoice '감자'는 도시의 문물, '닭'은 농촌의 일상을 상징한다.
\kfChoice '감자'는 '나'의 애정, '닭'은 점순이의 애정을 의미한다.
\kfChoice '감자'는 화해의 매개체, '닭'은 갈등의 매개체를 의미한다.
\end{kfChoices}
\end{kfQBlock}
\begin{kfQBlock}{05}{객관식}
\kfQStem{㉠\textasciitilde{}㉤ 중, '나'의 어수룩한 모습을 알 수 있는 부분이 \uline{\underline{아닌}} 것은?}
\begin{kfChoices}
\kfChoice ㉤ 이제는 땅도 떨어지고 집도 내쫓겨야 될는지 모른다.
\kfChoice ㉡ 계집애가 나물을 캐러 가면 갔지 남 울타리 엮는 데 훼방을 놓는 것은 다 뭐냐.
\kfChoice ㉢ 오늘 갑작스레 다정하게 구는 것은 웬일인가.
\kfChoice ㉣ 별로 우스울 것도 없는데 날씨가 풀리더니 이 놈의 계집애가 미쳤나 하고 의심하였다.
\kfChoice ㉠ 고놈의 계집애가 요새 들어 왜 나를 그렇게 못살게 구는지 모른다.
\end{kfChoices}
\end{kfQBlock}
\begin{kfQBlock}{06}{객관식}
\kfQStem{'나'와 점순이가 '동백꽃' 속으로 쓰러지는 장면의 기능으로 가장 적절한 것은?}
\begin{kfChoices}
\kfChoice 두 사람의 비극적인 사랑의 결말을 암시한다.
\kfChoice 앞으로 두 사람에게 닥칠 더 큰 시련을 예고한다.
\kfChoice 자연의 아름다움과 인간 세상의 갈등을 대비시킨다.
\kfChoice 인물들의 갈등이 해소되고 사랑이 시작되는 계기를 마련한다.
\kfChoice 현실의 어려움을 잊고 이상 세계로 도피하고 싶은 마음을 보여준다.
\end{kfChoices}
\end{kfQBlock}
\begin{kfQBlock}{07}{객관식}
\kfQStem{다음 <보기>를 참고할 때, 점순이의 말하기 방식으로 가장 적절한 것은?}
\begin{kfEvidence}사람들은 때로 자신의 속마음을 그대로 드러내지 않고 반대로 표현함으로써 상대방의 반응을 살피거나 관심을 유도한다. 좋아하는 마음을 일부러 퉁명스럽게 말하거나, 화난 마음을 오히려 웃으며 말하는 것이 그 예이다.\end{kfEvidence}
\begin{kfChoices}
\kfChoice 자신의 감정을 솔직하고 직접적으로 표현하고 있다.
\kfChoice 논리적인 근거를 들어 상대방을 설득하고 있다.
\kfChoice 속마음과 반대되는 행동으로 상대의 관심을 끌고 있다.
\kfChoice 어려운 단어를 사용하여 자신의 지식을 뽐내고 있다.
\kfChoice 상대방의 기분을 좋게 하는 칭찬을 주로 사용하고 있다.
\end{kfChoices}
\end{kfQBlock}
\begin{kfQBlock}{08}{서술형}
\kfQStem{'나'가 점순이의 호의를 계속해서 거절하는 이유를 <조건>에 맞게 서술하시오.}
\kfAnswerNote{답안}{4}
\end{kfQBlock}
\begin{kfQBlock}{09}{서술형}
\kfQStem{이 소설이 '1인칭 주인공 시점'을 사용하여 얻는 효과를 <조건>에 맞게 서술하시오.}
\kfAnswerNote{답안}{4}
\end{kfQBlock}
\kfEndMC


\kfStartNonlit{[사회] 미성년자 보호를 위한 계약 취소 제도}


\kfPassageNote{[사회] 미성년자 보호를 위한 계약 취소 제도}{%
　17세 고등학생이 부모 동의 없이 60만 원짜리 다이어트 식품을 할부로 사서 절반 정도 먹었지만 효과가 없어 계약을 취소하려고 한다. 판매업자는 "사용한 만큼 돈을 내야 한다"며 이미 받은 20만 원에 10만 원을 더 달라고 요구했다. 과연 학생이 추가로 돈을 내야 할까?

　최근 10대들의 구매력이 높아지면서 이런 분쟁이 늘고 있다. 민법에 따르면 19세 미만 미성년자는 원칙적으로 부모 같은 법정대리인의 동의 없이 계약을 할 수 없다. 동의 없이 계약했다면 나중에라도 부모가 계약을 취소할 수 있다. 이는 미성년자가 어른보다 경험과 판단력이 부족해서 자신에게 손해가 되는 계약을 하지 않도록 보호하기 위한 제도다.

　보통 계약을 취소하면 계약 전 상태로 되돌아가는 것이 원칙이다. 판매업자는 받은 돈을 돌려주고, 구매자는 상품을 반환한다. 상품을 이미 사용했다면 구매자가 사용한 만큼의 돈을 판매업자에게 돌려주는 것이 일반적이다.

　하지만 미성년자의 경우는 다르다. 민법 제141조는 미성년자가 계약을 취소할 때 '받은 이익이 현존하는 범위에서만 돈을 돌려주면 된다'고 정하고 있다. 즉, 생활필수품을 샀다면 실제로 도움이 되었으므로 사용한 만큼 돈을 내야 한다. 하지만 다이어트 식품처럼 꼭 필요하지 않은 물건을 샀다면 사용한 만큼의 돈을 낼 필요가 없다. 대신 남은 상품만 돌려주고 이미 낸 돈은 모두 돌려받을 수 있다.

　따라서 위 사례에서 학생이 계약을 취소하고 남은 다이어트 식품을 돌려줬다면, 판매업자는 받은 20만 원을 모두 돌려줘야 한다. 학생은 추가로 10만 원을 낼 필요가 없다.

　그런데 미성년자가 마음대로 계약을 취소하면 판매업자가 손해를 볼 수 있어서 이를 막는 제도도 있다. '취소권 배제'는 미성년자가 거짓말로 자신이 성년인 척하거나 부모 동의가 있는 척 속인 경우 취소권을 없애는 것이다. '확답 촉구권'은 판매업자가 1개월 이상 기간을 정해 미성년자 부모에게 계약 취소 여부를 물어볼 수 있는 권리다. 정해진 기간에 답하지 않으면 계약을 인정한 것으로 본다. '철회권'은 판매업자가 부모 동의 전까지 먼저 계약을 취소할 수 있는 권리다. 단, 처음부터 미성년자인 걸 알고 계약했다면 철회권을 쓸 수 없다.

　한편 미성년자도 부모 동의 없이 할 수 있는 계약이 있다. 버스나 지하철 타기, 김밥이나 과자 사기 등 일상적인 거래는 자유롭게 할 수 있다.
}


\kfPassageNote{문장 독해 — 본문 문장}{%
1. 17세 고등학생이 부모 동의 없이 60만 원짜리 다이어트 식품을 할부로 사서 절반 정도 먹었지만 효과가 없어 계약을 취소하려고 한다.

2. 판매업자는 "사용한 만큼 돈을 내야 한다"며 이미 받은 20만 원에 10만 원을 더 달라고 요구했다.

3. 과연 학생이 추가로 돈을 내야 할까?

4. 최근 10대들의 구매력이 높아지면서 이런 분쟁이 늘고 있다.

5. 민법에 따르면 19세 미만 미성년자는 원칙적으로 부모 같은 법정대리인의 동의 없이 계약을 할 수 없다.

6. 동의 없이 계약했다면 나중에라도 부모가 계약을 취소할 수 있다.

7. 이는 미성년자가 어른보다 경험과 판단력이 부족해서 자신에게 손해가 되는 계약을 하지 않도록 보호하기 위한 제도다.

8. 보통 계약을 취소하면 계약 전 상태로 되돌아가는 것이 원칙이다.

9. 판매업자는 받은 돈을 돌려주고, 구매자는 상품을 반환한다.

10. 상품을 이미 사용했다면 구매자가 사용한 만큼의 돈을 판매업자에게 돌려주는 것이 일반적이다.

11. 하지만 미성년자의 경우는 다르다.

12. 민법 제141조는 미성년자가 계약을 취소할 때 '받은 이익이 현존하는 범위에서만 돈을 돌려주면 된다'고 정하고 있다.

13. 즉, 생활필수품을 샀다면 실제로 도움이 되었으므로 사용한 만큼 돈을 내야 한다.

14. 하지만 다이어트 식품처럼 꼭 필요하지 않은 물건을 샀다면 사용한 만큼의 돈을 낼 필요가 없다.

15. 대신 남은 상품만 돌려주고 이미 낸 돈은 모두 돌려받을 수 있다.

16. 따라서 위 사례에서 학생이 계약을 취소하고 남은 다이어트 식품을 돌려줬다면, 판매업자는 받은 20만 원을 모두 돌려줘야 한다.

17. 학생은 추가로 10만 원을 낼 필요가 없다.

18. 그런데 미성년자가 마음대로 계약을 취소하면 판매업자가 손해를 볼 수 있어서 이를 막는 제도도 있다.

19. '취소권 배제'는 미성년자가 거짓말로 자신이 성년인 척하거나 부모 동의가 있는 척 속인 경우 취소권을 없애는 것이다.

20. '확답 촉구권'은 판매업자가 1개월 이상 기간을 정해 미성년자 부모에게 계약 취소 여부를 물어볼 수 있는 권리다.

21. 정해진 기간에 답하지 않으면 계약을 인정한 것으로 본다.

22. '철회권'은 판매업자가 부모 동의 전까지 먼저 계약을 취소할 수 있는 권리다.

23. 단, 처음부터 미성년자인 걸 알고 계약했다면 철회권을 쓸 수 없다.

24. 한편 미성년자도 부모 동의 없이 할 수 있는 계약이 있다.

25. 버스나 지하철 타기, 김밥이나 과자 사기 등 일상적인 거래는 자유롭게 할 수 있다.
}


\kfActivityBg \par\kfMaybeNeedspace{32\baselineskip}\noindent{\kfTipMark\,\,}{\small\color{kfInk} 다음 표의 '의미'를 읽고, 그에 해당하는 '어휘'를 지문에서 찾아 적으시오.}\par\nopagebreak[4]\smallskip\nopagebreak[4]
\begin{kfVocab}
\kfVocabItem{1}{다툼이나 싸움}
\kfVocabItem{2}{개인 사이의 재산이나 가족 관계 등을 다루는 법}
\kfVocabItem{3}{만 19세가 되지 않은 사람}
\kfVocabItem{4}{법률에 따라 미성년자 등을 대신하여 법률 행위를 하는 사람(부모 등)}
\kfVocabItem{5}{법적인 효력이 있는 약속}
\kfVocabItem{6}{한번 결정한 것을 무효로 하여 없애는 것}
\kfVocabItem{7}{원래 주인에게 다시 돌려주는 것}
\kfVocabItem{8}{현재 실제로 존재하는 것}
\kfVocabItem{9}{먹고 입고 자는 등 사람이 살아가는 데 꼭 필요한 물건}
\kfVocabItem{10}{어떤 권리나 자격을 갖지 못하게 막는 것}
\kfVocabItem{11}{상대방에게 확실한 대답을 재촉하여 요구할 수 있는 권리}
\kfVocabItem{12}{효력이 생기기 전에 계약을 없던 일로 할 수 있는 권리}
\end{kfVocab}

\kfActivityBg \par\needspace{15\baselineskip}
\kfSecHeader{kfAreaNonlit}{활동}{문장 독해}
\par\kfMaybeNeedspace{32\baselineskip}\noindent{\kfTipMark\,\,}{\small\color{kfInk} 다음은 본문의 번호 매긴 문장입니다. 각 문장을 읽고 물음에 답하시오.}\par\nopagebreak[4]\smallskip\nopagebreak[4]
\kfBeginMC
\begin{kfSentenceUnit}{1}{}
\kfSentQ{1.}{사례 속 학생이 계약을 할 때 없었던 것은 무엇인가?}
\begin{kfChoices}
\kfChoice 판매자의 적극적 권유
\kfChoice 구매 물품에 대한 정보
\kfChoice 다이어트 식품의 효과
\kfChoice 법정대리인의 동의
\end{kfChoices}
\end{kfSentenceUnit}

\begin{kfSentenceUnit}{2}{}
\kfSentQ{2.}{학생이 다이어트 식품에 효과가 없자 무엇을 하려 하는가?}
\begin{kfChoices}
\kfChoice 상품 교환
\kfChoice 계약 취소
\kfChoice 가격 할인
\kfChoice 할부 연장
\end{kfChoices}
\end{kfSentenceUnit}

\begin{kfSentenceUnit}{3}{}
\kfSentQ{3.}{판매업자가 학생에게 돈을 더 달라고 요구하는 근거는 무엇인가?}
\begin{kfChoices}
\kfChoice 학생이 부모의 동의를 받았다고 속였기 때문
\kfChoice 학생이 이미 사용한 만큼 돈을 내야 한다는 것
\kfChoice 계약서에 위약금 조항이 명시되어 있기 때문
\kfChoice 다이어트 식품의 가격이 구매 후 올랐기 때문
\end{kfChoices}
\end{kfSentenceUnit}

\begin{kfSentenceUnit}{4}{}
\kfSentQ{4.}{판매업자는 학생에게 총 얼마를 내라고 요구했는가?}
\begin{kfChoices}
\kfChoice 20만 원
\kfChoice 30만 원
\kfChoice 40만 원
\kfChoice 60만 원
\end{kfChoices}
\end{kfSentenceUnit}

\begin{kfSentenceUnit}{5}{}
\kfSentQ{5.}{이 문장은 글 전체에서 어떤 역할을 하는가?}
\begin{kfChoices}
\kfChoice 글의 결론을 미리 제시하는 역할
\kfChoice 앞으로 다룰 핵심 질문을 던지는 역할
\kfChoice 상반된 두 주장을 중재하는 역할
\kfChoice 구체적인 해결 방안을 나열하는 역할
\end{kfChoices}
\end{kfSentenceUnit}

\begin{kfSentenceUnit}{6}{}
\kfSentQ{6.}{'이런 분쟁'이 늘어나는 이유는 무엇인가?}
\begin{kfChoices}
\kfChoice 10대들의 구매력이 높아졌기 때문
\kfChoice 판매업자의 상술이 교묘해졌기 때문
\kfChoice 미성년자 보호법이 강화되었기 때문
\kfChoice 부모들의 경제 교육이 부족하기 때문
\end{kfChoices}
\end{kfSentenceUnit}

\begin{kfSentenceUnit}{7}{}
\kfSentQ{7.}{민법상 미성년자는 몇 세 미만을 가리키는가?}
\begin{kfChoices}
\kfChoice 15세 미만
\kfChoice 17세 미만
\kfChoice 19세 미만
\kfChoice 20세 미만
\end{kfChoices}
\end{kfSentenceUnit}

\begin{kfSentenceUnit}{8}{}
\kfSentQ{8.}{미성년자가 계약을 하려면 원칙적으로 누구의 동의가 필요한가?}
\begin{kfChoices}
\kfChoice 학교 선생님
\kfChoice 법정대리인
\kfChoice 형제자매
\kfChoice 친한 친구
\end{kfChoices}
\end{kfSentenceUnit}

\begin{kfSentenceUnit}{9}{}
\kfSentQ{9.}{부모는 어떤 경우에 미성년 자녀의 계약을 취소할 수 있는가?}
\begin{kfChoices}
\kfChoice 자녀가 물건을 마음에 들어하지 않을 때
\kfChoice 자녀가 물건을 이미 사용해 버렸을 때
\kfChoice 자녀가 부모 동의 없이 계약했을 때
\kfChoice 판매자가 계약 내용을 설명하지 않았을 때
\end{kfChoices}
\end{kfSentenceUnit}

\begin{kfSentenceUnit}{10}{}
\kfSentQ{10.}{미성년자 계약 취소 제도가 있는 이유는 무엇인가?}
\begin{kfChoices}
\kfChoice 손해를 볼 수 있는 판매업자의 이익을 보장하기 위해
\kfChoice 경제적으로 어려운 부모의 금전적 부담을 줄여주기 위해
\kfChoice 사회 경험이 필요한 미성년자의 과소비를 조장하기 위해
\kfChoice 경험과 판단력이 부족한 미성년자를 보호하기 위해
\end{kfChoices}
\end{kfSentenceUnit}

\begin{kfSentenceUnit}{11}{}
\kfSentQ{11.}{일반적인 계약 취소의 원칙은 무엇인가?}
\begin{kfChoices}
\kfChoice 위약금을 지불하고 계약을 끝내는 것
\kfChoice 받은 물건을 다른 것으로 교환하는 것
\kfChoice 계약 전 상태로 완전히 되돌아가는 것
\kfChoice 계약 기간을 연장하여 다시 맺는 것
\end{kfChoices}
\end{kfSentenceUnit}

\begin{kfSentenceUnit}{12}{}
\kfSentQ{12.}{계약 전 상태로 돌아가기 위해 판매업자는 무엇을 해야 하는가?}
\begin{kfChoices}
\kfChoice 받은 돈을 돌려주어야 한다.
\kfChoice 상품을 수리해주어야 한다.
\kfChoice 새 상품으로 교환해주어야 한다.
\kfChoice 위약금을 청구해야 한다.
\end{kfChoices}
\end{kfSentenceUnit}

\begin{kfSentenceUnit}{13}{}
\kfSentQ{13.}{계약 전 상태로 돌아가기 위해 구매자는 무엇을 해야 하는가?}
\begin{kfChoices}
\kfChoice 영수증을 폐기해야 한다.
\kfChoice 상품을 반환해야 한다.
\kfChoice 대금을 추가 지급해야 한다.
\kfChoice 사유서를 작성해야 한다.
\end{kfChoices}
\end{kfSentenceUnit}

\begin{kfSentenceUnit}{14}{}
\kfSentQ{14.}{상품을 이미 사용한 경우, 구매자는 보통 어떻게 해야 하는가?}
\begin{kfChoices}
\kfChoice 사용한 물건을 즉시 폐기 처분해야 한다.
\kfChoice 새 상품 가격 전액을 현금으로 배상해야 한다.
\kfChoice 사용 여부와 상관없이 무조건 반품하면 된다.
\kfChoice 사용한 만큼의 돈을 판매업자에게 돌려주어야 한다.
\end{kfChoices}
\end{kfSentenceUnit}

\begin{kfSentenceUnit}{15}{}
\kfSentQ{15.}{'하지만'이라는 표현을 볼 때, 앞으로 어떤 내용이 나올 것을 예상할 수 있는가?}
\begin{kfChoices}
\kfChoice 미성년자도 성인과 똑같이 책임을 져야 한다는 내용
\kfChoice 일반적인 계약 취소 원칙이 강화된다는 내용
\kfChoice 일반적인 경우와 다른 미성년자 계약 취소의 특징
\kfChoice 판매업자를 보호하기 위한 특별 조항에 대한 내용
\end{kfChoices}
\end{kfSentenceUnit}

\begin{kfSentenceUnit}{16}{}
\kfSentQ{16.}{미성년자가 계약을 취소할 때 돈을 돌려줘야 하는 범위는 어디까지인가?}
\begin{kfChoices}
\kfChoice 받은 이익이 현존하는 범위
\kfChoice 계약 당시 지불한 금액 전체
\kfChoice 상품의 정가에 해당하는 금액
\kfChoice 판매자가 입은 손해액 전체
\end{kfChoices}
\end{kfSentenceUnit}

\begin{kfSentenceUnit}{17}{}
\kfSentQ{17.}{미성년자가 사용한 만큼 돈을 내야 하는 경우는 어떤 물건을 샀을 때인가?}
\begin{kfChoices}
\kfChoice 기호식품
\kfChoice 오락용품
\kfChoice 생활필수품
\kfChoice 사치품
\end{kfChoices}
\end{kfSentenceUnit}

\begin{kfSentenceUnit}{18}{}
\kfSentQ{18.}{사례의 다이어트 식품은 이 문장에서 말하는 어떤 물건에 해당하는가?}
\begin{kfChoices}
\kfChoice 건강 증진을 위해 필수적인 물건
\kfChoice 생활하는 데 없어서는 안 될 물건
\kfChoice 꼭 필요하지 않은 물건
\kfChoice 법적으로 구매가 금지된 물건
\end{kfChoices}
\end{kfSentenceUnit}

\begin{kfSentenceUnit}{19}{}
\kfSentQ{19.}{꼭 필요하지 않은 물건을 산 미성년자가 계약을 취소할 때 해야 할 일은 무엇인가?}
\begin{kfChoices}
\kfChoice 남은 상품만 돌려주는 것
\kfChoice 사용한 만큼 비용을 지불하는 것
\kfChoice 새 상품으로 보상해 주는 것
\kfChoice 부모가 대신 사과하는 것
\end{kfChoices}
\end{kfSentenceUnit}

\begin{kfSentenceUnit}{20}{}
\kfSentQ{20.}{이때 미성년자가 이미 낸 돈은 어떻게 되는가?}
\begin{kfChoices}
\kfChoice 절반만 돌려받을 수 있다.
\kfChoice 모두 돌려받을 수 있다.
\kfChoice 포인트로 적립받는다.
\kfChoice 기부금으로 처리된다.
\end{kfChoices}
\end{kfSentenceUnit}

\begin{kfSentenceUnit}{21}{}
\kfSentQ{21.}{사례의 판매업자가 학생에게 해야 할 일은 무엇인가?}
\begin{kfChoices}
\kfChoice 받은 20만 원을 모두 돌려주는 것
\kfChoice 10만 원을 제외하고 돌려주는 것
\kfChoice 남은 상품을 폐기 처분하는 것
\kfChoice 학생의 부모에게 항의하는 것
\end{kfChoices}
\end{kfSentenceUnit}

\begin{kfSentenceUnit}{22}{}
\kfSentQ{22.}{사례의 학생이 판매업자에게 추가로 내야 할 금액은 얼마인가?}
\begin{kfChoices}
\kfChoice 10만 원
\kfChoice 20만 원
\kfChoice 30만 원
\kfChoice 0원
\end{kfChoices}
\end{kfSentenceUnit}

\begin{kfSentenceUnit}{23}{}
\kfSentQ{23.}{이 문장에서 소개하는 제도는 누구를 보호하기 위한 것인가?}
\begin{kfChoices}
\kfChoice 미성년자
\kfChoice 법정대리인
\kfChoice 판매업자
\kfChoice 소비자보호원
\end{kfChoices}
\end{kfSentenceUnit}

\begin{kfSentenceUnit}{24}{}
\kfSentQ{24.}{어떤 경우에 미성년자의 취소권이 없어지는가?}
\begin{kfChoices}
\kfChoice 판매업자가 미성년자에게 물건을 터무니없이 비싸게 판 경우
\kfChoice 자신이 성년이거나 부모 동의가 있는 것처럼 속인 경우
\kfChoice 부모가 자녀의 계약 사실을 나중에야 비로소 알게 된 경우
\kfChoice 물건을 구매한 지 1년이 지나서 계약의 효력이 사라진 경우
\end{kfChoices}
\end{kfSentenceUnit}

\begin{kfSentenceUnit}{25}{}
\kfSentQ{25.}{'확답 촉구권'은 누가 행사하는 권리인가?}
\begin{kfChoices}
\kfChoice 미성년자
\kfChoice 부모님
\kfChoice 판매업자
\kfChoice 경찰관
\end{kfChoices}
\end{kfSentenceUnit}

\begin{kfSentenceUnit}{26}{}
\kfSentQ{26.}{'확답 촉구권'은 누구를 대상으로 행사하는 권리인가?}
\begin{kfChoices}
\kfChoice 미성년자 본인
\kfChoice 미성년자의 부모
\kfChoice 동료 판매업자
\kfChoice 소비자 분쟁 위원회
\end{kfChoices}
\end{kfSentenceUnit}

\begin{kfSentenceUnit}{27}{}
\kfSentQ{27.}{'확답 촉구권'을 통해 판매업자가 부모에게 묻는 내용은 무엇인가?}
\begin{kfChoices}
\kfChoice 상품 구매 이유
\kfChoice 대금 지불 능력
\kfChoice 자녀 교육 방식
\kfChoice 계약 취소 여부
\end{kfChoices}
\end{kfSentenceUnit}

\begin{kfSentenceUnit}{28}{}
\kfSentQ{28.}{판매업자의 확답 촉구에 부모가 정해진 기간 동안 답을 하지 않으면 계약은 어떻게 되는가?}
\begin{kfChoices}
\kfChoice 계약이 자동으로 취소된다.
\kfChoice 계약을 인정한 것으로 본다.
\kfChoice 계약 기간이 연장된다.
\kfChoice 법적 소송으로 이어진다.
\end{kfChoices}
\end{kfSentenceUnit}

\begin{kfSentenceUnit}{29}{}
\kfSentQ{29.}{'철회권'은 누가 행사하는 권리인가?}
\begin{kfChoices}
\kfChoice 미성년자
\kfChoice 법정대리인
\kfChoice 판매업자
\kfChoice 제3자
\end{kfChoices}
\end{kfSentenceUnit}

\begin{kfSentenceUnit}{30}{}
\kfSentQ{30.}{'철회권'을 통해 판매업자는 무엇을 할 수 있는가?}
\begin{kfChoices}
\kfChoice 상품 가격을 올릴 수 있다.
\kfChoice 먼저 계약을 취소할 수 있다.
\kfChoice 강제로 대금을 받을 수 있다.
\kfChoice 부모의 동의를 요구할 수 있다.
\end{kfChoices}
\end{kfSentenceUnit}

\begin{kfSentenceUnit}{31}{}
\kfSentQ{31.}{판매업자는 '철회권'을 언제까지 행사할 수 있는가?}
\begin{kfChoices}
\kfChoice 계약 후 1주일 이내
\kfChoice 물건을 배송하기 전까지
\kfChoice 부모의 동의가 있기 전까지
\kfChoice 미성년자가 성인이 될 때까지
\end{kfChoices}
\end{kfSentenceUnit}

\begin{kfSentenceUnit}{32}{}
\kfSentQ{32.}{판매업자가 '철회권'을 쓸 수 \uline{없는} 경우는 언제인가?}
\begin{kfChoices}
\kfChoice 계약 당사자인 미성년자가 자신의 신분에 대해 거짓말을 했을 때
\kfChoice 미성년자의 법정대리인인 부모가 계약에 대한 동의를 거부했을 때
\kfChoice 계약 취소로 인해 판매자가 감당할 수 없을 만큼 손해를 보았을 때
\kfChoice 계약할 때부터 상대방이 미성년자인 것을 알고 있었던 경우
\end{kfChoices}
\end{kfSentenceUnit}

\begin{kfSentenceUnit}{33}{}
\kfSentQ{33.}{이 문장은 미성년자 계약의 어떤 측면을 설명하고 있는가?}
\begin{kfChoices}
\kfChoice 엄격하게 금지되는 측면
\kfChoice 예외적으로 허용되는 측면
\kfChoice 부작용이 발생하는 측면
\kfChoice 법적으로 처벌받는 측면
\end{kfChoices}
\end{kfSentenceUnit}

\begin{kfSentenceUnit}{34}{}
\kfSentQ{34.}{미성년자가 부모 동의 없이 자유롭게 할 수 있는 거래의 종류는 무엇인가?}
\begin{kfChoices}
\kfChoice 고가의 전자제품 구매
\kfChoice 부동산 임대차 계약
\kfChoice 일상적인 거래
\kfChoice 장기 할부 계약
\end{kfChoices}
\end{kfSentenceUnit}

\kfEndMC

\kfActivityBg \par\kfMaybeNeedspace{32\baselineskip}\noindent{\kfTipMark\,\,}{\small\color{kfInk} 글의 전체적인 논리 흐름을 파악하며 빈칸에 알맞은 내용을 채워 보세요.}\par\nopagebreak[4]\smallskip\nopagebreak[4]
\begin{kfStructure}
\kfNode{0}{}{}
\end{kfStructure}

\kfSecHeader{kfAreaNonlit}{문제}{}

\par\medskip
\kfBeginMC
\begin{kfQBlock}{01}{객관식}
\kfQStem{윗글에 대한 설명으로 적절하지 \uline{\underline{않은}} 것은?}
\begin{kfChoices}
\kfChoice 청소년 구매력 증가가 시장에 가져온 긍정적 변화를 통계로 보여 준다.
\kfChoice 계약 취소가 판매자에게 줄 수 있는 손해 문제를 함께 짚는다.
\kfChoice 판매자를 보호하기 위한 ‘확답 촉구권’과 ‘철회권’ 등을 소개한다.
\kfChoice 미성년자 계약 취소권과 판매자 보호 제도를 균형 있게 설명한다.
\kfChoice 미성년자가 부모 동의 없이 한 계약을 취소할 수 있는 권리를 다룬다.
\end{kfChoices}
\end{kfQBlock}
\begin{kfQBlock}{02}{객관식}
\kfQStem{윗글의 내용과 일치하지 \uline{않는} 것은?}
\begin{kfChoices}
\kfChoice 미성년자가 법정대리인의 동의 없이 한 계약은 취소할 수 있다.
\kfChoice 계약이 취소되면 미성년자는 사용한 물건의 가치만큼 돈을 내야 한다.
\kfChoice 미성년자가 자신을 성년이라고 속인 경우, 판매자는 계약을 취소하지 않을 수 있다.
\kfChoice 판매자는 미성년자의 부모에게 계약을 인정할 것인지 물어볼 수 있다.
\kfChoice 미성년자도 교통카드를 충전하고 사용하는 등의 일상적인 거래는 직접 할 수 있다.
\end{kfChoices}
\end{kfQBlock}
\begin{kfQBlock}{03}{객관식}
\kfQStem{㉠\textasciitilde{}㉢에 대한 설명으로 적절하지 \uline{\underline{않은}} 것은?}
\begin{kfEvidence}㉠ 취소권 배제 \\\relax
㉡ 확답 촉구권 \\\relax
㉢ 철회권\end{kfEvidence}
\begin{kfChoices}
\kfChoice ㉠은 미성년자의 속임수가 있었을 때 판매자를 보호하기 위한 제도이다.
\kfChoice ㉡은 판매자가 계약의 불확실한 상태를 해소하기 위해 사용할 수 있다.
\kfChoice ㉢은 판매자가 미성년자임을 몰랐을 경우에 계약에서 먼저 벗어날 수 있게 한다.
\kfChoice ㉡과 ㉢은 판매자가 미성년자의 법정대리인에게 행사하는 권리에 해당한다.
\kfChoice ㉠, ㉡, ㉢은 모두 미성년자의 취소권 남용으로부터 판매자를 보호하는 역할을 한다.
\end{kfChoices}
\end{kfQBlock}
\begin{kfQBlock}{04}{객관식}
\kfQStem{<보기>의 상황에 윗글을 적용한 판단으로 적절하지 \uline{\underline{않은}} 것은?}
\begin{kfEvidence}18세 A는 부모님 몰래 100만 원짜리 최신형 스마트폰을 구매하여 사용했다. 한 달 뒤 이 사실을 알게 된 A의 부모는 계약을 취소하려고 한다. A는 그동안 스마트폰을 사용하여 약간의 흠집이 생겼다.\end{kfEvidence}
\begin{kfChoices}
\kfChoice A의 부모는 흠집 난 스마트폰을 반납하고 이미 낸 돈을 돌려받을 수 있다.
\kfChoice A가 스마트폰을 한 번이라도 사용했다면 계약 취소 자체가 불가능하다.
\kfChoice 판매자는 A의 부모에게 ‘확답 촉구권’으로 계약 인정 여부를 물을 수 있다.
\kfChoice 한 달간 사용한 ‘사용 이익’은 따로 돌려주지 않아도 된다.
\kfChoice A가 미성년자임을 알고 판매한 판매자는 ‘철회권’을 행사할 수 없다.
\end{kfChoices}
\end{kfQBlock}
\begin{kfQBlock}{05}{객관식}
\kfQStem{사례의 학생이 추가로 돈을 내지 않아도 되는 법적 근거로 가장 타당한 것은?}
\begin{kfChoices}
\kfChoice 미성년자는 경제 활동의 주체가 될 수 없기 때문이다.
\kfChoice 다이어트 식품이 학생의 건강에 해로운 영향을 미쳤기 때문이다.
\kfChoice 판매업자가 학생에게 상품의 효능을 과장하여 광고했기 때문이다.
\kfChoice 미성년자는 현존하는 이익의 한도 내에서만 반환할 책임이 있기 때문이다.
\kfChoice 모든 계약은 14일 이내에 아무런 조건 없이 취소할 수 있기 때문이다.
\end{kfChoices}
\end{kfQBlock}
\begin{kfQBlock}{06}{단답형}
\kfQStem{우리나라 민법에서 정하는 미성년자의 기준 나이는 몇 세 미만인가?}
\kfAnswerLines{1}
\end{kfQBlock}
\begin{kfQBlock}{07}{단답형}
\kfQStem{미성년자 보호 제도의 목적은 미성년자가 어른보다 부족한 무엇 때문에 손해 보는 계약을 하지 않도록 하기 위함인가? (2가지)}
\kfAnswerLines{1}
\end{kfQBlock}
\begin{kfQBlock}{08}{단답형}
\kfQStem{미성년자가 계약을 취소할 때, 반환해야 할 돈의 범위를 무엇이라고 하는가?}
\kfAnswerLines{1}
\end{kfQBlock}
\begin{kfQBlock}{09}{단답형}
\kfQStem{미성년자가 거짓말로 성년인 척 속여서 계약한 경우, 계약 취소 권리가 없어지는 제도는 무엇인가?}
\kfAnswerLines{1}
\end{kfQBlock}
\begin{kfQBlock}{10}{단답형}
\kfQStem{판매업자가 미성년자의 부모에게 계약 취소 여부를 물어볼 수 있는 권리는 무엇인가?}
\kfAnswerLines{1}
\end{kfQBlock}
\begin{kfQBlock}{11}{서술형}
\kfQStem{본문 첫 문단의 사례에서, 학생과 판매업자가 각각 어떻게 행동해야 하는지 그 법적 근거와 함께 서술하시오.}
\kfAnswerNote{답안}{4}
\end{kfQBlock}
\kfEndMC



\clearpage

\kfChapterCover{3}{러셀2 (챕터3)}{Chapter 3}{이번 챕터의 학습 내용을 시작합니다.}
\begin{kfChapterAreaList}
\kfChapterAreaItem{문법}{kfAreaGrammar}{활용하는 용언, 꾸며주는 수식언, 홀로 쓰는 독립언}
\kfChapterAreaItem{개념}{kfAreaConcept}{마음속에 그려지는 그림, 심상}
\kfChapterAreaItem{문학}{kfAreaLiterature}{문장 독해 — 발췌 문장}
\kfChapterAreaItem{비문학}{kfAreaNonlit}{[과학] 중력을 이기는 식물의 물 운반 비밀}
\end{kfChapterAreaList}

\kfStartGrammar{활용하는 용언, 꾸며주는 수식언, 홀로 쓰는 독립언}


\kfPassageNoteNT{%
　문장이라는 요리의 핵심 재료인 체언과, 재료 간의 관계를 정해주는 관계언을 배웠다면, 이번에는 그 요리의 맛을 더하고 풍미를 살리는 단어들을 만나볼 차례이다. 움직임이나 상태를 서술하는 용언, 다른 말을 꾸며 의미를 풍부하게 하는 수식언, 그리고 문장 속에서 홀로 쓰이는 독립언이 바로 그 주인공이다.

　용언은 문장의 서술어 역할을 담당하며 동작이나 상태, 성질을 풀이하는 단어 그룹이다. 용언의 가장 큰 특징은 문장 속에서 형태가 고정되지 않고 어미가 바뀌며 다양하게 쓰인다는 점인데, 이를 ‘활용’이라고 한다. 예를 들어 ‘먹다’가 ‘먹고, 먹으니, 먹어서’ 등으로 변하는 것이다. 용언은 의미에 따라 두 갈래로 나뉜다. ‘달리다’, ‘웃다’처럼 움직임이나 작용을 나타내 ‘어찌한다’에 해당하는 것을 동사라 하고, ‘예쁘다’, ‘착하다’와 같이 성질이나 상태를 나타내 ‘어떠하다’에 해당하는 것을 형용사라 한다.

　수식언은 이름 그대로 다른 단어를 꾸며주는 역할을 전문적으로 수행한다. 수식언은 활용을 하지 않아 형태가 고정되어 있으며, 무엇을 꾸미는지에 따라 관형사와 부사로 나뉜다. 관형사는 ‘새 신발’, ‘온 세상’, ‘이 사람’의 ‘새’, ‘온’, ‘이’처럼 명사, 대명사, 수사 같은 체언 앞에서 ‘어떤’이라는 의미를 더해 체언을 꾸민다. 반면 부사는 주로 동사나 형용사 같은 용언을 꾸며 ‘어떻게, 언제, 어디서’ 등의 의미를 더하지만, 때로는 다른 부사나 관형사, 혹은 문장 전체를 꾸미기도 하는 꾸미기의 멀티플레이어다. ‘매우 높다’의 ‘매우’, ‘빨리 달린다’의 ‘빨리’ 그리고 ‘과연 그는 똑똑하다’의 ‘과연’이 모두 부사에 해당한다.

　독립언은 문장 안의 다른 성분들과 직접적인 문법적 관계를 맺지 않고 독립적으로 쓰이는 단어이다. 여기에는 감탄사가 유일하게 속한다. ‘아!’, ‘여보세요’, ‘네’처럼 말하는 이의 놀람이나 느낌, 부름, 대답 등을 나타내는 말이 감탄사이며, 보통 문장의 맨 앞에 놓이는 경우가 많다. 이처럼 용언, 수식언, 독립언은 각자의 자리에서 문장을 더욱 정교하고 다채롭게 만드는 필수적인 요소이다.
\par\medskip\begin{itemize}[leftmargin=18pt, itemsep=2pt, topsep=2pt, label=\textbullet]
\item 용언(활용) — '먹다→먹고·먹으니·먹어서' / 동·형(달리다·예쁘다)
\item 수식언 — 관형사(새·온·이) / 부사(매우·빨리·과연)
\item 독립언 — '아!·여보세요·네'(감탄사)
\end{itemize}
}

\kfActivityBg \par\kfMaybeNeedspace{32\baselineskip}\noindent{\kfTipMark\,\,}{\small\color{kfInk} 지문의 핵심 내용을 떠올리며 빈칸에 알맞은 말을 채워 보세요.}\par\nopagebreak[4]\smallskip\nopagebreak[4]

\kfActivityBg \kfSecHeader{kfAreaGrammar}{활동}{분석 훈련}
\par\kfMaybeNeedspace{24\baselineskip}\noindent{\kfTipMark\,\,}{\small\color{kfInk} 다음 문장을 읽고, 설명에 해당하는 단어를 모두 찾아 쓰시오.}\par\nopagebreak[4]\smallskip\nopagebreak[4]
\begin{enumerate}[leftmargin=22pt, itemsep=6pt, topsep=4pt, label=\textbf{\arabic*.}]
\item 다음 문장에서 주어의 움직임이나 작용을 나타내는 동사를 모두 찾아 기본형으로 쓰시오.
\begin{kfEvidence}나는 어제 친구를 만나서 밥을 먹었다.\end{kfEvidence}
\item 다음 문장에서 주어의 상태나 성질을 나타내는 형용사를 모두 찾아 기본형으로 쓰시오.
\begin{kfEvidence}가을 하늘은 높고 푸르며 공기가 맑다.\end{kfEvidence}
\item 다음 문장에서 체언(명사, 대명사, 수사)을 꾸며주는 관형사를 모두 찾아 쓰시오.
\begin{kfEvidence}그 학생은 헌 책상 위에 놓인 모든 책을 읽었다.\end{kfEvidence}
\item 다음 문장에서 용언(동사, 형용사)이나 다른 말을 꾸며주는 부사를 모두 찾아 쓰시오.
\begin{kfEvidence}차가 아주 빨리 달리니 금방 도착하겠다.\end{kfEvidence}
\item 다음 문장에서 놀람이나 부름, 대답을 나타내는 감탄사를 찾아 쓰시오.
\begin{kfEvidence}어머나, 벌써 갔구나! 응, 나도 곧 갈게.\end{kfEvidence}
\item 다음 문장에서 형태가 변하는 용언을 모두 찾아 기본형으로 쓰시오.
\begin{kfEvidence}날씨가 좋으니, 우리는 즐겁게 놀았다.\end{kfEvidence}
\item 다음 문장에서 형태가 변하지 않고 다른 말을 꾸며주는 수식언을 모두 찾아 쓰시오.
\begin{kfEvidence}저 새 자동차는 과연 매우 잘 달린다.\end{kfEvidence}
\item 다음 문장에서 움직임을 나타내는 동사와, 상태를 나타내는 형용사를 각각 찾아 기본형으로 쓰시오.
\begin{kfEvidence}맑은 하늘에 비행기가 빠르게 날아간다.\end{kfEvidence}
\item 다음 문장에서 체언을 꾸미는 관형사와 용언을 꾸미는 부사를 구분하여 찾으시오.
\begin{kfEvidence}한 아이가 문득 옛 생각을 떠올렸다.\end{kfEvidence}
\item 다음 문장에서 독립적으로 쓰이는 독립언을 찾아 쓰시오.
\begin{kfEvidence}아이고, 이걸 어쩌지? 얘야, 내 말 좀 들어보렴.\end{kfEvidence}
\end{enumerate}

\kfSecHeader{kfAreaGrammar}{문제}{}

\par\medskip
\kfBeginMC
\begin{kfQBlock}{01}{객관식}
\kfQStem{다음 중 용언에 대한 설명으로 적절하지 \uline{\underline{않은}} 것은?}
\begin{kfChoices}
\kfChoice 주어의 움직임이나 상태를 풀이해 주는 서술어 역할을 한다.
\kfChoice 크게 동사와 형용사 두 가지 품사가 여기에 속한다.
\kfChoice 문장에서 활용을 통해 형태가 자유롭게 변하는 가변어이다.
\kfChoice 어간에 다양한 어미가 붙어 여러 문법적 기능을 수행한다.
\kfChoice 형태가 변하지 않는 불변어로, 체언을 꾸며 주는 역할을 한다.
\end{kfChoices}
\end{kfQBlock}
\begin{kfQBlock}{02}{객관식}
\kfQStem{다음 중 품사의 종류가 \uline{다른} 하나는?}
\begin{kfChoices}
\kfChoice 슬프다
\kfChoice 읽다
\kfChoice 깨끗하다
\kfChoice 작다
\kfChoice 높다
\end{kfChoices}
\end{kfQBlock}
\begin{kfQBlock}{03}{객관식}
\kfQStem{밑줄 친 단어 중 품사가 관형사인 것은?}
\begin{kfChoices}
\kfChoice \uline{빨리} 달린다.
\kfChoice 하늘이 \uline{매우} 파랗다.
\kfChoice \uline{저} 사람이 내 친구다.
\kfChoice \uline{와}, 정말 멋지다!
\kfChoice 밥을 \uline{먹고} 간다.
\end{kfChoices}
\end{kfQBlock}
\begin{kfQBlock}{04}{객관식}
\kfQStem{다음 단어 중 문장에서 쓰일 때 형태가 변하는 '활용'을 하지 \uline{않는} 것은?}
\begin{kfChoices}
\kfChoice 자다
\kfChoice 좋다
\kfChoice 빨리
\kfChoice 푸르다
\kfChoice 읽다
\end{kfChoices}
\end{kfQBlock}
\begin{kfQBlock}{05}{객관식}
\kfQStem{다음 중 수식언이 \uline{\underline{아닌}} 것은?}
\begin{kfChoices}
\kfChoice 새
\kfChoice 헌
\kfChoice 매우
\kfChoice 우리
\kfChoice 과연
\end{kfChoices}
\end{kfQBlock}
\begin{kfQBlock}{06}{객관식}
\kfQStem{밑줄 친 단어의 품사가 부사가 \uline{\underline{아닌}} 것은?}
\begin{kfChoices}
\kfChoice 날씨가 \uline{정말} 덥다.
\kfChoice 밥을 \uline{빨리} 먹어라.
\kfChoice \uline{새} 학기가 시작되었다.
\kfChoice \uline{아주} 높이 나는 새.
\kfChoice \uline{결코} 포기하지 않겠다.
\end{kfChoices}
\end{kfQBlock}
\begin{kfQBlock}{07}{객관식}
\kfQStem{다음 중 독립언에 해당하는 단어는?}
\begin{kfChoices}
\kfChoice 그리고
\kfChoice 여보세요
\kfChoice 에게
\kfChoice 모든
\kfChoice 생각하다
\end{kfChoices}
\end{kfQBlock}
\begin{kfQBlock}{08}{객관식}
\kfQStem{‘활용’에 대한 설명으로 적절하지 \uline{\underline{않은}} 것은?}
\begin{kfChoices}
\kfChoice 용언의 어간에 다양한 어미가 결합하여 형태가 변하는 현상이다.
\kfChoice 동사와 형용사, 서술격 조사 ‘이다’에서 일어난다.
\kfChoice 어미가 교체되며 문법적 기능이 달라진다.
\kfChoice ‘먹다 → 먹고·먹어·먹으니’처럼 어미만 바뀌는 모습으로 나타난다.
\kfChoice 명사·대명사·수사 등 모든 품사에서 공통적으로 일어나는 현상이다.
\end{kfChoices}
\end{kfQBlock}
\begin{kfQBlock}{09}{객관식}
\kfQStem{다음 문장에 모두 포함된 품사끼리 짝지은 것은?}
\begin{kfEvidence}와, 저 아이는 정말 착하구나!\end{kfEvidence}
\begin{kfChoices}
\kfChoice 명사, 동사
\kfChoice 관형사, 동사
\kfChoice 부사, 수사
\kfChoice 대명사, 조사
\kfChoice 감탄사, 형용사
\end{kfChoices}
\end{kfQBlock}
\begin{kfQBlock}{10}{객관식}
\kfQStem{밑줄 친 단어의 품사 구별이 바른 것은?}
\begin{kfChoices}
\kfChoice \uline{한} 사람이 온다. (수사)
\kfChoice 그는 \uline{매우} 성실하다. (관형사)
\kfChoice \uline{아니}, 그럴 리가! (감탄사)
\kfChoice 산이 \uline{높다}. (동사)
\kfChoice \uline{웃음}을 참기 힘들다. (동사)
\end{kfChoices}
\end{kfQBlock}
\begin{kfQBlock}{11}{객관식}
\kfQStem{다음 중 문장에서 독립적으로 쓰일 수 있는 품사는?}
\begin{kfChoices}
\kfChoice 명사
\kfChoice 조사
\kfChoice 부사
\kfChoice 감탄사
\kfChoice 형용사
\end{kfChoices}
\end{kfQBlock}
\begin{kfQBlock}{12}{객관식}
\kfQStem{밑줄 친 단어 중 성격이 \uline{다른} 하나는?}
\begin{kfChoices}
\kfChoice \uline{달리는} 기차
\kfChoice \uline{예쁜} 얼굴
\kfChoice \uline{작은} 거인
\kfChoice \uline{밝은} 미래
\kfChoice \uline{맑은} 하늘
\end{kfChoices}
\end{kfQBlock}
\begin{kfQBlock}{13}{객관식}
\kfQStem{"체언을 꾸며준다."라는 설명에 해당하는 품사는?}
\begin{kfChoices}
\kfChoice 부사
\kfChoice 관형사
\kfChoice 감탄사
\kfChoice 동사
\kfChoice 조사
\end{kfChoices}
\end{kfQBlock}
\begin{kfQBlock}{14}{객관식}
\kfQStem{<보기>의 빈칸에 들어갈 말로 적절한 것은?}
\begin{kfEvidence}* 수식언은 활용을 \kfFillBlank[20mm]{}. \\\relax
* 문장 속에서 형태가 고정되어 있다.\end{kfEvidence}
\begin{kfChoices}
\kfChoice 한다
\kfChoice 하지 않는다
\kfChoice 할 수도 있다
\kfChoice 해야만 한다
\kfChoice 하는 것이 좋다
\end{kfChoices}
\end{kfQBlock}
\begin{kfQBlock}{15}{객관식}
\kfQStem{다음 문장에 대한 분석으로 옳지 \uline{\underline{않은}} 것은?}
\begin{kfEvidence}아, 그 새 옷은 정말 예쁘구나!\end{kfEvidence}
\begin{kfChoices}
\kfChoice 독립언 '아'가 사용되었다.
\kfChoice 수식언은 '그, 새, 정말' 3개이다.
\kfChoice 용언 '예쁘구나'는 형용사이다.
\kfChoice 체언은 '그, 옷' 2개이다.
\kfChoice 형태가 변하는 가변어가 사용되었다.
\end{kfChoices}
\end{kfQBlock}
\begin{kfQBlock}{16}{단답형}
\kfQStem{문장에서 주어의 동작이나 상태를 풀이하며 '활용'을 하는 단어의 묶음을 무엇이라고 하는가?}
\kfAnswerLines{1}
\end{kfQBlock}
\begin{kfQBlock}{17}{단답형}
\kfQStem{용언이 문장 속에서 어미를 바꾸는 것을 무엇이라고 하는가?}
\kfAnswerLines{1}
\end{kfQBlock}
\begin{kfQBlock}{18}{단답형}
\kfQStem{'걷다', '보다'처럼 주어의 움직임을 나타내는 품사는 무엇인가?}
\kfAnswerLines{1}
\end{kfQBlock}
\begin{kfQBlock}{19}{단답형}
\kfQStem{'예쁘다', '덥다'처럼 주어의 상태나 성질을 나타내는 품사는 무엇인가?}
\kfAnswerLines{1}
\end{kfQBlock}
\begin{kfQBlock}{20}{단답형}
\kfQStem{다른 말을 꾸며주는 역할을 하는 관형사와 부사를 묶어 무엇이라고 하는가?}
\kfAnswerLines{1}
\end{kfQBlock}
\begin{kfQBlock}{21}{단답형}
\kfQStem{'새', '헌', '모든'처럼 체언을 꾸며주는 품사는 무엇인가?}
\kfAnswerLines{1}
\end{kfQBlock}
\begin{kfQBlock}{22}{단답형}
\kfQStem{'아주', '매우', '빨리'처럼 주로 용언을 꾸며주는 품사는 무엇인가?}
\kfAnswerLines{1}
\end{kfQBlock}
\begin{kfQBlock}{23}{단답형}
\kfQStem{문장의 다른 성분과 관계없이 독립적으로 쓰이는 단어의 묶음을 무엇이라고 하는가?}
\kfAnswerLines{1}
\end{kfQBlock}
\begin{kfQBlock}{24}{단답형}
\kfQStem{'와', '네', '아니요'처럼 놀람, 느낌, 대답을 나타내는 품사는 무엇인가?}
\kfAnswerLines{1}
\end{kfQBlock}
\begin{kfQBlock}{25}{단답형}
\kfQStem{‘어찌한다?’라는 물음에 답이 되며, 움직임을 나타내는 품사는 무엇인가?}
\kfAnswerLines{1}
\end{kfQBlock}
\begin{kfQBlock}{26}{서술형}
\kfQStem{다음 문장에서 용언, 수식언, 독립언에 해당하는 단어를 각각 모두 찾아 쓰시오.}
\begin{kfEvidence}어머나, 저 헌 자동차가 아주 빨리 달린다!\end{kfEvidence}
\kfAnswerNote{답안}{4}
\end{kfQBlock}
\kfEndMC


\kfStartConcept{마음속에 그려지는 그림, 심상}


\kfPassageNoteNT{%
　문학 작품을 읽을 때, 머릿속에 생생한 그림이 그려지거나 소리가 들리는 듯한 감각적 경험을 하게 된다. 이처럼 언어를 통해 우리 마음속에 떠오르는 구체적이고 감각적인 인상을 심상, 즉 '마음의 그림'이라고 한다. 심상은 감각의 종류에 따라 시각적 심상, 청각적 심상, 후각적 심상, 미각적 심상, 촉각적 심상, 공감각적 심상으로 나뉜다.

　시각적 심상은 색채, 모양, 움직임 등 눈으로 보는 듯한 느낌을 주는 것이다. 곽재구의 「사평역에서」에서  "청색의 손바닥을 불빛 속에 적셔두고"는 추위에 얼어 파랗게 보이는 손의 '청색'과 난로의 '불빛'이 선명한 색채 대비를 통해 힘겨운 삶을 시각적으로 형상화한다.

　청각적 심상은 소리를 귀로 듣는 듯한 느낌을 주는 것이다. 김소월의 「가는 길」에서 "저 산에도 까마귀, 들에 까마귀, 서산(西山)에는 해 진다고 지저귑니다."는 까마귀 울음소리로 떠나기 싫은 사람의 아쉬운 마음을 강조한다.

　후각적 심상은 냄새를 코로 맡는 듯한 느낌을 주는 것이다. 이육사의 「광야」에서  "지금 눈 내리고 / 매화 향기 홀로 아득하니"는 추운 계절 속에서도 피어나는 '매화 향기'를 통해 고결한 이미지를 후각적으로 형상화했다.

　미각적 심상은 맛을 혀로 느끼는 듯한 느낌을 주는 것이다. 정지용의 「고향」에서 "메마른 입술이 쓰디쓰다"는 '쓴맛'을 통해 화자의 괴롭고 힘든 감정을 전달한다.

　촉각적 심상은 온도나 감촉 등 피부로 느끼는 듯한 느낌을 주는 것이다. 이상화의 「빼앗긴 들에도 봄은 오는가」에서 "발목이 시리도록 밟아도 보고"는 '시리다'라는 차가운 감각을 통해 화자가 딛고 선 땅의 느낌을 생생하게 전달한다.

　공감각적 심상은 하나의 감각을 다른 종류의 감각으로 바꾸어 표현하는 것이다. 김광균의 「외인촌」에서 "분수처럼 흩어지는 푸른 종소리"는 청각적 대상인 ‘종소리’를 ‘푸른’이라는 시각적 표현으로 옮겨 표현한 것이다. 

　심상은 작가가 전달하고자 하는 정서와 분위기를 효과적으로 표현하며, 독자가 작품을 생생하게 체험하도록 돕는다.
\par\medskip\begin{itemize}[leftmargin=18pt, itemsep=2pt, topsep=2pt, label=\textbullet]
\item 시각·청각 — 곽재구 「사평역에서」의 '청색 손바닥·불빛' / 「가는 길」의 '까마귀 지저귐'
\item 후각·미각 — 이육사 「광야」의 '매화 향기' / 정지용 「고향」의 '쓰디쓴 입술'
\item 공감각 — 김광균 「외인촌」의 '분수처럼 흩어지는 푸른 종소리'(청→시)
\end{itemize}
}

\kfSecHeader{kfAreaConcept}{문제}{}

\par\medskip
\kfBeginMC
\begin{kfQBlock}{01}{객관식}
\kfQStem{<보기>의 시구에서 중심적으로 사용된 심상은?}
\begin{kfEvidence}해야 솟아라. 해야 솟아라. 말갛게 씻은 얼굴 고운 해야 솟아라.

- 박두진, 「해」\end{kfEvidence}
\begin{kfChoices}
\kfChoice 시각적 심상
\kfChoice 청각적 심상
\kfChoice 후각적 심상
\kfChoice 미각적 심상
\kfChoice 촉각적 심상
\end{kfChoices}
\end{kfQBlock}
\begin{kfQBlock}{02}{객관식}
\kfQStem{<보기>의 시구에서 중심적으로 사용된 심상은?}
\begin{kfEvidence}바람맛도 짭짤한 물맛도 짭짤한

- 백석, 「통영」\end{kfEvidence}
\begin{kfChoices}
\kfChoice 시각적 심상
\kfChoice 청각적 심상
\kfChoice 후각적 심상
\kfChoice 미각적 심상
\kfChoice 촉각적 심상
\end{kfChoices}
\end{kfQBlock}
\begin{kfQBlock}{03}{객관식}
\kfQStem{<보기>의 시 구절에서 공통적으로 사용된 심상은 무엇인가?}
\begin{kfEvidence}• 서늘한 옷자락에 열로 상기한 볼을 말없이 부비는 것이었다.

 - 김종길, 「성탄제」

• 유리창엔 차고 슬픈 것이 어른거린다.

 - 정지용 「유리창 1」\end{kfEvidence}
\begin{kfChoices}
\kfChoice 촉각적 심상
\kfChoice 청각적 심상
\kfChoice 후각적 심상
\kfChoice 미각적 심상
\kfChoice 시각적 심상
\end{kfChoices}
\end{kfQBlock}
\begin{kfQBlock}{04}{객관식}
\kfQStem{<보기>의 (가)\textasciitilde{}(마)에 대한 설명으로 적절하지 \uline{\underline{않은}} 것은?}
\begin{kfEvidence}(가) 청색의 손바닥을 불빛 속에 적셔두고

(나) 저 산에도 까마귀, 들에 까마귀

(다) 매화 향기 홀로 아득하니

(라) 메마른 입술이 쓰디쓰다

(마) 발목이 시리도록 밟아도 보고\end{kfEvidence}
\begin{kfChoices}
\kfChoice (가) ‘청색의 손바닥’은 색채를 나타내는 시각적 심상이다.
\kfChoice (마) ‘발목이 시리도록’은 미각적 심상이다.
\kfChoice (다) ‘매화 향기’는 후각적 심상이다.
\kfChoice (라) ‘쓰디쓰다’는 미각적 심상이다.
\kfChoice (나) ‘까마귀’의 울음소리는 청각적 심상이다.
\end{kfChoices}
\end{kfQBlock}
\begin{kfQBlock}{05}{객관식}
\kfQStem{<보기>의 심상에 대한 분석으로 가장 적절한 것은?}
\begin{kfEvidence}금빛 게으른 울음을 우는 곳

- 정지용 「향수」\end{kfEvidence}
\begin{kfChoices}
\kfChoice 시각을 청각으로 바꾸어 표현했다.
\kfChoice 청각을 시각으로 바꾸어 표현했다.
\kfChoice 후각을 촉각으로 바꾸어 표현했다.
\kfChoice 촉각을 후각으로 바꾸어 표현했다.
\kfChoice 미각을 시각으로 바꾸어 표현했다.
\end{kfChoices}
\end{kfQBlock}
\begin{kfQBlock}{06}{객관식}
\kfQStem{<보기>의 시에 대한 감상으로 적절하지 \uline{\underline{않은}} 것은?}
\begin{kfEvidence}해와 하늘빛이

문둥이는 서러워

보리밭에 달 뜨면

애기 하나 먹고

꽃처럼 붉은 울음을 밤새 울었다.

- 서정주 「문둥이」\end{kfEvidence}
\begin{kfChoices}
\kfChoice ‘붉은 울음’은 공감각적 심상을 활용한 표현이군.
\kfChoice 청각 대상인 ‘울음’을 시각적으로 형상화한 것이군.
\kfChoice ‘꽃처럼’이란 표현을 보니 후각적 심상이 두드러지는군.
\kfChoice 시각 이미지를 통해 화자의 슬픔을 강렬히 전달하고 있군.
\kfChoice ‘보리밭에 달 뜨면’에서 서정적인 시각적 심상이 드러나는군.
\end{kfChoices}
\end{kfQBlock}
\begin{kfQBlock}{07}{객관식}
\kfQStem{다음 중 심상의 종류가 나머지 넷과 \uline{다른} 하나는?}
\begin{kfChoices}
\kfChoice 새벽부터 돌 깨는 산울림에 떨다가
\kfChoice 내 마음은 호수요,  그대 저어 오오.
\kfChoice 길은 구겨진 넥타이처럼 한 줄로 풀려 있다.
\kfChoice 아 아버지가 눈을 헤치고 따 오신 그 붉은 산수유 열매
\kfChoice 해야 솟아라. 해야 솟아라. 말갛게 씻은 얼굴 고운 해야 솟아라.
\end{kfChoices}
\end{kfQBlock}
\begin{kfQBlock}{08}{객관식}
\kfQStem{<보기>의 밑줄 친 부분과 동일한 유형의 심상이 사용된 것은?}
\begin{kfEvidence}아, 강낭콩꽃보다도 더 푸른

그 물결 위에

양귀비꽃보다도 더 붉은

그 마음 흘러라.

- 변영로, 「논개」\end{kfEvidence}
\begin{kfChoices}
\kfChoice 매운 계절의 채찍에 갈겨
\kfChoice 얼음장같이 차가운 방바닥
\kfChoice 썩은 땀의 냄새에 절어 있었다.
\kfChoice 쨍그랑, 하고 깨지는 유리병 소리
\kfChoice 눈을 감으면 환히 보이는 그대 얼굴
\end{kfChoices}
\end{kfQBlock}
\begin{kfQBlock}{09}{객관식}
\kfQStem{<보기>의 시 구절에서 공통적으로 활용된 심상과 그 효과로 가장 적절한 것은?}
\begin{kfEvidence}• 어두운 방 안엔 바알간 숯불이 피고

 - 김종길, 「성탄제」

• 자하문 밖으론 초승달이 하얗게 걸려 있다.

 - 김광균, 「데생」\end{kfEvidence}
\begin{kfChoices}
\kfChoice 청각적 심상을 통해 고요한 분위기를 조성한다.
\kfChoice 후각적 심상을 통해 계절의 변화를 암시한다.
\kfChoice 촉각적 심상을 통해 대상의 질감을 강조한다.
\kfChoice 미각적 심상을 통해 화자의 정서를 환기한다.
\kfChoice 시각적 심상을 통해 선명한 인상을 전달한다.
\end{kfChoices}
\end{kfQBlock}
\begin{kfQBlock}{10}{단답형}
\kfQStem{언어를 통해 우리 마음속에 떠오르는 구체적이고 감각적인 인상을 무엇이라고 하는가?}
\kfAnswerLines{1}
\end{kfQBlock}
\begin{kfQBlock}{11}{단답형}
\kfQStem{색채, 모양, 움직임 등 눈으로 보는 듯한 느낌을 주는 심상은 무엇인가?}
\kfAnswerLines{1}
\end{kfQBlock}
\begin{kfQBlock}{12}{단답형}
\kfQStem{소리를 귀로 듣는 듯한 느낌을 주는 심상은 무엇인가?}
\kfAnswerLines{1}
\end{kfQBlock}
\begin{kfQBlock}{13}{단답형}
\kfQStem{맛을 혀로 느끼는 듯한 느낌을 주는 심상은 무엇인가?}
\kfAnswerLines{1}
\end{kfQBlock}
\begin{kfQBlock}{14}{단답형}
\kfQStem{냄새를 코로 맡는 듯한 느낌을 주는 심상은 무엇인가?}
\kfAnswerLines{1}
\end{kfQBlock}
\begin{kfQBlock}{15}{단답형}
\kfQStem{온도나 감촉 등 피부로 느끼는 듯한 느낌을 주는 심상은 무엇인가?}
\kfAnswerLines{1}
\end{kfQBlock}
\begin{kfQBlock}{16}{단답형}
\kfQStem{하나의 감각을 다른 종류의 감각으로 바꾸어 표현하는 심상은 무엇인가?}
\kfAnswerLines{1}
\end{kfQBlock}
\begin{kfQBlock}{17}{서술형}
\kfQStem{<보기>의 시 구절에서 중심적으로 사용된 심상이 무엇인지 쓰고, 그 심상이 시의 분위기를 어떻게 만들어 주는지 서술하시오.}
\begin{kfEvidence}해야 솟아라. 해야 솟아라. \\\relax
말갛게 씻은 얼굴 고운 해야 솟아라. \\\relax
산 넘어 산 넘어서 어둠을 살라 먹고, \\\relax
산 넘어서 밤새도록 어둠을 살라 먹고, \\\relax
이글이글 애띤 얼굴 고운 해야 솟아라. \\\relax
- 박두진, 「해」\end{kfEvidence}
\kfAnswerNote{답안}{4}
\end{kfQBlock}
\kfEndMC


\kfStartLit{문장 독해 — 발췌 문장}


\kfPassageNote{}{%
열무 삼십 단을 이고 \\[4pt]
시장에 간 우리 엄마 \\[4pt]
안 오시네, 해는 시든 지 오래 \\[4pt]
나는 찬밥처럼 방에 담겨 \\[4pt]
아무리 천천히 숙제를 해도 \\[4pt]
엄마 안 오시네, 배춧잎 같은 발소리 타박타박 \\[4pt]
안 들리네, 어둡고 무서워 \\[4pt]
금간 창틈으로 고요히 빗소리 \\[4pt]
빈방에 혼자 엎드려 훌쩍거리던 \\[14pt]
아주 먼 옛날 \\[4pt]
지금도 내 눈시울을 뜨겁게 하는 \\[4pt]
그 시절, 내 유년의 윗목
}


\kfPassageNote{문장 독해 — 발췌 문장}{%
1. 3행 '안 오시네', 6행 '안 오시네', 7행 '안 들리네' \\[14pt]
2. 아무리 천천히 숙제를 해도 엄마 안 오시네 \\[14pt]
3. 찬밥처럼 \\[14pt]
4. 배춧잎 같은 발소리 \\[14pt]
5. 어둡고 무서워 \\[14pt]
6. 빈방에 혼자 엎드려 훌쩍거리던 \\[14pt]
7. 지금도 내 눈시울을 뜨겁게 하는 \\[14pt]
8. 그 시절
}


\kfPassageNote{해설 학습 — 본문}{%
[단원 1] 이 시는 어른이 된 화자가 어린 시절의 특정 장면을 회상하는 방식으로 구성된다. 시의 내용은 열무를 팔러 시장에 간 어머니를 어린아이가 해가 저물도록 방에서 홀로 기다리는 상황을 묘사하고 있다.

이 짧은 장면 속에는 아이의 여러 가지 감정이 담겨 있다. 엄마가 빨리 돌아왔으면 하는 간절한 기다림, 혼자 남겨진 외로움과 쓸쓸함, 그리고 날이 어두워지면서 느끼는 두려움이다. 

화자는 이 기억을 ‘아주 먼 옛날’의 일이라고 말하지만, 그 기억은 여전히 생생하게 남아 ‘지금도 내 눈시울을 뜨겁게’ 만든다. 

이처럼 이 시는 한 아이의 평범한 저녁 시간을 통해 엄마에 대한 그리움과 어린 시절의 아픈 기억이 한 사람의 마음에 얼마나 깊이 새겨질 수 있는지를 보여준다. 고된 삶을 살았던 어머니에 대한 안타까움과 사랑 또한 느낄 수 있는 작품이다.

[단원 2] 이 시는 감각적인 표현을 사용하여 아이의 감정을 생생하게 전달하는 것이 특징이다. 시인은 ‘아이가 외로웠다’라고 말하는 대신, 아이가 보고 듣고 느꼈던 것들을 생생하게 보여줌으로써 그 감정을 더욱 깊이 있게 전달한다.

먼저 시각적 이미지를 통해 시간이 흐르고 분위기가 어두워지는 것을 보여준다. ‘해는 시든 지 오래’라는 표현은 해가 진 지 오래되어 세상이 완전히 어두워졌음을 보여주며 아이의 불안감을 커지게 한다.

청각적 이미지 역시 아이의 외로움을 깊게 만든다. 아이가 간절히 기다리는 엄마의 발소리는 들리지 않고, ‘금간 창틈으로 고요히 빗소리’만이 들릴 뿐이다. 바깥세상과 단절된 방 안에서 들려오는 이 빗소리는 기다리는 대상을 더욱 애타게 만들고 아이의 슬픈 정서를 한층 더 깊게 만든다.

또한 촉각적 이미지는 아이의 서러운 처지를 느끼게 한다. ‘찬밥처럼 방에 담겨’ 있다는 표현은 혼자 남겨진 아이의 차가운 외로움을 느끼게 한다. 또한 ‘내 유년의 윗목’이라는 구절은 따뜻한 아랫목과 반대되는 공간으로, 따뜻한 온기가 부족했던 어린 시절의 서러움을 상징한다.

그뿐만 아니라 이 시는 하나의 감각을 다른 감각으로 바꾸어 표현하는 공감각적 심상을 사용하여 감정을 섬세하게 그려낸다. 엄마의 ‘배춧잎 같은 발소리’가 대표적인 예로, 지쳐서 힘없는 발소리인 청각 정보를 축 늘어진 배춧잎이라는 시각적 모습으로 표현한 것이다.

이처럼 시인은 시각, 청각, 촉각과 같은 다양한 감각적 이미지는 물론, 감각을 넘나드는 공감각적 표현까지 사용하여 아이가 처한 상황과 감정을 효과적으로 그려낸다.

[단원 3] 시는 감정이나 상황을 다른 대상에 빗대어 표현하는 비유법을 통해 독자에게 깊은 감동을 준다. 이 시는 특히 직유법과 은유법을 효과적으로 사용하여 화자의 정서를 섬세하게 드러낸다.

직유법은 ‘\textasciitilde{}같이’, ‘\textasciitilde{}처럼’ 등을 사용하여 두 대상을 직접 연결하는 방식이다. "나는 찬밥처럼 방에 담겨"라는 구절은 화자의 외로운 처지를 보여주는 대표적인 예다. 자신의 신세를 온기 없는 ‘찬밥’에 빗대어, 가족의 관심에서 벗어난 소외감과 서러움을 절실하게 표현했다. 또한 "배춧잎 같은 발소리"는 지친 어머니의 모습을 감각적으로 묘사한다. 힘없는 발소리를 축 늘어진 ‘배춧잎’에 비유하여, 고된 하루를 보낸 어머니의 피곤함을 짐작하게 한다.

은유법은 ‘A는 B이다’의 형태로 두 대상을 암시적으로 동일시하는 방법이다. "내 유년의 윗목"이라는 표현이 이에 해당한다. 이는 ‘내 유년은 윗목이다’라는 의미로, 온돌방의 차가운 공간인 ‘윗목’을 통해 사랑이 부족했던 가난하고 외로운 어린 시절을 상징적으로 드러낸다.

이처럼 시인은 직유와 은유라는 비유적 장치를 통해 화자의 슬픔, 어머니의 고단함, 불행했던 유년 시절의 기억을 압축적으로 그려낸다. 이러한 표현은 단순한 묘사를 넘어 시의 의미를 풍부하게 만들고, 독자가 화자의 감정에 깊이 공감하도록 이끄는 핵심 역할을 한다.
}


\kfActivityBg \par\needspace{15\baselineskip}
\kfSecHeader{kfAreaLiterature}{활동}{문장 독해}
\par\kfMaybeNeedspace{32\baselineskip}\noindent{\kfTipMark\,\,}{\small\color{kfInk} 다음은 작품에서 발췌한 문장입니다. 각 문장을 읽고 물음에 답하시오.}\par\nopagebreak[4]\smallskip\nopagebreak[4]
\kfBeginMC
\begin{kfSentenceUnit}{1}{}
\kfSentQ{1.}{다음 행에서 생략된 주어를 찾아 바르게 짝지은 것은? (3행 '안 오시네', 6행 '안 오시네', 7행 '안 들리네' 순서)}
\begin{kfChoices}
\kfChoice 엄마, 빗소리, 엄마
\kfChoice 엄마, 엄마, 빗소리
\kfChoice 나, 엄마, 엄마의 발소리
\kfChoice 엄마, 엄마, 엄마의 발소리
\end{kfChoices}
\end{kfSentenceUnit}

\begin{kfSentenceUnit}{2}{}
\kfSentQ{2.}{다음 문장에서 서술어 '해도'와 '안 오시네'의 주어를 바르게 짝지은 것은? (아무리 천천히 숙제를 해도 엄마 안 오시네)}
\begin{kfChoices}
\kfChoice 나, 엄마
\kfChoice 엄마, 나
\kfChoice 엄마, 빗소리
\kfChoice 엄마, 엄마의 발소리
\end{kfChoices}
\end{kfSentenceUnit}

\begin{kfSentenceUnit}{3}{}
\kfSentQ{3.}{다음 시구인 '찬밥처럼'에 사용된 비유 표현에서 '원관념'과 '보조 관념'을 바르게 짝지은 것은?}
\begin{kfChoices}
\kfChoice 나, 방
\kfChoice 나, 찬밥
\kfChoice 방, 찬밥
\kfChoice 찬밥, 어둠
\end{kfChoices}
\end{kfSentenceUnit}

\begin{kfSentenceUnit}{4}{}
\kfSentQ{4.}{다음 시구인 '배춧잎 같은 발소리'에 사용된 비유 표현에서 '원관념'과 '보조 관념'을 바르게 짝지은 것은?}
\begin{kfChoices}
\kfChoice 발소리, 숙제
\kfChoice 배춧잎, 엄마
\kfChoice 발소리, 배춧잎
\kfChoice 엄마의 발소리, 찬밥
\end{kfChoices}
\end{kfSentenceUnit}

\begin{kfSentenceUnit}{5}{}
\kfSentQ{5.}{다음 문장에서 서술어 '훌쩍거리던'의 주체로 알맞은 것은?}
\begin{kfChoices}
\kfChoice 나
\kfChoice 엄마
\kfChoice 빗소리
\kfChoice 발소리
\end{kfChoices}
\end{kfSentenceUnit}

\begin{kfSentenceUnit}{6}{}
\kfSentQ{6.}{에서 '내'가 가리키는 대상으로 가장 적절한 것은?}
\begin{kfChoices}
\kfChoice 나
\kfChoice 어린 아이
\kfChoice 어른이 된 나
\kfChoice 시장에 간 엄마
\end{kfChoices}
\end{kfSentenceUnit}

\begin{kfSentenceUnit}{7}{}
\kfSentQ{7.}{"그 시절"이라는 표현이 가리키는 때로 가장 적절한 것은?}
\begin{kfChoices}
\kfChoice 숙제를 하던 시절
\kfChoice 엄마를 기다리던 어린 시절
\kfChoice 시장에서 열무를 팔던 시절
\kfChoice 빗소리를 들으며 울던 어린 시절
\end{kfChoices}
\end{kfSentenceUnit}

\kfEndMC

\kfActivityBg \kfSecHeader{kfAreaLiterature}{활동}{해설 문제}
\par\kfMaybeNeedspace{18\baselineskip}\noindent{\kfTipMark\,\,}{\small\color{kfInk} 해설 본문을 읽고, 단원별 단답형 물음에 답하시오.}\par\nopagebreak[4]\smallskip\nopagebreak[4]
\kfBeginMC
\begin{enumerate}[leftmargin=22pt, itemsep=6pt, topsep=4pt, label=\textbf{\arabic*.}]
\item 이 시의 화자는 어떤 시점에서 과거를 돌아보고 있는가?
\item 화자가 떠올리는 것은 어떤 시절의 기억인가?
\item 시 속의 어린 아이는 누구를 기다리고 있는가?
\item 엄마는 무엇을 하러 어디에 갔는가?
\item 아이가 혼자 방에 남겨져 느끼는 감정은 무엇인가?
\item 시간이 흐르고 날이 어두워지면서 아이가 느끼는 주된 감정은 무엇인가?
\item '아주 먼 옛날'의 기억이 현재의 화자에게 어떤 영향을 미치고 있는가?
\item 이 시는 어떤 표현을 통해 감정을 생생하게 전달하는가?
\item ‘해는 시든 지 오래’라는 표현은 어떤 감각적 이미지에 해당하는가?
\item 해가 졌다는 표현은 아이의 어떤 감정을 고조시키는가?
\item 아이가 간절히 기다리는 소리는 무엇인가?
\item '금간 창틈으로 고요히 빗소리'에 나타나는 감각적 이미지는 무엇인가?
\item ‘찬밥’과 ‘윗목’은 공통적으로 어떤 감각적 이미지를 느끼게 하는가?
\item 하나의 감각을 다른 감각으로 바꾸어 표현하는 기법을 무엇이라고 하는가?
\item 이 시에서 화자의 처지를 빗대어 표현한 사물은 무엇인가?
\item ‘찬밥처럼 방에 담겨’ 있다는 것은 어떤 마음을 표현한 것인가?
\item 엄마의 지친 발걸음을 무엇에 비유했는가?
\item ‘배춧잎 같은 발소리’라는 표현을 통해 짐작할 수 있는 엄마의 상태는 어떠한가?
\item 화자의 어린 시절 전체를 상징하는 시어는 무엇인가?
\item ‘윗목’은 어떤 특징을 가진 공간인가?
\item ‘내 유년의 윗목’이라는 표현은 화자의 어린 시절이 어떠했음을 상징하는가?
\item 'A처럼'과 같이 표현하는 비유법을 무엇이라고 하는가?
\item '내 유년은 윗목이다'와 같이 표현하는 비유법을 무엇이라고 하는가?
\end{enumerate}
\kfEndMC

\kfSecHeader{kfAreaLiterature}{문제}{}

\par\medskip
\kfBeginMC
\begin{kfQBlock}{01}{객관식}
\kfQStem{이 시에 대한 설명으로 적절하지 \uline{\underline{않은}} 것은?}
\begin{kfChoices}
\kfChoice 어른이 된 화자가 어린 시절을 회상하는 형식이다.
\kfChoice 시장에 간 엄마를 홀로 기다리던 외로운 기억을 노래한다.
\kfChoice ‘지금도 내 눈시울을 뜨겁게 하는’ 표현으로 현재까지 이어지는 슬픔을 보여 준다.
\kfChoice 어린 시절의 아픈 기억과 그 정서가 시 전체를 관통한다.
\kfChoice 미래에 대한 희망적 전망을 밝고 활기찬 분위기로 그려 낸다.
\end{kfChoices}
\end{kfQBlock}
\begin{kfQBlock}{02}{객관식}
\kfQStem{시적 화자에 대한 설명으로 적절하지 \uline{\underline{않은}} 것은?}
\begin{kfChoices}
\kfChoice 어린 시절, 홀로 방에 남겨진 경험이 있다.
\kfChoice 시장에 간 어머니를 초조하게 기다리고 있다.
\kfChoice 과거의 기억 때문에 현재에도 슬픔을 느낀다.
\kfChoice 빈방에서 홀로 외로움과 두려움을 함께 느끼고 있다.
\kfChoice 시간이 흘러 이제는 과거의 기억을 잊고 지내고 있다.
\end{kfChoices}
\end{kfQBlock}
\begin{kfQBlock}{03}{객관식}
\kfQStem{이 시의 표현 방식에 대한 설명으로 적절하지 \uline{\underline{않은}} 것은?}
\begin{kfChoices}
\kfChoice 화자의 정서를 직접 제시하고 있다.
\kfChoice 밝고 어두운 색채의 대비를 통해 희망을 강조하고 있다.
\kfChoice 독백적인 어조를 사용하여 화자의 내면을 진솔하게 고백한다.
\kfChoice 감각적인 이미지를 활용하여 시의 분위기를 생생하게 만든다.
\kfChoice 비유적인 표현을 통해 화자의 정서를 효과적으로 드러낸다.
\end{kfChoices}
\end{kfQBlock}
\begin{kfQBlock}{04}{객관식}
\kfQStem{㉠\textasciitilde{}㉤ 중, 늦게까지 일하고 돌아온 어머니의 지친 모습을 비유적으로 표현한 구절로 가장 적절한 것은?}
\begin{kfChoices}
\kfChoice ㉠
\kfChoice ㉡
\kfChoice ㉢
\kfChoice ㉣
\kfChoice ㉤
\end{kfChoices}
\end{kfQBlock}
\begin{kfQBlock}{05}{객관식}
\kfQStem{<보기>를 참고하여 이 시를 감상한 내용으로 적절하지 \uline{\underline{않은}} 것은?}
\begin{kfEvidence}시에서 ‘윗목’은 아궁이에서 불을 땔 때 온기가 잘 전달되지 않는 방의 찬 부분을 가리킨다. 문학에서는 이러한 공간의 특성을 활용하여, 따뜻한 사랑이나 관심이 부족했던 서럽고 힘든 상태나 시절을 상징하는 의미로 사용되고는 한다.\end{kfEvidence}
\begin{kfChoices}
\kfChoice '찬밥'이라는 시어는 '윗목'이 주는 차가운 이미지와 연결되는군.
\kfChoice '윗목'과 반대되는 '아랫목'은 따뜻한 사랑을 의미한다고 볼 수 있군.
\kfChoice '윗목'은 화자가 느꼈던 어린 시절의 가난과 서러움을 함축하고 있군.
\kfChoice 화자는 '윗목' 같았던 유년 시절을 이제는 따뜻한 추억으로 여기고 있군.
\kfChoice 결국 이 시는 '윗목'으로 상징되는 유년의 아픔에 대한 고백이라고 할 수 있군.
\end{kfChoices}
\end{kfQBlock}
\begin{kfQBlock}{06}{객관식}
\kfQStem{이 시의 화자를 ‘나’에서 ‘엄마’로 바꾸어 다시 쓸 때 나타나는 효과로 적절하지 \uline{\underline{않은}} 것은?}
\begin{kfChoices}
\kfChoice 시장에서 일하는 엄마의 힘겹고 고된 상황이 한층 잘 드러난다.
\kfChoice 홀로 집에 있을 아이를 걱정하는 엄마의 마음이 중심이 된다.
\kfChoice 시선이 아이의 외로움에서 엄마의 고단함으로 자연스럽게 옮겨 간다.
\kfChoice 아이의 속마음은 엄마의 시선을 통해 간접적으로 추측되게 된다.
\kfChoice 한 인물의 시점에서 벗어나 사건을 객관적으로 전달하는 효과가 생긴다.
\end{kfChoices}
\end{kfQBlock}
\begin{kfQBlock}{07}{객관식}
\kfQStem{다음은 시를 감상한 학생들의 대화이다. 내용이 적절하지 \uline{\underline{않은}} 학생은?}
\begin{kfEvidence}민수: ‘해는 시든 지 오래’라는 표현에서 시간이 많이 흘렀다는 것과 함께 어두운 분위기를 시각적으로 느낄 수 있었어. \\\relax
지혜: 화자를 ‘찬밥’에 비유하고 시의 마지막을 ‘윗목’으로 끝낸 것에서 화자가 느꼈을 차가운 외로움이 잘 전해졌어. \\\relax
준호: ‘타박타박’이나 ‘빗소리’ 같은 소리를 나타내는 말을 사용해서, 혼자 있는 빈방의 적막함이 더 생생하게 느껴졌어. \\\relax
서연: ‘열무’나 ‘찬밥’ 같은 음식을 소재로 사용해서 배고팠던 어린 시절의 서러움을 미각적으로 표현하고 있어. \\\relax
현우: ‘아주 먼 옛날’과 ‘지금도’라는 말을 통해, 과거의 아픈 기억이 현재까지 이어지고 있다는 것을 알 수 있었어.\end{kfEvidence}
\begin{kfChoices}
\kfChoice 민수
\kfChoice 지혜
\kfChoice 준호
\kfChoice 서연
\kfChoice 현우
\end{kfChoices}
\end{kfQBlock}
\begin{kfQBlock}{08}{서술형}
\kfQStem{"나는 찬밥처럼 방에 담겨"라는 표현에 담긴 화자의 정서와 처지를 <조건>에 맞게 서술하시오.}
\kfAnswerNote{답안}{4}
\end{kfQBlock}
\begin{kfQBlock}{09}{서술형}
\kfQStem{어른이 된 화자가 여전히 어린 시절의 기억을 떠올리며 눈물을 흘리는 이유는 무엇일지, 시의 내용을 바탕으로 <조건>에 맞게 서술하시오.}
\kfAnswerNote{답안}{4}
\end{kfQBlock}
\kfEndMC


\kfStartNonlit{[과학] 중력을 이기는 식물의 물 운반 비밀}


\kfPassageNote{[과학] 중력을 이기는 식물의 물 운반 비밀}{%
　식물이 살아가려면 물이 꼭 필요하다. 식물은 잎에서 광합성을 통해 양분을 만드는데, 물이 바로 그 재료가 되기 때문이다. 그런데 물은 중력 때문에 높은 곳에서 낮은 곳으로 흐르는데, 식물은 반대로 위쪽으로 자란다. 그렇다면 식물은 어떻게 뿌리에서 흡수한 물을 높은 곳에 있는 잎까지 보낼 수 있을까?

　미국에는 키가 112m나 되는 거대한 세쿼이아 나무가 있다. 뿌리가 땅속 15m까지 뻗어 있어서 물이 뿌리에서 나무 꼭대기까지 가려면 무려 127m를 올라가야 한다. 펌프 같은 장치도 없는데 어떻게 이런 일이 가능할까? 식물은 뿌리압, 모세관 현상, 증산 작용이라는 세 가지 힘을 함께 사용해서 물을 끌어 올린다.

　뿌리압은 뿌리에서 물을 위로 밀어 올리는 힘이다. 뿌리 안쪽은 여러 가지 물질들이 녹아 있어서 흙 속 물보다 농도가 높다. 이때 농도 균형을 맞추려고 흙 속의 물이 뿌리 안으로 들어오면서 압력이 생긴다. 호박이나 수세미의 줄기를 자르고 뿌리를 물에 넣어보면 잘린 부분에서 물이 계속 올라오는 것을 볼 수 있는데, 이것이 바로 뿌리압 때문이다.

　모세관 현상은 가는 관을 따라 액체가 올라가는 현상이다. 물 분자와 관 벽이 달라붙으려는 힘이 물 분자끼리 붙으려는 힘보다 크기 때문에 일어난다. 관이 가늘수록 물이 더 높이 올라간다. 식물 안에는 뿌리에서 잎까지 연결된 물관이라는 통로가 있는데, 그 지름이 0.075mm로 매우 가늘다. 이렇게 가는 물관 덕분에 모세관 현상으로 물을 밀어 올릴 수 있다.

　증산 작용은 식물이 물을 끌어 올리는 가장 큰 힘이다. 식물의 잎에는 기공이라는 작은 구멍이 있어서 여기서 물이 수증기로 증발한다. 나무 그늘이 건물 그늘보다 시원한 이유도 잎에서 물이 증발하면서 주위 열을 흡수하기 때문이다. 10cm×10cm 크기의 잔디밭에서 1년 동안 증발하는 물의 양은 무려 55kg이나 된다.

　증산 작용이 일어나면 물 분자들이 사슬처럼 연결된 물 기둥을 만든다. 잎의 기공을 통해 맨 위쪽 물 분자가 빠져나가면 아래쪽 물 분자들이 연달아 끌려 올라온다. 이렇게 잡아당기는 힘이 식물이 물을 높은 곳까지 끌어 올리는 가장 중요한 힘이다.

　이처럼 식물은 뿌리압, 모세관 현상, 증산 작용이라는 세 가지 힘을 함께 사용해서 중력을 이기고 물을 높은 곳까지 끌어 올릴 수 있다.
}


\kfPassageNote{문장 독해 — 본문 문장}{%
1. 식물이 살아가려면 물이 꼭 필요하다.

2. 식물은 잎에서 광합성을 통해 양분을 만드는데, 물이 바로 그 재료가 되기 때문이다.

3. 그런데 물은 중력 때문에 높은 곳에서 낮은 곳으로 흐르는데, 식물은 반대로 위쪽으로 자란다.

4. 그렇다면 식물은 어떻게 뿌리에서 흡수한 물을 높은 곳에 있는 잎까지 보낼 수 있을까?

5. 미국에는 키가 112m나 되는 거대한 세쿼이아 나무가 있다.

6. 뿌리가 땅속 15m까지 뻗어 있어서 물이 뿌리에서 나무 꼭대기까지 가려면 무려 127m를 올라가야 한다.

7. 펌프 같은 장치도 없는데 어떻게 이런 일이 가능할까?

8. 식물은 뿌리압, 모세관 현상, 증산 작용이라는 세 가지 힘을 함께 사용해서 물을 끌어 올린다.

9. 뿌리압은 뿌리에서 물을 위로 밀어 올리는 힘이다.

10. 뿌리 안쪽은 여러 가지 물질들이 녹아 있어서 흙 속 물보다 농도가 높다.

11. 이때 농도 균형을 맞추려고 흙 속의 물이 뿌리 안으로 들어오면서 압력이 생긴다.

12. 호박이나 수세미의 줄기를 자르고 뿌리를 물에 넣어보면 잘린 부분에서 물이 계속 올라오는 것을 볼 수 있는데, 이것이 바로 뿌리압 때문이다.

13. 모세관 현상은 가는 관을 따라 액체가 올라가는 현상이다.

14. 물 분자와 관 벽이 달라붙으려는 힘이 물 분자끼리 붙으려는 힘보다 크기 때문에 일어난다.

15. 관이 가늘수록 물이 더 높이 올라간다.

16. 식물 안에는 뿌리에서 잎까지 연결된 물관이라는 통로가 있는데, 그 지름이 0.075mm로 매우 가늘다.

17. 이렇게 가는 물관 덕분에 모세관 현상으로 물을 밀어 올릴 수 있다.

18. 증산 작용은 식물이 물을 끌어 올리는 가장 큰 힘이다.

19. 식물의 잎에는 기공이라는 작은 구멍이 있어서 여기서 물이 수증기로 증발한다.

20. 나무 그늘이 건물 그늘보다 시원한 이유도 잎에서 물이 증발하면서 주위 열을 흡수하기 때문이다.

21. 10cm×10cm 크기의 잔디밭에서 1년 동안 증발하는 물의 양은 무려 55kg이나 된다.

22. 증산 작용이 일어나면 물 분자들이 사슬처럼 연결된 물 기둥을 만든다.

23. 잎의 기공을 통해 맨 위쪽 물 분자가 빠져나가면 아래쪽 물 분자들이 연달아 끌려 올라온다.

24. 이렇게 잡아당기는 힘이 식물이 물을 높은 곳까지 끌어 올리는 가장 중요한 힘이다.

25. 이처럼 식물은 뿌리압, 모세관 현상, 증산 작용이라는 세 가지 힘을 함께 사용해서 중력을 이기고 물을 높은 곳까지 끌어 올릴 수 있다.
}


\kfActivityBg \par\kfMaybeNeedspace{32\baselineskip}\noindent{\kfTipMark\,\,}{\small\color{kfInk} 다음 표의 '의미'를 읽고, 그에 해당하는 '어휘'를 지문에서 찾아 적으시오.}\par\nopagebreak[4]\smallskip\nopagebreak[4]
\begin{kfVocab}
\kfVocabItem{1}{식물이 빛을 이용하여 양분을 스스로 만드는 과정}
\kfVocabItem{2}{지구가 물체를 끌어당기는 힘}
\kfVocabItem{3}{뿌리에서 물을 위로 밀어 올리는 힘}
\kfVocabItem{4}{아주 가는 관을 따라 액체가 저절로 올라가는 현상}
\kfVocabItem{5}{식물 잎의 기공을 통해 물이 수증기 형태로 빠져나가는 현상}
\kfVocabItem{6}{용액에 특정 성분이 얼마나 포함되어 있는지 나타내는 정도}
\kfVocabItem{7}{일정한 넓이의 면에 수직으로 작용하는 힘}
\kfVocabItem{8}{물질의 고유한 성질을 가지는 가장 작은 입자}
\kfVocabItem{9}{식물 뿌리에서 줄기, 잎까지 연결되어 물을 운반하는 통로}
\kfVocabItem{10}{식물 잎에 있는, 기체가 드나드는 작은 구멍}
\kfVocabItem{11}{물이 기체 상태로 변한 것}
\kfVocabItem{12}{액체가 표면에서 기체로 변하는 현상}
\end{kfVocab}

\kfActivityBg \par\needspace{15\baselineskip}
\kfSecHeader{kfAreaNonlit}{활동}{문장 독해}
\par\kfMaybeNeedspace{32\baselineskip}\noindent{\kfTipMark\,\,}{\small\color{kfInk} 다음은 본문의 번호 매긴 문장입니다. 각 문장을 읽고 물음에 답하시오.}\par\nopagebreak[4]\smallskip\nopagebreak[4]
\kfBeginMC
\begin{kfSentenceUnit}{1}{}
\kfSentQ{1.}{식물이 살아가기 위해 꼭 필요한 것은 무엇인가?}
\begin{kfChoices}
\kfChoice 흙
\kfChoice 빛
\kfChoice 물
\kfChoice 공기
\end{kfChoices}
\end{kfSentenceUnit}

\begin{kfSentenceUnit}{2}{}
\kfSentQ{2.}{광합성의 재료가 되는 것은 무엇인가?}
\begin{kfChoices}
\kfChoice 물
\kfChoice 산소
\kfChoice 양분
\kfChoice 포도당
\end{kfChoices}
\end{kfSentenceUnit}

\begin{kfSentenceUnit}{3}{}
\kfSentQ{3.}{물이 높은 곳에서 낮은 곳으로 흐르는 이유는 무엇인가?}
\begin{kfChoices}
\kfChoice 압력 때문이다
\kfChoice 증발 때문이다
\kfChoice 농도 때문이다
\kfChoice 중력 때문이다
\end{kfChoices}
\end{kfSentenceUnit}

\begin{kfSentenceUnit}{4}{}
\kfSentQ{4.}{물과 반대로, 식물은 어느 방향으로 자라는가?}
\begin{kfChoices}
\kfChoice 옆쪽으로 자란다
\kfChoice 땅속으로 자란다
\kfChoice 아래쪽으로 자란다
\kfChoice 위쪽으로 자란다
\end{kfChoices}
\end{kfSentenceUnit}

\begin{kfSentenceUnit}{5}{}
\kfSentQ{5.}{이 문장에서 제기하고 있는 핵심 의문은 무엇인가?}
\begin{kfChoices}
\kfChoice 식물은 왜 물이 필요한가?
\kfChoice 물은 왜 낮은 곳으로 흐르는가?
\kfChoice 식물은 어떻게 물을 거꾸로 끌어올리는가?
\kfChoice 뿌리의 역할은 무엇인가?
\end{kfChoices}
\end{kfSentenceUnit}

\begin{kfSentenceUnit}{6}{}
\kfSentQ{6.}{이 글에서 예시로 든 나무의 이름과 키는 얼마인가?}
\begin{kfChoices}
\kfChoice 소나무, 15m
\kfChoice 참나무, 55m
\kfChoice 대나무, 80m
\kfChoice 세쿼이아 나무, 112m
\end{kfChoices}
\end{kfSentenceUnit}

\begin{kfSentenceUnit}{7}{}
\kfSentQ{7.}{물이 이동해야 하는 총 거리(127m)는 나무의 키와 뿌리 깊이를 어떻게 계산한 값인가?}
\begin{kfChoices}
\kfChoice 나무의 키(112m) × 2
\kfChoice 뿌리 깊이(15m) × 10
\kfChoice 나무의 키(112m) - 뿌리 깊이(15m)
\kfChoice 나무의 키(112m) + 뿌리 깊이(15m)
\end{kfChoices}
\end{kfSentenceUnit}

\begin{kfSentenceUnit}{8}{}
\kfSentQ{8.}{이 문장에서 '이런 일'이 구체적으로 가리키는 내용은 무엇인가?}
\begin{kfChoices}
\kfChoice 잎에서 양분을 만들어 저장하는 일
\kfChoice 뿌리가 땅속 깊이 뻗어 나가는 일
\kfChoice 중력을 이용하여 물을 내리는 일
\kfChoice 펌프 없이 127m 높이로 물을 끌어 올리는 일
\end{kfChoices}
\end{kfSentenceUnit}

\begin{kfSentenceUnit}{9}{}
\kfSentQ{9.}{식물이 물을 끌어올리기 위해 사용하는 세 가지 힘은 무엇인가?}
\begin{kfChoices}
\kfChoice 중력, 마찰력, 탄성력
\kfChoice 광합성, 호흡 작용, 증산 작용
\kfChoice 응집력, 부착력, 표면 장력
\kfChoice 뿌리압, 모세관 현상, 증산 작용
\end{kfChoices}
\end{kfSentenceUnit}

\begin{kfSentenceUnit}{10}{}
\kfSentQ{10.}{'뿌리압'이란 무엇인가?}
\begin{kfChoices}
\kfChoice 잎에서 물을 당기는 힘
\kfChoice 물을 뿌리 쪽에 가두는 힘
\kfChoice 뿌리에서 물을 위로 밀어 올리는 힘
\kfChoice 줄기를 튼튼하게 지탱하는 힘
\end{kfChoices}
\end{kfSentenceUnit}

\begin{kfSentenceUnit}{11}{}
\kfSentQ{11.}{뿌리 안쪽과 흙 속 물의 농도는 어느 쪽이 더 높은가?}
\begin{kfChoices}
\kfChoice 흙 속
\kfChoice 뿌리 안쪽
\kfChoice 둘이 같다
\kfChoice 알 수 없다
\end{kfChoices}
\end{kfSentenceUnit}

\begin{kfSentenceUnit}{12}{}
\kfSentQ{12.}{흙 속의 물이 뿌리 안으로 들어오는 이유는 무엇인가?}
\begin{kfChoices}
\kfChoice 뿌리압을 없애려고
\kfChoice 중력의 영향을 받아서
\kfChoice 광합성을 하기 위해서
\kfChoice 농도 균형을 맞추려고
\end{kfChoices}
\end{kfSentenceUnit}

\begin{kfSentenceUnit}{13}{}
\kfSentQ{13.}{'이것'이 가리키는 현상은 무엇인가?}
\begin{kfChoices}
\kfChoice 잘린 줄기 부분에서 물이 올라오는 현상
\kfChoice 뿌리가 썩어서 물러지는 현상
\kfChoice 잎이 노랗게 변하는 현상
\kfChoice 줄기가 말라 비틀어지는 현상
\end{kfChoices}
\end{kfSentenceUnit}

\begin{kfSentenceUnit}{14}{}
\kfSentQ{14.}{'모세관 현상'이란 무엇인가?}
\begin{kfChoices}
\kfChoice 굵은 관을 따라 액체가 내려가는 현상
\kfChoice 가는 관을 따라 액체가 올라가는 현상
\kfChoice 액체가 기체로 변하여 날아가는 현상
\kfChoice 고체가 액체 위로 떠오르는 현상
\end{kfChoices}
\end{kfSentenceUnit}

\begin{kfSentenceUnit}{15}{}
\kfSentQ{15.}{모세관 현상이 일어나는 이유는 무엇인가?}
\begin{kfChoices}
\kfChoice 중력이 작용하기 때문에
\kfChoice 기압 차이가 발생해서
\kfChoice 물 분자끼리 붙는 힘이 더 커서
\kfChoice 물 분자와 관 벽이 붙는 힘이 더 커서
\end{kfChoices}
\end{kfSentenceUnit}

\begin{kfSentenceUnit}{16}{}
\kfSentQ{16.}{관의 굵기와 물이 올라가는 높이는 어떤 관계인가?}
\begin{kfChoices}
\kfChoice 비례 관계
\kfChoice 관계 없다
\kfChoice 반비례 관계
\kfChoice 일정하다
\end{kfChoices}
\end{kfSentenceUnit}

\begin{kfSentenceUnit}{17}{}
\kfSentQ{17.}{식물 속에서 물이 이동하는 통로의 이름은 무엇인가?}
\begin{kfChoices}
\kfChoice 물관
\kfChoice 체관
\kfChoice 기공
\kfChoice 엽록체
\end{kfChoices}
\end{kfSentenceUnit}

\begin{kfSentenceUnit}{18}{}
\kfSentQ{18.}{물관의 지름은 얼마인가?}
\begin{kfChoices}
\kfChoice 0.005mm
\kfChoice 0.075mm
\kfChoice 0.75mm
\kfChoice 7.5mm
\end{kfChoices}
\end{kfSentenceUnit}

\begin{kfSentenceUnit}{19}{}
\kfSentQ{19.}{식물에서 모세관 현상이 일어날 수 있는 이유는 무엇인가?}
\begin{kfChoices}
\kfChoice 물관이 매우 넓기 때문
\kfChoice 물관이 막혀 있기 때문
\kfChoice 물관이 매우 가늘기 때문
\kfChoice 잎이 넓기 때문
\end{kfChoices}
\end{kfSentenceUnit}

\begin{kfSentenceUnit}{20}{}
\kfSentQ{20.}{식물이 물을 끌어올리는 세 가지 힘 중 가장 큰 힘은 무엇인가?}
\begin{kfChoices}
\kfChoice 중력
\kfChoice 뿌리압
\kfChoice 모세관 현상
\kfChoice 증산 작용
\end{kfChoices}
\end{kfSentenceUnit}

\begin{kfSentenceUnit}{21}{}
\kfSentQ{21.}{증산 작용은 잎의 어디에서 일어나는가?}
\begin{kfChoices}
\kfChoice 기공
\kfChoice 엽록체
\kfChoice 세포벽
\kfChoice 뿌리털
\end{kfChoices}
\end{kfSentenceUnit}

\begin{kfSentenceUnit}{22}{}
\kfSentQ{22.}{나무 그늘이 더 시원한 이유는 무엇인가?}
\begin{kfChoices}
\kfChoice 잎이 햇빛을 가리기 때문
\kfChoice 광합성으로 열을 방출하기 때문
\kfChoice 물이 증발하며 열을 흡수하기 때문
\kfChoice 뿌리에서 찬물을 올리기 때문
\end{kfChoices}
\end{kfSentenceUnit}

\begin{kfSentenceUnit}{23}{}
\kfSentQ{23.}{이 문장은 증산 작용의 어떤 점을 보여주기 위한 예시인가?}
\begin{kfChoices}
\kfChoice 잔디가 잘 자란다는 점
\kfChoice 물이 적게 든다는 점
\kfChoice 많은 양의 물이 증발한다는 점
\kfChoice 증발 속도가 느리다는 점
\end{kfChoices}
\end{kfSentenceUnit}

\begin{kfSentenceUnit}{24}{}
\kfSentQ{24.}{증산 작용이 일어날 때 물 분자들은 어떤 형태를 만드는가?}
\begin{kfChoices}
\kfChoice 흩어진 물방울
\kfChoice 얼음 결정
\kfChoice 수증기 구름
\kfChoice 사슬처럼 연결된 물 기둥
\end{kfChoices}
\end{kfSentenceUnit}

\begin{kfSentenceUnit}{25}{}
\kfSentQ{25.}{물기둥이 위로 끌려 올라오는 직접적인 원인은 무엇인가?}
\begin{kfChoices}
\kfChoice 맨 위쪽 물 분자가 빠져나가는 것
\kfChoice 아래쪽에서 밀어 올리는 것
\kfChoice 중간에서 끊어지는 것
\kfChoice 중력이 당기는 것
\end{kfChoices}
\end{kfSentenceUnit}

\begin{kfSentenceUnit}{26}{}
\kfSentQ{26.}{증산 작용이 물을 끌어올리는 가장 중요한 힘인 이유는 무엇 때문인가?}
\begin{kfChoices}
\kfChoice '미는 힘'으로 물을 올리기 때문
\kfChoice '잡아당기는 힘'으로 물기둥 전체를 끌어올리기 때문
\kfChoice 물을 저장하는 힘이므로
\kfChoice 뿌리를 튼튼하게 하기 때문
\end{kfChoices}
\end{kfSentenceUnit}

\begin{kfSentenceUnit}{27}{}
\kfSentQ{27.}{이 문장이 식물이 물을 끌어올리기 위해 사용한다고 요약한 세 가지 힘은 무엇인가?}
\begin{kfChoices}
\kfChoice 중력, 마찰력, 탄성력
\kfChoice 응집력, 부착력, 장력
\kfChoice 뿌리압, 모세관 현상, 증산 작용
\kfChoice 광합성, 호흡, 소화
\end{kfChoices}
\end{kfSentenceUnit}

\begin{kfSentenceUnit}{28}{}
\kfSentQ{28.}{식물이 물을 위로 끌어올리기 위해 이겨내야 하는 힘은 무엇인가?}
\begin{kfChoices}
\kfChoice 마찰력
\kfChoice 탄성력
\kfChoice 자기력
\kfChoice 중력
\end{kfChoices}
\end{kfSentenceUnit}

\kfEndMC

\kfActivityBg \par\kfMaybeNeedspace{32\baselineskip}\noindent{\kfTipMark\,\,}{\small\color{kfInk} 글의 전체적인 논리 흐름을 파악하며 빈칸에 알맞은 내용을 채워 보세요.}\par\nopagebreak[4]\smallskip\nopagebreak[4]
\begin{kfStructure}
\kfNode{0}{}{}
\end{kfStructure}

\kfSecHeader{kfAreaNonlit}{문제}{}

\par\medskip
\kfBeginMC
\begin{kfQBlock}{01}{객관식}
\kfQStem{윗글의 내용과 일치하지 \uline{않는} 것은?}
\begin{kfChoices}
\kfChoice 식물은 중력과 반대 방향으로 물을 이동시켜야 한다.
\kfChoice 뿌리압은 농도 차이 때문에 발생하는 힘이다.
\kfChoice 모세관 현상은 관이 굵을수록 더 강하게 일어난다.
\kfChoice 나무 그늘이 시원한 것은 증산 작용과 관련이 있다.
\kfChoice 증산 작용은 물을 밀어 올리기보다 끌어당기는 원리에 가깝다.
\end{kfChoices}
\end{kfQBlock}
\begin{kfQBlock}{02}{객관식}
\kfQStem{윗글에 나타난 세 가지 힘에 대한 설명으로 적절하지 \uline{\underline{않은}} 것은?}
\begin{kfChoices}
\kfChoice 식물은 뿌리압·모세관 현상·증산 작용을 함께 사용해 물을 끌어 올린다.
\kfChoice 뿌리압은 뿌리에서 물을 위로 밀어 올리는 힘이다.
\kfChoice 모세관 현상은 물과 관 벽이 붙으려는 힘으로 물을 끌어 올린다.
\kfChoice 증산 작용은 잎의 기공에서 물이 증발하며 물을 잡아당기는 힘이다.
\kfChoice 세 가지 힘은 식물 안에서 각각 독립적으로만 작용한다.
\end{kfChoices}
\end{kfQBlock}
\begin{kfQBlock}{03}{객관식}
\kfQStem{<보기>의 현상과 가장 관련 깊은 원리는?}
\begin{kfEvidence}컵에 물을 담고 휴지 한쪽 끝을 담가 두었더니, 잠시 후 물이 휴지를 타고 올라와 젖지 않았던 부분까지 번져 나갔다.\end{kfEvidence}
\begin{kfChoices}
\kfChoice 광합성
\kfChoice 뿌리압
\kfChoice 모세관 현상
\kfChoice 증산 작용
\kfChoice 중력
\end{kfChoices}
\end{kfQBlock}
\begin{kfQBlock}{04}{객관식}
\kfQStem{윗글을 바탕으로 추론한 내용으로 적절하지 \uline{\underline{않은}} 것은?}
\begin{kfChoices}
\kfChoice 잎이 없는 겨울철에는 증산 작용이 활발하지 않을 것이다.
\kfChoice 사막처럼 건조한 지역의 식물은 기공의 수가 적을 수 있다.
\kfChoice 습도가 높은 날에는 증산 작용이 더 활발하게 일어날 것이다.
\kfChoice 비가 많이 와서 흙의 습도가 매우 높은 날에는 뿌리압이 강하게 작용할 것이다.
\kfChoice 키가 큰 나무일수록 물을 꼭대기까지 올리는 데 더 많은 힘이 필요할 것이다.
\end{kfChoices}
\end{kfQBlock}
\begin{kfQBlock}{05}{객관식}
\kfQStem{식물이 물을 끌어 올리는 힘을 작용 방식에 따라 분류한 설명으로 적절하지 \uline{\underline{않은}} 것은?}
\begin{kfChoices}
\kfChoice 뿌리압은 물을 위로 ‘밀어 올리는 힘’으로 분류된다.
\kfChoice 증산 작용은 잎에서 물을 위로 ‘잡아당기는 힘’으로 분류된다.
\kfChoice 모세관 현상은 가는 관을 따라 물을 끌어 올리는 힘이다.
\kfChoice ‘당기는 힘’에 해당하는 것은 증산 작용이다.
\kfChoice 증산 작용은 물을 위로 ‘밀어 올리는 힘’에 해당한다.
\end{kfChoices}
\end{kfQBlock}
\begin{kfQBlock}{06}{단답형}
\kfQStem{식물이 잎에서 양분을 만드는 과정에서 물을 재료로 사용하는 활동은 무엇인가?}
\kfAnswerLines{1}
\end{kfQBlock}
\begin{kfQBlock}{07}{단답형}
\kfQStem{식물이 물을 끌어 올리기 위해 사용하는 세 가지 힘은 무엇인가?}
\kfAnswerLines{1}
\end{kfQBlock}
\begin{kfQBlock}{08}{단답형}
\kfQStem{흙 속의 물이 뿌리 안으로 들어오면서 생기는 압력으로 물을 위로 밀어 올리는 힘은 무엇인가?}
\kfAnswerLines{1}
\end{kfQBlock}
\begin{kfQBlock}{09}{단답형}
\kfQStem{뿌리 안쪽과 흙 속 물의 무엇이 다르기 때문에 뿌리압이 생기는가?}
\kfAnswerLines{1}
\end{kfQBlock}
\begin{kfQBlock}{10}{단답형}
\kfQStem{식물 안에서 뿌리부터 잎까지 물이 이동하는 매우 가는 통로의 이름은 무엇인가?}
\kfAnswerLines{1}
\end{kfQBlock}
\begin{kfQBlock}{11}{서술형}
\kfQStem{식물이 물을 끌어 올리는 가장 중요한 힘인 '증산 작용'의 과정을 <조건>에 맞게 서술하시오.}
\kfAnswerNote{답안}{4}
\end{kfQBlock}
\kfEndMC



\clearpage

\kfChapterCover{4}{러셀2 (챕터4)}{Chapter 4}{이번 챕터의 학습 내용을 시작합니다.}
\begin{kfChapterAreaList}
\kfChapterAreaItem{문법}{kfAreaGrammar}{단어의 구조를 파헤치다: 형태소, 어근, 접사}
\kfChapterAreaItem{개념}{kfAreaConcept}{소설 구성의 3요소: 인물, 사건, 배경}
\kfChapterAreaItem{문학}{kfAreaLiterature}{「꺼삐딴 리」 — 전광용}
\kfChapterAreaItem{비문학}{kfAreaNonlit}{[기술] 말소리를 글자로 바꾸는 컴퓨터의 비밀}
\end{kfChapterAreaList}

\kfStartGrammar{단어의 구조를 파헤치다: 형태소, 어근, 접사}


\kfPassageNoteNT{%
　우리가 사용하는 단어는 더 작은 단위로 쪼갤 수 있다. 예를 들어 ‘책가방’은 ‘책’과 ‘가방’으로 나눌 수 있다. 만약 ‘책’을 ‘ㅊ, ㅐ, ㄱ’처럼 자음과 모음으로 더 쪼개면 ‘글을 읽는 물건’이라는 본래의 뜻은 사라지고 의미 없는 소리의 조각만 남게 된다. 이처럼 뜻을 가진 채로 더는 쪼갤 수 없는 가장 작은 말의 단위를 형태소라고 한다. 형태소는 단어를 구성하는 기본 재료가 된다는 점에서 문법의 원자와 같다.

　형태소는 크게 자립성 여부와 의미의 유형이라는 두 가지 기준으로 나뉜다. 먼저 자립성을 기준으로, 홀로 쓰일 수 있는 자립 형태소와 반드시 다른 말에 기대어야만 쓰이는 의존 형태소로 나뉜다. 명사, 대명사, 수사, 관형사, 부사, 감탄사는 자립 형태소에 속하며, 조사, 어미, 접사와 용언의 어간은 의존 형태소에 속한다. 다음으로 의미를 기준으로, 실질적인 의미를 가진 실질 형태소와 문법적 기능을 하는 형식 형태소로 나눌 수 있다. 대부분의 자립 형태소와 용언의 어간, 합성어의 어근 등은 실질 형태소이며, 형식 형태소에는 조사, 어미, 접사만이 해당한다.

　여기서 학생들이 가장 많이 혼동하는 부분이 바로 용언의 어간이다. ‘먹다’의 ‘먹-’과 같은 용언의 어간은, ‘먹는다’는 구체적인 뜻을 가진 실질 형태소이면서도 홀로 쓰일 수는 없어 반드시 어미와 결합해야 하는 의존 형태소의 성격을 동시에 지니기 때문이다. 이처럼 하나의 형태소가 두 가지 분류 기준에 따라 어떤 위치를 차지하는지 정확히 이해하는 것이 중요하다.

　이러한 형태소 개념은 단어의 구조를 분석하는 핵심 열쇠가 된다. 단어를 구성할 때, 실질적 의미를 나타내는 중심 부분을 어근이라고 한다. ‘덮개’의 ‘덮-’이나 ‘맨손’의 ‘손’이 어근에 해당한다. 반면, 어근에 붙어 특정한 뜻을 더하는 부분을 접사라고 한다. ‘맨손’처럼 어근 앞에 붙으면 접두사, ‘덮개’처럼 어근 뒤에 붙으면 접미사이다.

　어근과 어간은 헷갈리기 쉽지만, 어근은 단어의 ‘형성’을 따지는 개념이고 어간은 용언의 ‘활용’을 따지는 개념이라는 결정적 차이가 있다. 예를 들어 ‘치솟다’의 경우, 의미의 중심인 ‘솟-’이 어근이고, 활용 시 변하지 않는 부분인 ‘치솟-’ 전체가 어간이 된다. 이처럼 단어를 구성하는 재료와 그것이 문장 속에서 활용되는 방식을 구분하여 이해하는 것은 우리말을 더욱 깊이 있게 탐구하는 중요한 바탕이 된다.
\par\medskip\begin{itemize}[leftmargin=18pt, itemsep=2pt, topsep=2pt, label=\textbullet]
\item 자립/의존 — 명사(책)=자립, 용언 어간(먹-)=의존
\item 실질/형식 — 용언 어간(먹-)=실질, 조사·어미·접사=형식
\item 어근 vs 어간 — '치솟다'의 어근='솟-' / 어간='치솟-'
\end{itemize}
}

\kfActivityBg \par\kfMaybeNeedspace{32\baselineskip}\noindent{\kfTipMark\,\,}{\small\color{kfInk} 지문의 핵심 내용을 떠올리며 빈칸에 알맞은 말을 채워 보세요.}\par\nopagebreak[4]\smallskip\nopagebreak[4]

\kfActivityBg \kfSecHeader{kfAreaGrammar}{활동}{분석 훈련}
\par\kfMaybeNeedspace{24\baselineskip}\noindent{\kfTipMark\,\,}{\small\color{kfInk} 다음 문장이나 단어를 읽고, 설명에 맞게 분석하시오.}\par\nopagebreak[4]\smallskip\nopagebreak[4]
\begin{enumerate}[leftmargin=22pt, itemsep=6pt, topsep=4pt, label=\textbf{\arabic*.}]
\item 다음 문장을 형태소 단위로 모두 나누시오.
\begin{kfEvidence}나는 풋사과를 먹었다.\end{kfEvidence}
\item 위 1번 문장의 형태소 중 '풋-'과 '먹-'을 자립성 여부와 의미 유형에 따라 각각 분석하시오.  (자립/의존 형태소, 실질/형식 형태소)
\begin{kfEvidence}나는 풋사과를 먹었다.\end{kfEvidence}
\item 다음 단어를 어근과 접사로 나누고, 접사의 종류(접두사/접미사)를 밝히시오.
\begin{kfEvidence}지우개\end{kfEvidence}
\item 다음 단어를 어간과 어미로 나누시오.
\begin{kfEvidence}아름답고\end{kfEvidence}
\item 다음 단어에서 어근과 어간을 각각 찾아 쓰시오.
\begin{kfEvidence}치솟았다\end{kfEvidence}
\item 다음 문장에서 실질 형태소를 모두 찾으시오.
\begin{kfEvidence}동생이 새 신을 신었다.\end{kfEvidence}
\item 다음 문장에서 형식 형태소를 모두 찾으시오.
\begin{kfEvidence}그가 웃음을 터뜨렸다.\end{kfEvidence}
\item 다음 단어에서 접두사를 찾아 쓰시오.
\begin{kfEvidence}군소리\end{kfEvidence}
\item 다음 단어에서 접미사를 찾아 쓰시오.
\begin{kfEvidence}장난꾸러기\end{kfEvidence}
\item 다음 문장을 형태소로 분석할 때, 그 개수가 가장 많은 단어는 무엇인가?
\begin{kfEvidence}나는 잠꾸러기 동생을 깨웠다.\end{kfEvidence}
\end{enumerate}

\kfSecHeader{kfAreaGrammar}{문제}{}

\par\medskip
\kfBeginMC
\begin{kfQBlock}{01}{객관식}
\kfQStem{다음 중 형태소에 대한 설명으로 적절하지 \uline{\underline{않은}} 것은?}
\begin{kfChoices}
\kfChoice 뜻을 가진 가장 작은 말의 단위이다.
\kfChoice 자립성 여부에 따라 자립 형태소와 의존 형태소로 나뉜다.
\kfChoice 모든 형태소는 홀로 쓰일 수 있다.
\kfChoice 실질적 의미 유무에 따라 실질 형태소와 형식 형태소로 나뉜다.
\kfChoice 단어를 구성하는 기본 재료가 된다.
\end{kfChoices}
\end{kfQBlock}
\begin{kfQBlock}{02}{객관식}
\kfQStem{다음 단어 중 형태소의 개수가 나머지와 \uline{다른} 하나는?}
\begin{kfChoices}
\kfChoice 구름
\kfChoice 손수건
\kfChoice 밤낮
\kfChoice 돌다리
\kfChoice 책가방
\end{kfChoices}
\end{kfQBlock}
\begin{kfQBlock}{03}{객관식}
\kfQStem{밑줄 친 형태소의 종류가 나머지와 \uline{다른} 하나는?}
\begin{kfChoices}
\kfChoice \uline{하늘}이 맑다.
\kfChoice \uline{꽃}이 피었다.
\kfChoice 밥을 \uline{먹}었다.
\kfChoice \uline{새} 책을 샀다.
\kfChoice \uline{강}이 흐른다.
\end{kfChoices}
\end{kfQBlock}
\begin{kfQBlock}{04}{객관식}
\kfQStem{다음 <보기>의 ㉠, ㉡에 대한 설명으로 옳은 것은?}
\begin{kfEvidence}아기가 ㉠놀이터에서 논다. 여기서 '놀이터'는 ㉡'놀-'이라는 중심 의미에 '-이-'와 '-터'가 붙어 만들어졌다.\end{kfEvidence}
\begin{kfChoices}
\kfChoice ㉠의 '놀이-'는 어간이다.
\kfChoice ㉠의 '-터'는 접두사이다.
\kfChoice ㉡은 의존 형태소이면서 형식 형태소이다.
\kfChoice ㉡은 단어의 중심 의미를 나타내는 어근이다.
\kfChoice ㉡은 활용 시 변하는 부분인 어미이다.
\end{kfChoices}
\end{kfQBlock}
\begin{kfQBlock}{05}{객관식}
\kfQStem{다음 밑줄 친 부분 중 문법적 기능을 하는 형식 형태소가 \uline{\underline{아닌}} 것은?}
\begin{kfChoices}
\kfChoice \uline{하늘}이 맑다.
\kfChoice 책을 읽\uline{는}다.
\kfChoice \uline{맨}손으로 잡았다.
\kfChoice 웃\uline{음}소리.
\kfChoice 학교\uline{에} 간다.
\end{kfChoices}
\end{kfQBlock}
\begin{kfQBlock}{06}{객관식}
\kfQStem{밑줄 친 부분이 접두사가 \uline{\underline{아닌}} 것은?}
\begin{kfChoices}
\kfChoice \uline{맨}손 체조
\kfChoice \uline{새}파란 하늘
\kfChoice \uline{군}것질거리
\kfChoice \uline{새} 신발
\kfChoice \uline{헛}걸음을 했다.
\end{kfChoices}
\end{kfQBlock}
\begin{kfQBlock}{07}{객관식}
\kfQStem{‘어근’과 ‘어간’에 대한 설명으로 적절하지 \uline{\underline{않은}} 것은?}
\begin{kfChoices}
\kfChoice 어근은 새로운 단어를 만드는 형성 과정과 관련하여 다루어지는 문법 개념이다.
\kfChoice 어간은 용언이 다양한 어미와 결합해 활용할 때 형태가 변하지 않는 부분을 가리킨다.
\kfChoice 어근은 한 단어 안에서 실질적 의미를 지닌 가장 중심이 되는 핵심 부분이다.
\kfChoice 단어에 따라 어근과 어간이 가리키는 범위가 일치할 수도, 서로 다를 수도 있다.
\kfChoice 어근은 활용할 때 변하고 어간은 변하지 않으며, 두 개념은 동일한 대상을 가리킨다.
\end{kfChoices}
\end{kfQBlock}
\begin{kfQBlock}{08}{객관식}
\kfQStem{다음 중 의존 형태소만으로 이루어진 단어는?}
\begin{kfChoices}
\kfChoice 먹다
\kfChoice 하늘
\kfChoice 손발
\kfChoice 책
\kfChoice 구름
\end{kfChoices}
\end{kfQBlock}
\begin{kfQBlock}{09}{객관식}
\kfQStem{밑줄 친 부분에 대한 분석이 바르지 \uline{\underline{않은}} 것은?}
\begin{kfEvidence}나는 치솟는 해를 보았다.\end{kfEvidence}
\begin{kfChoices}
\kfChoice '치솟는'의 어근은 '솟-'이다.
\kfChoice '치솟는'의 '치-'는 접두사이다.
\kfChoice '치솟는'의 어간은 '치솟-'이다.
\kfChoice '치솟는'의 어미는 '-는'이다.
\kfChoice '치솟는'의 어근과 어간은 같다.
\end{kfChoices}
\end{kfQBlock}
\begin{kfQBlock}{10}{객관식}
\kfQStem{다음 중 실질 형태소의 개수가 가장 많은 것은?}
\begin{kfChoices}
\kfChoice 아이가 잠을 잔다.
\kfChoice 그는 맨손으로 싸웠다.
\kfChoice 우리는 맑은 물을 다 마셨다.
\kfChoice 나는 일곱 개의 사과를 샀다.
\kfChoice 저기 큰 집이 보인다.
\end{kfChoices}
\end{kfQBlock}
\begin{kfQBlock}{11}{객관식}
\kfQStem{다음 중 접미사가 붙어 만들어진 단어가 \uline{\underline{아닌}} 것은?}
\begin{kfChoices}
\kfChoice 구경꾼
\kfChoice 낚시꾼
\kfChoice 장난꾸러기
\kfChoice 미닫이
\kfChoice 가을하늘
\end{kfChoices}
\end{kfQBlock}
\begin{kfQBlock}{12}{객관식}
\kfQStem{다음 <보기>와 같이 형태소 분석이 되는 단어는?}
\begin{kfEvidence}의존 형태소 + 의존 형태소\end{kfEvidence}
\begin{kfChoices}
\kfChoice 하늘
\kfChoice 밥물
\kfChoice 손잡이
\kfChoice 잠
\kfChoice 덥다
\end{kfChoices}
\end{kfQBlock}
\begin{kfQBlock}{13}{객관식}
\kfQStem{다음 중 어간과 어미의 구분이 \uline{불가능한} 단어는?}
\begin{kfChoices}
\kfChoice 읽는다
\kfChoice 슬프다
\kfChoice 갑자기
\kfChoice 가시고
\kfChoice 예뻤다
\end{kfChoices}
\end{kfQBlock}
\begin{kfQBlock}{14}{객관식}
\kfQStem{'실질 형태소이면서 의존 형태소'인 것만으로 짝지어진 것은?}
\begin{kfChoices}
\kfChoice 하늘, 바다
\kfChoice 가-, -고
\kfChoice 먹-, 잡-
\kfChoice 이, 가
\kfChoice 풋-, -개
\end{kfChoices}
\end{kfQBlock}
\begin{kfQBlock}{15}{객관식}
\kfQStem{단어 '짓밟히다'에 대한 분석으로 옳지 \uline{\underline{않은}} 것은?}
\begin{kfChoices}
\kfChoice 어근은 '밟-'이다.
\kfChoice '짓-'은 접두사이다.
\kfChoice '-히-'는 접미사이다.
\kfChoice 어간은 '짓밟히-'이다.
\kfChoice 어근과 어간이 같다.
\end{kfChoices}
\end{kfQBlock}
\begin{kfQBlock}{16}{단답형}
\kfQStem{뜻을 가진 가장 작은 말의 단위를 무엇이라고 하는가?}
\kfAnswerLines{1}
\end{kfQBlock}
\begin{kfQBlock}{17}{단답형}
\kfQStem{형태소 중, '하늘', '구름'처럼 홀로 쓰일 수 있는 형태소는 무엇인가?}
\kfAnswerLines{1}
\end{kfQBlock}
\begin{kfQBlock}{18}{단답형}
\kfQStem{형태소 중, '-다', '어-'처럼 반드시 다른 말과 함께 쓰이는 형태소는 무엇인가?}
\kfAnswerLines{1}
\end{kfQBlock}
\begin{kfQBlock}{19}{단답형}
\kfQStem{실질적인 의미를 나타내는 형태소는 무엇인가?}
\kfAnswerLines{1}
\end{kfQBlock}
\begin{kfQBlock}{20}{단답형}
\kfQStem{문법적인 기능을 하는 조사, 어미, 접사를 통틀어 무엇이라고 하는가?}
\kfAnswerLines{1}
\end{kfQBlock}
\begin{kfQBlock}{21}{단답형}
\kfQStem{단어에서 실질적인 의미를 나타내는 중심 부분을 무엇이라고 하는가?}
\kfAnswerLines{1}
\end{kfQBlock}
\begin{kfQBlock}{22}{단답형}
\kfQStem{어근에 붙어 의미를 더하거나 제한하는 부분을 무엇이라고 하는가?}
\kfAnswerLines{1}
\end{kfQBlock}
\begin{kfQBlock}{23}{단답형}
\kfQStem{'맨손'의 '맨-'처럼 어근 앞에 붙는 접사는 무엇인가?}
\kfAnswerLines{1}
\end{kfQBlock}
\begin{kfQBlock}{24}{단답형}
\kfQStem{'지우개'의 '-개'처럼 어근 뒤에 붙는 접사는 무엇인가?}
\kfAnswerLines{1}
\end{kfQBlock}
\begin{kfQBlock}{25}{단답형}
\kfQStem{용언이 활용할 때 변하지 않는 부분을 무엇이라고 하는가?}
\kfAnswerLines{1}
\end{kfQBlock}
\begin{kfQBlock}{26}{서술형}
\kfQStem{다음 문장을 형태소 단위로 나누고, 각 형태소를 자립성 여부(자립/의존)와 의미 유형(실질/형식)에 따라 분류하시오.}
\begin{kfEvidence}그가 새 책을 읽었다.\end{kfEvidence}
\kfAnswerNote{답안}{4}
\end{kfQBlock}
\kfEndMC


\kfStartConcept{소설 구성의 3요소: 인물, 사건, 배경}


\kfPassageNoteNT{%
　소설은 산문으로 쓰인 서사 문학의 한 갈래로, 인물, 사건, 배경이라는 세 가지 요소로 구성된다. 이 요소들은 서로 긴밀하게 결합하여 하나의 완결된 이야기를 만들어낸다.

　인물은 소설 속에서 행동하고 생각하며 감정을 느끼는 존재이다. 인물은 역할에 따라 주동 인물과 반동 인물로 구분할 수 있는데, 주동 인물은 사건을 이끌어가는 중심인물이고 반동 인물은 주동 인물과 대립하거나 갈등을 일으키는 인물이다. 또한 성격의 변화 여부에 따라 평면적 인물과 입체적 인물로 나뉜다. 평면적 인물은 처음부터 끝까지 성격이 변하지 않는 인물로 단순하고 뚜렷한 성격을 보여주며, 입체적 인물은 사건을 겪으면서 성격이 변화하고 성장하는 인물로 복잡하고 여러 면을 가진 모습을 드러낸다. 인물의 성격은 외모, 말투, 행동, 생각, 다른 인물과의 관계 등을 통해 드러나는데, 작가는 직접적 제시 방법으로 서술자가 인물의 성격을 직접 설명하거나, 간접적 제시 방법으로 인물의 행동이나 대화를 통해 독자 스스로 성격을 파악하도록 한다.

　사건은 소설 속에서 일어나는 여러 사건들이 연결된 줄거리를 말한다. 사건의 전개는 일반적으로 발단, 전개, 위기, 절정, 결말의 5단계 구조를 따른다. 발단에서는 인물과 배경이 소개되고, 전개에서는 사건이 본격적으로 진행되며 갈등이 형성된다. 위기 단계에서는 갈등이 심화되어 긴장감이 높아지고, 절정에서는 갈등이 최고조에 달하며 이야기의 방향이 결정된다. 결말에서는 갈등이 해소되고 이야기가 마무리된다. 이러한 사건들은 원인과 결과의 관계를 가지고 서로 긴밀하게 연결되어야 한다. 갈등은 사건 전개의 중심이 되는 힘인데, 인물과 인물 사이의 외적 갈등뿐만 아니라 인물과 사회의 갈등, 인물과 운명의 갈등, 인물 내면의 심리적 갈등 등 다양한 형태로 나타난다.

　배경은 사건이 일어나는 시간적, 공간적, 사회적 상황을 의미한다. 시간적 배경은 계절, 시대, 특정 연도, 하루 중 특정 시각 등을 포함하며, 공간적 배경은 도시, 시골, 집 안, 학교, 전쟁터 등 사건이 벌어지는 장소를 가리킨다. 사회적 배경은 당시의 시대 상황을 말하며 인물의 사고방식과 가치관에 큰 영향을 미친다. 배경은 단순히 이야기가 펼쳐지는 무대 역할에 그치지 않고 인물의 성격 형성에 영향을 주고, 인물의 심리 상태를 반영하며, 사건 전개의 동기나 원인이 되기도 한다. 또한 작품의 분위기를 만들고 주제를 효과적으로 드러내는 역할을 한다.

　세 요소는 서로 깊게 연결되어 작품의 전체적인 의미를 완성한다. 인물은 배경의 영향을 받아 특정한 성격을 형성하고, 그러한 인물이 행동함으로써 사건이 발생한다. 동시에 사건을 겪으면서 인물은 변화하고 성장한다. 따라서 소설을 읽을 때는 인물, 사건, 배경이라는 요소들이 어떻게 결합하고 조화를 이루는지 전체적으로 분석해야 한다.
\par\medskip\begin{itemize}[leftmargin=18pt, itemsep=2pt, topsep=2pt, label=\textbullet]
\item 인물 — 주동·반동 / 평면·입체 / 직접·간접 제시
\item 사건 — 발단·전개·위기·절정·결말 5단계
\item 배경 — 시간(계절·시대) / 공간(도시·전쟁터) / 사회(시대 상황)
\end{itemize}
}

\kfSecHeader{kfAreaConcept}{문제}{}

\par\medskip
\kfBeginMC
\begin{kfQBlock}{01}{객관식}
\kfQStem{소설 구성의 3요소에 대한 설명으로 적절하지 \uline{\underline{않은}} 것은?}
\begin{kfChoices}
\kfChoice 인물·사건·배경이 소설을 이루는 핵심 요소이다.
\kfChoice 세 요소가 유기적으로 결합하여 작품의 주제를 형성한다.
\kfChoice 인물은 사건을 일으키고 사건은 배경 위에서 전개된다.
\kfChoice 세 요소는 모두 고전 소설과 현대 소설에 공통으로 적용된다.
\kfChoice 세 요소는 서로 영향을 주지 않고 각각 독립적으로만 기능한다.
\end{kfChoices}
\end{kfQBlock}
\begin{kfQBlock}{02}{객관식}
\kfQStem{소설의 ‘인물’에 대한 설명으로 적절하지 \uline{\underline{않은}} 것은 무엇인가?}
\begin{kfChoices}
\kfChoice 작가의 사상이나 가치관을 대변하는 역할을 한다.
\kfChoice 성격이 변하지 않는 평면적 인물과 변하는 입체적 인물이 있다.
\kfChoice 갈등을 유발하고 사건을 전개시키는 원동력이 된다.
\kfChoice 사건이 일어나는 시간적, 공간적 환경을 의미한다.
\kfChoice 독자는 인물에 감정을 이입하며 작품에 몰입하게 된다.
\end{kfChoices}
\end{kfQBlock}
\begin{kfQBlock}{03}{객관식}
\kfQStem{소설의 ‘배경’이 수행하는 기능으로 거리가 \uline{먼} 것은 무엇인가?}
\begin{kfChoices}
\kfChoice 작품의 전반적인 분위기를 조성한다.
\kfChoice 이야기에 사실성과 현장감을 부여한다.
\kfChoice 인물의 성격과 행동의 이유를 설명한다.
\kfChoice 사건 전개의 필연성을 더욱 강화하는 장치가 된다.
\kfChoice 인물 간의 대화를 통해 갈등을 직접적으로 표출한다.
\end{kfChoices}
\end{kfQBlock}
\begin{kfQBlock}{04}{객관식}
\kfQStem{다음 <보기>의 소설에서 밑줄 친 ‘비’가 지니는 배경으로서의 효과는 무엇인가?}
\begin{kfEvidence}새침하게 흐린 품이 눈이 올 듯하더니, 눈은 아니 오고 얼다가 만 비가 추적추적 내리는 날이었다. 이 날이야말로 동소문 안에서 인력거꾼 노릇을 하는 김 첨지에게는 오래간만에도 닥친 운수 좋은 날이었다.

- 현진건 「운수 좋은 날」\end{kfEvidence}
\begin{kfChoices}
\kfChoice 희망적이고 밝은 분위기를 조성한다.
\kfChoice 앞으로 벌어질 비극적 결말을 암시한다.
\kfChoice 인물 간의 갈등이 곧 해소될 것을 보여준다.
\kfChoice 주인공의 유머러스한 성격을 부각한다.
\kfChoice 과거의 행복했던 시절을 회상하게 한다.
\end{kfChoices}
\end{kfQBlock}
\begin{kfQBlock}{05}{객관식}
\kfQStem{‘반동 인물’의 역할에 대한 설명으로 적절하지 \uline{\underline{않은}} 것은?}
\begin{kfChoices}
\kfChoice 주인공과 정면으로 대립하며 갈등의 한 축을 이루는 인물이다.
\kfChoice 주동 인물이 목표를 이루지 못하도록 끊임없이 방해하는 인물이다.
\kfChoice 사건의 긴장감과 흥미를 만들어 내는 핵심적인 역할을 한다.
\kfChoice 「흥부전」의 ‘놀부’가 가장 대표적인 반동 인물의 예이다.
\kfChoice 주인공을 도와 시련을 함께 이겨 내는 조력자 역할을 한다.
\end{kfChoices}
\end{kfQBlock}
\begin{kfQBlock}{06}{객관식}
\kfQStem{다음 <보기>의 상황은 일반적인 사건 전개 과정 중 어느 단계에 해당하는가?}
\begin{kfEvidence}춘향이 신관 사또의 수청을 거부하여 옥에 갇히고, 이몽룡은 과거에 급제하여 암행어사가 되어 남원으로 향한다.\end{kfEvidence}
\begin{kfChoices}
\kfChoice 발단
\kfChoice 전개
\kfChoice 위기
\kfChoice 절정
\kfChoice 결말
\end{kfChoices}
\end{kfQBlock}
\begin{kfQBlock}{07}{객관식}
\kfQStem{소설 「흥부전」에서 '흥부'는 이야기의 중심에서 사건을 이끌어 간다. 이러한 인물 유형을 무엇이라고 하는가?}
\begin{kfChoices}
\kfChoice 평면적 인물
\kfChoice 입체적 인물
\kfChoice 주동 인물
\kfChoice 반동 인물
\kfChoice 전형적 인물
\end{kfChoices}
\end{kfQBlock}
\begin{kfQBlock}{08}{객관식}
\kfQStem{소설 「흥부전」에서 '놀부'는 '흥부'를 괴롭히고 대립하며 갈등을 일으킨다. 이러한 인물 유형을 무엇이라고 하는가?}
\begin{kfChoices}
\kfChoice 주동 인물
\kfChoice 반동 인물
\kfChoice 입체적 인물
\kfChoice 보조적 인물
\kfChoice 평면적 인물
\end{kfChoices}
\end{kfQBlock}
\begin{kfQBlock}{09}{객관식}
\kfQStem{소설 「춘향전」에서 '춘향'은 변 사또의 모진 고문에도 불구하고 끝까지 이 도령에 대한 지조를 지킨다. 이처럼 성격이 변하지 않는 인물 유형은 무엇인가?}
\begin{kfChoices}
\kfChoice 입체적 인물
\kfChoice 평면적 인물
\kfChoice 주동 인물
\kfChoice 반동 인물
\kfChoice 사회적 인물
\end{kfChoices}
\end{kfQBlock}
\begin{kfQBlock}{10}{객관식}
\kfQStem{<보기>의 인물 성격 제시 방법에 대한 설명으로 적절하지 \uline{\underline{않은}} 것은?}
\begin{kfEvidence}그리고 뭣에 떠다밀렸는지 나의 어깨를 짚은 채 그대로 퍽 쓰러진다. 그 바람에 나의 몸뚱이도 겹쳐서 쓰러지며, 한창 피어 퍼드러진 노란 동백꽃 속으로 폭 파묻혀 버렸다.

* 김유정 「동백꽃」\end{kfEvidence}
\begin{kfChoices}
\kfChoice 인물의 행동과 상황 묘사를 통해 성격을 짐작하게 한다.
\kfChoice 서술자가 인물의 성격을 직접 말하지 않고 ‘보여 주는’ 방식이다.
\kfChoice 독자에게 해석의 여지를 주는 ‘간접적 제시’에 해당한다.
\kfChoice 인물의 행동·말투에서 성격을 추측하도록 유도한다.
\kfChoice 서술자가 ‘점순이는 짓궂다’와 같이 성격을 직접 설명하는 방법이다.
\end{kfChoices}
\end{kfQBlock}
\begin{kfQBlock}{11}{객관식}
\kfQStem{「홍길동전」에서 ‘홍길동’이 서자라는 신분 때문에 겪는 갈등에 대한 설명으로 적절하지 \uline{\underline{않은}} 것은?}
\begin{kfChoices}
\kfChoice 신분 제도라는 사회 제약과 부딪히는 ‘인물과 사회의 갈등’이다.
\kfChoice 개인이 시대의 제도와 충돌하면서 발생한 갈등이다.
\kfChoice 작가는 이 갈등을 통해 신분 제도의 부조리를 비판하고자 한다.
\kfChoice 홍길동이 율도국으로 떠나는 결말과도 연결되는 갈등이다.
\kfChoice 홍길동 한 사람의 마음속에서만 일어나는 ‘인물 내면의 갈등’이다.
\end{kfChoices}
\end{kfQBlock}
\begin{kfQBlock}{12}{객관식}
\kfQStem{황순원의 「소나기」에서 '소나기'는 소년과 소녀가 수숫단에 함께 숨으며 가까워지는 계기를 만든다. 이는 배경의 어떤 역할에 해당하는가?}
\begin{kfChoices}
\kfChoice 인물의 심리 상태를 반영한다.
\kfChoice 사건 전개의 동기나 원인이 된다.
\kfChoice 인물의 성격 형성에 영향을 준다.
\kfChoice 작품의 주제를 효과적으로 드러낸다.
\kfChoice 사회적 상황을 암시한다.
\end{kfChoices}
\end{kfQBlock}
\begin{kfQBlock}{13}{단답형}
\kfQStem{소설 속에서 행동하고 생각하며 감정을 느끼는 주체를 무엇이라고 하는가?}
\kfAnswerLines{1}
\end{kfQBlock}
\begin{kfQBlock}{14}{단답형}
\kfQStem{소설에서 사건을 이끌어가는 중심인물을 무엇이라고 하는가?}
\kfAnswerLines{1}
\end{kfQBlock}
\begin{kfQBlock}{15}{단답형}
\kfQStem{주동 인물과 대립하거나 갈등을 일으키는 인물을 무엇이라고 하는가?}
\kfAnswerLines{1}
\end{kfQBlock}
\begin{kfQBlock}{16}{단답형}
\kfQStem{사건을 겪으면서 성격이 변화하고 성장하는 복잡한 인물을 무엇이라고 하는가?}
\kfAnswerLines{1}
\end{kfQBlock}
\begin{kfQBlock}{17}{단답형}
\kfQStem{서술자가 인물의 성격을 직접 설명하는 제시 방법을 무엇이라고 하는가?}
\kfAnswerLines{1}
\end{kfQBlock}
\begin{kfQBlock}{18}{단답형}
\kfQStem{소설 속에서 일어나는 일련의 사건들이 연결된 줄거리를 무엇이라고 하는가?}
\kfAnswerLines{1}
\end{kfQBlock}
\begin{kfQBlock}{19}{단답형}
\kfQStem{소설의 5단계 구성 중 갈등이 최고조에 달하며 이야기의 방향이 결정되는 단계는 무엇인가?}
\kfAnswerLines{1}
\end{kfQBlock}
\begin{kfQBlock}{20}{서술형}
\kfQStem{<보기>는 「흥부전」의 한 부분이다. '놀부'의 성격을 제시하는 방식이 무엇에 해당하는지 <조건>에 맞게 서술하시오.}
\begin{kfEvidence}놀부는 마음보가 사나웠다. 부모에게 물려받은 재산을 독차지하고도 아우를 내쫓았으니, 그는 욕심이 많고 인색하기 짝이 없는 사람이었다.\end{kfEvidence}
\kfAnswerNote{답안}{4}
\end{kfQBlock}
\kfEndMC


\kfStartLit{「꺼삐딴 리」 — 전광용}


\kfPassageNote{꺼삐딴 리}{%
[앞부분 줄거리] 이인국 박사는 뛰어난 실력을 가졌지만, 시대의 권력에 빌붙어 부와 명예를 쌓아온 기회주의적인 외과 의사이다. 어느 날 미국 유학 중인 딸이 외국인과 결혼하겠다는 소식을 전해오자, 그는 자신의 이익을 따져 이를 받아들인다. 미국으로 떠날 준비를 하던 그는 도움을 준 미국 대사관의 브라운 씨를 만나러 가던 중, 해방 직후 자신이 겪었던 위험한 순간을 떠올린다.

　1945년 8월 말, 해방의 기쁨이 온 나라를 뒤덮고 있을 때였다. 말복도 지났지만 날씨는 여전히 무더웠다. 이인국 박사는 며칠 동안 불안하고 초조한 마음에 잠도 제대로 자지 못했다. 무엇인가 좋지 않은 일이 닥쳐올 것만 같아 벌벌 떨었다. 그렇게 붐비던 환자도 찾아오지 않고 쉴 새 없던 전화도 잠잠해졌다. 입원실은 마지막 환자였던 도청의 일본인 과장이 끌려간 후 텅 비었다. 조수와 약제사는 고향에 다녀오겠다고 떠나갔고, 서울 출신인 간호사 혜숙만이 남아 빈집 같은 병원을 지키고 있었다.

　이층 일본식 방에서 얇은 옷만 입은 채 뒹굴고 있던 이인국 박사는 더위를 견디지 못하고 부채를 내던지며 일어났다. 그는 목욕탕으로 갔다. 찬물을 대야째로 머리부터 몇 번이고 끼얹었다. 등줄기가 시원해지며 몸이 가뿐해졌다. 그러나 수건으로 몸을 닦으면서도 무언가에 꽉 눌린 것처럼 가슴이 답답한 것은 사라지지 않았다.

　그는 창문으로 조심스럽게 길가를 내려다보았다. 수많은 사람들이 여전히 시끄럽게 몰려다니고 있었다. 굳게 닫힌 은행 철문에 붙은 벽보가 길 건너편에서 하얗게 보였다. 그곳에 쓰인 구절, ‘친일파, 민족 반역자를 타도하자.’ 라는 글자가 바로 눈앞에 생생하게 보이는 듯했다. 어제 저녁 그것을 처음 보았을 때의 소름 끼치는 공포가 되살아났다. 순간 이인국 박사는 방 쪽으로 머리를 홱 돌렸다. ‘나야 괜찮겠지…….’ 혼잣말을 하면서 그는 다시 부채를 들었다.

　그러나 벽보를 볼 때 자신과 눈이 마주치자, 깔보는 것인지 통쾌해하는 것인지 모를 웃음을 지으며 자신을 위아래로 훑어보던 춘석이라는 청년의 모습이 자꾸만 떠올랐다. 어두운 밤에 거미줄을 뒤집어쓴 것처럼 불쾌했다. 그까짓 놈 하고 머리에서 지워 버리려 해도 거머리처럼 끈질기게 달라붙는 것만 같았다.

　벌써 육 개월 전의 일이다. 감옥에서 병 때문에 잠시 풀려났다는 중환자가 사람에게 업혀 병원에 실려 왔다. 퀭한 눈에 뼈만 앙상하게 남은 몸을 제대로 가누지도 못하는 환자였다. 그는 간호사의 부축으로 겨우 진찰을 받았다. 이인국 박사는 청진기로 환자의 숨소리를 들으면서도, 머릿속으로는 어떻게 할지 고민하고 있었다. 입원시킬 것인가, 거절할 것인가……. 환자의 초라한 모습이나 옷차림을 보니 돈이 많지 않은 것은 분명했다. 하지만 그보다 더 마음에 걸리는 것이 있었다. 일본의 높은 관리들이 자기 집처럼 드나드는 이 병원에 ‘불온한 생각을 가졌다고 잡혀간 사람’을 입원시킨다는 것은, 시의원이라는 자신의 체면에도 맞지 않을뿐더러, 모범적인 일본 국민으로서 쌓아 올린 공이 하루아침에 무너지는 일이라고 생각했다. 순간 그는 이런 경우 단번에 결정하는 자신의 방식대로 즉시 결정을 내렸다. 그는 응급 치료만 해주고는 입원실이 없다는 가장 그럴듯하고 떳떳한 핑계를 대며 애원하는 환자를 돌려보냈다.

　그런데 며칠 전, 해방을 축하하는 거리 행진을 구경하러 나갔다가, 완장을 두른 젊은이와 눈이 마주쳤다. 자신을 노려보는 청년의 눈에서 불꽃이 튀는 듯한 기세를 느꼈다. 무슨 영문인지 몰라 어리둥절하던 이인국 박사는, 그가 바로 예전에 입원을 거절했던 환자 춘석이라는 것을 혜숙에게서 듣고야 슬금슬금 주위 눈치를 살피며 집으로 기어 들어왔다. 그 후 되도록 거리에 나가지 않았지만, 운 나쁘게도 어제 저녁 그 벽보 앞에서 다시 마주쳤던 것이다.

　갑자기 밖이 시끄러워졌다. 머리에 깍지를 낀 채 비스듬히 누워 어지러운 생각에 잠겨 있던 이인국 박사는 일어나 앉아 길가에 귀를 기울였다. 시끄러운 소리는 점점 커졌다. 궁금증을 견디지 못해 그는 엉거주춤 몸을 구부려 밖을 내다보았다. 길에 가득 찬 사람들은 손에 태극기와 붉은 깃발을 들고 고함을 지르고 있었다. ‘무엇일까?’ 그는 고개를 갸웃하며 다시 자리에 주저앉았다. 계단을 급히 뛰어 올라오는 발자국 소리가 들려왔다. 혜숙이었다.

　“아마 소련 군대가 들어오나 봐요. 모두들 난리가 났어요…….”

　숨을 헐떡이며 이야기하는 혜숙의 말에 이인국 박사는 아무 대꾸도 없이 눈만 껌벅이며 앉아 있었다. 며칠 전 라디오에서 오늘 들어올 예정이라고 했으니 이제 정말 오는구나 싶었다. 혜숙이 내려간 뒤에도 이인국 박사는 한참 동안 아무 움직임도 없이 바깥쪽만 내다보았다. 그러다 무엇을 생각했는지 자리에서 벌떡 일어나 벽장 문을 열었다. 안쪽에서 액자들을 꺼냈다. ‘일본어만 쓰는 집[國語常用의 家]’ 해방되던 날 떼어서 넣어 두고는 깜박 잊고 있었다. 그는 액자 뒤를 열어 두꺼운 종이를 빼내어, 글자 하나 남지 않도록 손끝에 힘을 주어 꼼꼼히 찢었다.

　이 종이 한 장만으로도 일본인들과 어울릴 때 얼마나 당당할 수 있었던가. 아쉬운 마음이 번개처럼 머릿속을 스쳐갔다. 그는 병원에서나 집에서나 일본말만 써왔다. 해방 뒤 어쩔 수 없이 쓰는 우리말이 오히려 어색하게 느껴질 정도였다. 아내가 앞장서고 아이들까지 잘 따라준 덕분에 이 상을 받은 것이 아니던가. 상을 받던 날, 온 집안이 큰 경사가 난 것처럼 기뻐했다. “잠꼬대까지 일본어로 할 정도가 아니면 이 영광스러운 상을 받을 수 있겠소.” 하던 일본인 관리의 웃음 띤 칭찬 소리가 떠올랐다. 그는 아이들을 일본 학교에 보낸 것을 얼마나 다행으로 여겼던가. 이인국 박사는 깊은 한숨을 내쉬었다.

　그리고는 자신의 통장에 있던 돈을 전부 내어준 은행 지점장의 친절에 새삼 고마움을 느꼈다. 그 돈마저 없었더라면…… 등골이 오싹해졌다. 무슨 세상이 오든 돈만 있으면 굶어 죽지는 않겠지. 그는 작은 금고가 있는 안방의 장롱을 생각하며 혼자 중얼거렸다. 이인국 박사는 무슨 일이 일어나도 자신만은 꼭 살아남을 것이라는 막연한 기대를 되새겼다.

　주위가 어두워졌다. 땅이 흔들리는 듯한 소리와 함께 공포가 가까워졌다. 사람들의 환호성이 터져 나왔다. 만세 소리가 계속 이어졌다. 세상 돌아가는 형편을 알아보려고 거리에 나갔던 아내가 돌아왔다.

　“여보, 탱크 부대가 들어왔어요. 거리는 온통 사람들로 난리가 났는데, 집안에 틀어박혀 뭘 하고 있어요…….”

　아내의 목소리는 흥분해 있었지만, 그것이 기쁨 때문인지 당황해서인지 알 수 없었다. ‘여자란 저렇게 어리석으면서도 대담한 것일까…….’ 이인국 박사는 엷은 어둠 속에서 아내를 지켜보며 입맛을 다셨다.

　“불도 안 켜고.” 아내가 전등 스위치를 켰다. 이인국 박사는 백 촉짜리 전구가 너무 밝아 못마땅했다.

　“불은 왜 켜는 거요?”

　“그럼 켜지 않고 캄캄하게 있어요? 자, 어서 나가 봅시다.”

　아내가 이끄는 대로 이인국 박사는 마지못해 아무렇지 않은 척하고 따라 나섰다. 눈부신 헤드라이트 불빛. 탱크 부대의 행렬은 끝없이 이어지고 있었다. 이인국 박사는 강한 불빛을 피하며 가로수에 기대어 섰다. 박수와 환호성, 만세 소리가 그치지 않는 길 양쪽을 따라 탱크는 물처럼 서서히 흘러갔다. 탱크 위로 몸을 반쯤 내민, 머리를 짧게 깎은 병사들은 가끔 ‘우라!’ 하고 외치며 손을 흔들었다. 이인국 박사는 자신과 아무 관련 없는 외국 군대라는 착각을 느끼며, 박수도 환호도 나오지 않는 어색한 마음으로 멍하니 쳐다보고만 있었다. 누가 자신의 행동을 지켜보지나 않나 싶어 주위를 두리번거렸다. 그러나 아무도 그에게는 관심을 두지 않고 탱크를 향해 목청이 터져라 만세만 부르고 있었다. ‘어떻게든 되겠지…….’ 그는 밑도 끝도 없는 한마디를 되뇌면서 조용히 집으로 들어왔다.

[중략 줄거리] 결국 이인국 박사는 친일파로 지목받아 소련군에게 붙잡힌다. 그는 감옥에서 의사로서의 기술을 이용해 살아날 기회를 잡는다. 소련군 장교 스텐코프의 혹을 제거하는 수술에 성공하여 그의 믿음을 얻고 풀려난 이인국은, 이제 소련말을 배우며 새로운 권력에 빌붙는다. 이후 6·25 전쟁이 터지고 남쪽으로 내려온 그는 서울에서 다시 병원을 차려 큰 성공을 거둔다.

　이제 이인국 박사는 미국이 새로운 기회의 땅이라고 판단하고 미국 대사관의 브라운 씨에게 접근한다. 차가 브라운 씨의 집 앞에 닿았다. ㉠\uline{미국 국기를 보면서} 이인국 박사는 과거 보았던 붉은 깃발을 떠올렸다. 브라운 씨가 나오자 이인국 박사는 웃으며 고려청자 화병을 선물로 내놓았다. 포장을 푼 브라운 씨는 기쁨을 참지 못하는 듯 ‘탱큐’를 거듭 외쳤다.

　“참 이거 귀중한 것입니다.”

　“뭐 대단한 것이 아닙니다만 그저 제 성의입니다.”

　이인국 박사는 마음이 놓이면서 만족감을 느꼈다. 다음 날 함께 사냥을 가기로 약속하고 브라운 씨의 대문을 나섰다. 그의 마음속에는 새로운 다짐과 희망이 부풀어올랐다.

　‘흥, 그 섬 오랑캐한테서도 살았고, 닥싸귀 같은 로스케한테서도 살아났는데, 미국놈이라고 다를까. 혁명이 일어나면 일어나고, 나라가 바뀌면 바뀌고, ㉡\uline{아직 이 이인국의 살 구멍은 막히지 않았다.}’

　그는 허공을 향하여 마음껏 소리치고 싶었다.

　‘그러면 우선 비행기 회사에 들러 형편이나 알아볼까.’

　이인국 박사는 시가를 비스듬히 문 채 지나가는 택시를 불러 세웠다.

“반도 호텔로.”

　차창 너머로 보이는 맑은 가을 하늘이 이인국 박사에게는 더욱 푸르고 드높게만 느껴졌다.

---
}


\kfPassageNote{문장 독해 — 발췌 문장}{%
1. 그곳에 쓰인 구절, ‘친일파, 민족 반역자를 타도하자.’ 라는 글자가 바로 눈앞에 생생하게 보이는 듯했다.

2. 어제 저녁 그것을 처음 보았을 때의 소름 끼치는 공포가 되살아났다.

3. 그까짓 놈 하고 머리에서 지워 버리려 해도 거머리처럼 끈질기게 달라붙는 것만 같았다.

4. 그는 간호사의 부축으로 겨우 진찰을 받았다. 이인국 박사는 청진기로 환자의 숨소리를 들으면서도, 머릿속으로는 어떻게 할지 고민하고 있었다.

5. "불온한 생각을 가졌다고 잡혀간 사람"

6. 아내가 앞장서고 아이들까지 잘 따라준 덕분에 이 상을 받은 것이 아니던가.

7. 이 종이 한 장만으로도 일본인들과 어울릴 때 얼마나 당당할 수 있었던가.

8. "참 이거 귀중한 것입니다."

9. '흥, 그 섬 오랑캐한테서도 살았고, 닥싸귀 같은 로스케한테서도 살아났는데, 미국놈이라고 다를까.'
}


\kfPassageNote{해설 학습 — 본문}{%
[단원 1] 전광용의 「꺼삐딴 리」를 이해하기 위해서는 작품의 시대적 배경을 파악하는 것이 중요하다. 이 소설은 일제강점기 말부터 1945년 해방, 남북 분단과 소련군 진주, 그리고 6.25 전쟁을 거쳐 1950년대에 이르는 한국 현대사의 가장 혼란한 시기를 배경으로 한다. 권력의 주인이 계속해서 바뀌는 격동의 시기였다.

주인공 이인국 박사는 이러한 시대적 혼란을 오히려 새로운 기회로 삼는 기회주의자의 전형이다.  그는 의사로서의 신념이나 민족에 대한 의리보다 오직 자신의 생존과 이익을 최우선으로 삼는 카멜레온 같은 인물이다. 일제강점기에는 ‘일본어만 쓰는 집’이라는 상을 받을 정도의 친일파로, 해방 후 북한에서는 재빨리 러시아어를 배워 친소파로, 6.25 전쟁 이후 남한에서는 미군에게 줄을 대는 친미파로 변신하며 새로운 권력에 의지하며 살아간다.

작가는 이인국 박사라는 인물을 통해 우리의 아픈 역사를 되돌아보게 한다. 나라와 민족이 위기에 처했을 때 지식인으로서 책임을 다하지 않고 개인의 이익만을 좇았던 인물을 풍자하고 비판하는 것이다. 결국 이인국 박사는 올바른 가치관 없이 살아남는 기술만을 터득한 인물이 어떻게 비인간적으로 변해가는지를 보여주는 시대의 거울이라 할 수 있다.

[단원 2] 소설에서 소재는 단순한 배경이 아니라 인물의 성격이나 주제 의식을 드러내는 중요한 역할을 한다. 「꺼삐딴 리」에는 주인공 이인국 박사의 기회주의적인 성격을 상징적으로 보여주는 여러 소재들이 등장한다.

첫째, '일본어만 쓰는 집'이라고 쓰인 액자이다. 이것은 이인국 박사가 얼마나 열심히 일본에 협력했는지를 보여주는 증거물이다. 그는 해방이 되자마자 이것을 찢어버리며 과거를 손쉽게 지워버렸다. 둘째, 소련군 장교 스텐코프의 턱에 난 혹이다. 이 혹은 이인국 박사에게는 절망적인 상황에서 벗어날 수 있는 기회가 되었다. 그는 자신의 의술을 이용해 혹을 제거해주고 새로운 권력의 편에 서게 되었다. 셋째, 미국인 브라운 씨에게 선물한 고려청자이다. 고려청자는 우리 민족의 소중한 문화유산이지만, 이인국 박사에게는 자신의 목적을 달성하기 위한 뇌물에 불과했다.

이처럼 작가는 여러 소재를 효과적으로 활용하여 이인국 박사의 비도덕적이고 기회주의적인 성격을 뚜렷하게 보여준다. 이를 통해 독자들은 인물의 내면을 더 깊이 이해할 수 있게 된다.

[단원 3] 소설의 제목은 작품의 주제 의식을 함축적으로 담고 있는 경우가 많다. 「꺼삐딴 리」라는 독특한 제목에는 작가의 깊은 의도가 숨어 있다.

'꺼삐딴 리'는 '카피탄 리'를 우리말로 표기한 것이다. '카피탄'은 러시아어로 우두머리를 뜻하는 말이다. 이인국 박사가 소련군 장교의 혹을 떼어준 후, 그를 신뢰하게 된 스텐코프가 이인국 박사를 '카피탄 리'라고 부르면서 얻게 된 별명이다. 일제강점기에는 일본식 이름으로 불렸을 그가 이제는 러시아식 별명으로 불리게 된 것이다. 이는 그가 얼마나 성공적으로 새로운 시대에 적응했는지를 보여준다. 소설의 마지막, 모든 일이 자신의 뜻대로 풀린 이인국 박사가 보는 하늘은 '더욱 푸르고 드높게만' 느껴진다. 나라를 배신하고 자신의 이익만을 챙긴 인물이 아무런 죄책감 없이 희망찬 미래를 꿈꾸는 이 장면은 매우 반어적이다.

작가는 이러한 제목과 결말을 통해 잘못된 시대 속에서 양심을 잃어버린 인물의 모습을 보여주며, 진정으로 올바른 삶이 무엇인지에 대한 질문을 독자에게 던진다.
}


\kfActivityBg \par\needspace{15\baselineskip}
\kfSecHeader{kfAreaLiterature}{활동}{문장 독해}
\par\kfMaybeNeedspace{32\baselineskip}\noindent{\kfTipMark\,\,}{\small\color{kfInk} 다음은 작품에서 발췌한 문장입니다. 각 문장을 읽고 물음에 답하시오.}\par\nopagebreak[4]\smallskip\nopagebreak[4]
\kfBeginMC
\begin{kfSentenceUnit}{1}{}
\kfSentQ{1.}{다음 문장에서 '그곳'이 가리키는 장소는 어디인가?}
\begin{kfChoices}
\kfChoice 벽시계
\kfChoice 작은 손거울
\kfChoice 은행 철문에 붙은 벽보
\kfChoice 소련군이 붙인 선전 포스터
\end{kfChoices}
\end{kfSentenceUnit}

\begin{kfSentenceUnit}{2}{}
\kfSentQ{2.}{다음 문장에서 '그것'이 가리키는 것은 무엇인가?}
\begin{kfChoices}
\kfChoice 벽보의 글자
\kfChoice 춘석의 비웃는 얼굴
\kfChoice 딸이 보낸 편지 내용
\kfChoice 소련군 탱크의 소리
\end{kfChoices}
\end{kfSentenceUnit}

\begin{kfSentenceUnit}{3}{}
\kfSentQ{3.}{다음 문장에서 '그까짓 놈'은 누구를 가리키는가?}
\begin{kfChoices}
\kfChoice 아내
\kfChoice 청년 춘석
\kfChoice 소련군 장교
\kfChoice 이인국 박사
\end{kfChoices}
\end{kfSentenceUnit}

\begin{kfSentenceUnit}{4}{}
\kfSentQ{4.}{다음 문장에서 서술어의 주어로 가장 적절한 것은?

* 서술어 '받았다'의 주어 * 서술어 '고민하고 있었다'의 주어}
\begin{kfChoices}
\kfChoice 혜숙, 아내
\kfChoice 이인국, 브라운
\kfChoice 스텐코프, 춘석
\kfChoice 춘석, 이인국 박사
\end{kfChoices}
\end{kfSentenceUnit}

\begin{kfSentenceUnit}{5}{}
\kfSentQ{5.}{다음 문장에서 '불온한 생각을 가졌다고 잡혀간 사람'은 누구를 가리키는가?}
\begin{kfChoices}
\kfChoice 혜숙
\kfChoice 아내
\kfChoice 환자 춘석
\kfChoice 이인국 박사
\end{kfChoices}
\end{kfSentenceUnit}

\begin{kfSentenceUnit}{6}{}
\kfSentQ{6.}{다음 문장에서 서술어의 주어로 가장 적절한 것은?

* 서술어 '앞장서고'의 주어 * 서술어 '따라준'의 주어}
\begin{kfChoices}
\kfChoice 혜숙, 춘석
\kfChoice 혜숙, 아내
\kfChoice 이인국, 아내
\kfChoice 아내, 아이들
\end{kfChoices}
\end{kfSentenceUnit}

\begin{kfSentenceUnit}{7}{}
\kfSentQ{7.}{다음 문장에서 '이 종이'가 가리키는 것은 무엇인가?}
\begin{kfChoices}
\kfChoice ‘국어 상용의 가’ 상장
\kfChoice 신문 기사
\kfChoice 예금 통장
\kfChoice 진료 기록부
\end{kfChoices}
\end{kfSentenceUnit}

\begin{kfSentenceUnit}{8}{}
\kfSentQ{8.}{다음 대화에서 브라운 씨가 말하는 '이거'가 가리키는 물건은 무엇인가?}
\begin{kfChoices}
\kfChoice 엽총
\kfChoice 고려청자 화병
\kfChoice 브라운의 사냥총
\kfChoice 이인국 박사의 수술 도구
\end{kfChoices}
\end{kfSentenceUnit}

\begin{kfSentenceUnit}{9}{}
\kfSentQ{9.}{다음은 이인국 박사의 독백이다. 생략된 주어를 포함하여 서술어 '살았고'와 '살아났는데'의 주어는 누구인가?}
\begin{kfChoices}
\kfChoice 혜숙
\kfChoice 브라운
\kfChoice 이인국 박사
\kfChoice 소련군 장교 스텐코프
\end{kfChoices}
\end{kfSentenceUnit}

\kfEndMC

\kfActivityBg \kfSecHeader{kfAreaLiterature}{활동}{해설 문제}
\par\kfMaybeNeedspace{18\baselineskip}\noindent{\kfTipMark\,\,}{\small\color{kfInk} 해설 본문을 읽고, 단원별 단답형 물음에 답하시오.}\par\nopagebreak[4]\smallskip\nopagebreak[4]
\kfBeginMC
\begin{enumerate}[leftmargin=22pt, itemsep=6pt, topsep=4pt, label=\textbf{\arabic*.}]
\item 이 소설의 주인공 이름은 무엇인가?
\item 시대의 변화에 따라 유리한 쪽을 따라다니는 사람을 무엇이라고 하는가?
\item 이 소설의 배경이 되는 시기는 언제부터 언제까지인가?
\item 이인국 박사가 의사로서의 신념이나 민족에 대한 의리보다 최우선으로 생각한 것은 무엇인가?
\item 일제 강점기 시절 이인국 박사의 삶의 태도를 한 단어로 무엇이라고 하는가?
\item 이인국 박사가 해방 후 새로운 권력으로 삼은 대상은 누구인가?
\item 6.25 전쟁 이후 우리나라에 큰 영향력을 미친 나라는 어디인가?
\item 작가는 이인국 박사라는 인물을 어떤 시선으로 바라보고 있는가?
\item 이인국 박사의 적극적인 친일 행위를 증명하는 소품은 무엇인가?
\item 이인국 박사가 위기에서 벗어날 기회가 되었던 것은 무엇인가?
\item 이인국 박사가 미국인에게 뇌물로 사용한 우리 문화유산은 무엇인가?
\item 이인국 박사에게 고려청자는 문화유산이 아닌 무엇으로 쓰였는가?
\item 이인국 박사가 액자를 찢어버리는 행동은 무엇을 의미하는가?
\item 소설에서 '소재'는 어떤 역할을 하는가?
\item '카피탄'은 어느 나라 말인가?
\item 이인국 박사에게 '카피탄 리'라는 별명을 붙여준 사람은 누구인가?
\item 이 별명은 이인국 박사가 어디에 성공적으로 적응했음을 보여주는가?
\item 반도덕적인 인물이 희망찬 미래를 꿈꾸는 마지막 장면의 표현 방식을 무엇이라고 하는가?
\item 작가는 반어적인 결말을 통해 이인국 박사의 어떤 모습을 강조하는가?
\item 작가는 제목과 결말을 통해 독자에게 무엇을 던지는가?
\end{enumerate}
\kfEndMC

\kfSecHeader{kfAreaLiterature}{문제}{}

\par\medskip
\kfBeginMC
\begin{kfQBlock}{01}{객관식}
\kfQStem{이 글의 서술상 특징에 대한 설명으로 적절하지 \uline{\underline{않은}} 것은?}
\begin{kfChoices}
\kfChoice 현재에서 과거를 회상하고 다시 현재로 돌아오는 역순행적 구성을 취한다.
\kfChoice 미국 대사관으로 가는 길이라는 현재 상황 속에서 과거가 끼어든다.
\kfChoice 시간 순서를 의도적으로 비틀어 인물의 일대기를 압축적으로 보여 준다.
\kfChoice 회상 장면을 통해 인물의 변신 과정을 효과적으로 드러낸다.
\kfChoice 여러 인물의 시점을 번갈아 사용해 사건을 입체적으로 서술한다.
\end{kfChoices}
\end{kfQBlock}
\begin{kfQBlock}{02}{객관식}
\kfQStem{이인국 박사가 춘석의 입원을 거절한 이유로 가장 적절한 것은?}
\begin{kfChoices}
\kfChoice 병원에 빈 입원실이 없었기 때문에
\kfChoice 춘석의 병이 너무 심해서 고칠 자신이 없었기 때문에
\kfChoice 치료비를 낼 돈이 없어 보이는 춘석이 마음에 들지 않아서
\kfChoice 독립운동을 하던 춘석을 입원시켰다가 자신에게 피해가 올까 봐
\kfChoice 춘석이 과거 자신에게 무례하게 굴었던 것에 복수하기 위해서
\end{kfChoices}
\end{kfQBlock}
\begin{kfQBlock}{03}{객관식}
\kfQStem{이인국 박사가 겪은 시대의 순서가 바르게 연결된 것은?}
\begin{kfChoices}
\kfChoice 미군 주둔기 → 일제 강점기 → 소련군 주둔기
\kfChoice 일제 강점기 → 소련군 주둔기 → 미군 주둔기
\kfChoice 소련군 주둔기 → 미군 주둔기 → 일제 강점기
\kfChoice 일제 강점기 → 미군 주둔기 → 소련군 주둔기
\kfChoice 미군 주둔기 → 소련군 주둔기 → 일제 강점기
\end{kfChoices}
\end{kfQBlock}
\begin{kfQBlock}{04}{객관식}
\kfQStem{이인국 박사에 대한 설명으로 적절하지 \uline{\underline{않은}} 것은?}
\begin{kfChoices}
\kfChoice 자신의 이익을 위해서라면 신념도 쉽게 바꾼다.
\kfChoice 시대의 변화를 빠르게 파악하고 적응하는 능력이 뛰어나다.
\kfChoice 과거의 잘못에 대해 양심의 가책을 느끼며 괴로워한다.
\kfChoice 의술을 자신의 생존과 출세를 위한 도구로 활용한다.
\kfChoice 새로운 권력자에게 접근하기 위해 뇌물을 사용한다.
\end{kfChoices}
\end{kfQBlock}
\begin{kfQBlock}{05}{객관식}
\kfQStem{<보기>의 관점에서 ㉠의 의미를 해석한 것으로 가장 적절한 것은?}
\begin{kfEvidence}소설에서 국기는 흔히 그 나라의 주권이나 권력을 상징한다. 따라서 인물이 어떤 국기를 바라보는지는 그가 어떤 권력에 의지하고 있는지를 보여주는 장치가 될 수 있다.\end{kfEvidence}
\begin{kfChoices}
\kfChoice 나라를 잃었던 그 슬픈 과거를 떠올리게 하는 매개체
\kfChoice 앞으로 자신이 새롭게 의지하고 복종해야 할 권력의 상징
\kfChoice 민족의 독립을 위해 싸웠던 투사들의 숭고한 희생
\kfChoice 전 세계의 평화와 화합을 바라는 인류 보편의 소망
\kfChoice 자신이 저지른 과거의 잘못을 반성하게 하는 계기
\end{kfChoices}
\end{kfQBlock}
\begin{kfQBlock}{06}{객관식}
\kfQStem{㉡에 담긴 이인국 박사의 생각으로 가장 적절한 것은?}
\begin{kfChoices}
\kfChoice 어떤 시련이 닥쳐도 정의는 반드시 승리할 것이다.
\kfChoice 의사로서의 실력만 있다면 어디서든 성공할 수 있다.
\kfChoice 시대가 어떻게 변하든 나는 살아남을 방법을 찾을 수 있다.
\kfChoice 진정한 친구 한 명만 있다면 인생은 실패하지 않는다.
\kfChoice 신념을 지키다 보면 언젠가 세상이 알아줄 날이 올 것이다.
\end{kfChoices}
\end{kfQBlock}
\begin{kfQBlock}{07}{객관식}
\kfQStem{이 글에 드러난 이인국 박사의 말과 행동으로 적절하지 \uline{\underline{않은}} 것은?}
\begin{kfChoices}
\kfChoice 해방이 되자 ‘일본어만 쓰는 집’ 액자를 찢어 버렸다.
\kfChoice 소련군 장교의 혹을 제거해주고 그의 신임을 얻었다.
\kfChoice 자신의 잘못을 비난하는 춘석에게 진심으로 사과했다.
\kfChoice 미국인 관리에게 고려청자를 선물하며 환심을 사려 했다.
\kfChoice 시대가 바뀌어도 자신은 살아남을 수 있다고 확신했다.
\end{kfChoices}
\end{kfQBlock}
\begin{kfQBlock}{08}{객관식}
\kfQStem{이인국 박사의 심리 변화를 바르게 제시한 것은?}
\begin{kfChoices}
\kfChoice 불안과 초조 → 안도와 만족감 → 새로운 희망
\kfChoice 분노와 증오 → 연민과 동정 → 체념과 절망
\kfChoice 기대와 설렘 → 깊은 실망 → 쓸쓸함과 후회
\kfChoice 자부심과 긍지 → 수치심과 죄책감 → 반성과 다짐
\kfChoice 평화와 안정 → 혼란과 갈등 → 비관적 전망
\end{kfChoices}
\end{kfQBlock}
\begin{kfQBlock}{09}{객관식}
\kfQStem{소설의 제목이 ‘꺼삐딴 리’인 이유에 대한 설명으로 적절하지 \uline{\underline{않은}} 것은?}
\begin{kfChoices}
\kfChoice ‘꺼삐딴’은 우두머리·대장을 뜻하는 러시아어 단어이다.
\kfChoice 소련군 점령 시기에 이인국 박사가 얻은 별명에서 따온 것이다.
\kfChoice 낯선 외래어 별명을 제목으로 삼아 인물에 대한 거리감을 만들어 낸다.
\kfChoice 권력에 기생하며 변신을 거듭해 온 인물의 모습을 풍자하기 위해서이다.
\kfChoice 주인공이 의술이 뛰어난 ‘최고의 의사’임을 칭찬하기 위해서이다.
\end{kfChoices}
\end{kfQBlock}
\begin{kfQBlock}{10}{객관식}
\kfQStem{이 소설을 통해 작가가 말하고자 하는 바에 대한 설명으로 적절하지 \uline{\underline{않은}} 것은?}
\begin{kfChoices}
\kfChoice 자신의 이익만 좇는 기회주의적 삶에 대한 비판이다.
\kfChoice 역사의 위기 속에서도 변하지 않는 양심의 가치를 강조한다.
\kfChoice 권력에 기생하는 인물의 모습을 통해 시대의 부조리를 풍자한다.
\kfChoice 부정적인 인물을 처벌 없이 그려 냄으로써 오히려 비판의 효과를 높인다.
\kfChoice 격동의 시대에는 이인국 박사처럼 사는 것이 현명하다는 점을 강조한다.
\end{kfChoices}
\end{kfQBlock}
\begin{kfQBlock}{11}{서술형}
\kfQStem{이인국 박사가 환자 '춘석'의 입원을 거절한 이유와, 그 결정을 통해 알 수 있는 그의 성격을 < 조 건 >에 맞게 서술하시오.}
\kfAnswerNote{답안}{4}
\end{kfQBlock}
\begin{kfQBlock}{12}{서술형}
\kfQStem{이인국 박사가 시대의 변화에 따라 어떻게 적응해 나갔는지, 그가 의지했던 ‘세력’의 변화를 중심으로 < 조 건 >에 맞게 서술하시오.}
\kfAnswerNote{답안}{4}
\end{kfQBlock}
\kfEndMC


\kfStartNonlit{[기술] 말소리를 글자로 바꾸는 컴퓨터의 비밀}


\kfPassageNote{[기술] 말소리를 글자로 바꾸는 컴퓨터의 비밀}{%
　음성 인식 기술은 컴퓨터가 사람이 말하는 소리를 듣고 그것을 글자로 바꾸는 기술이다. 우리가 스마트폰에 "날씨 알려줘"라고 말하면 컴퓨터가 이를 인식해서 날씨 정보를 보여주는 것이 바로 이 기술 덕분이다.

　사람의 말은 'ㄱ', 'ㅏ', 'ㄴ' 같은 음소들이 시간 순서대로 이어진 것으로 볼 수 있다. 음소란 말의 가장 작은 소리 단위로, 자음과 모음을 말한다. 컴퓨터는 각 단어의 음소 배열을 '기준 패턴'으로 미리 저장해 두고, 사람이 말한 소리에서 뽑아낸 '입력 패턴'과 비교해서 단어를 알아낸다.

　컴퓨터가 음성을 인식하는 과정은 다음과 같다. 먼저 입력된 소리에서 잡음을 제거하고 음성 신호만 골라낸다. 그 다음 음성 신호를 하나의 음소로 판단되는 구간인 '음소 추정 구간'들로 나눈다. 하지만 음성을 정확히 음소 단위로 나누기는 어렵다. 그래서 먼저 음성을 일정한 시간 간격의 '단위 구간'으로 나누고, 이 구간들을 하나씩 또는 여러 개를 이어 붙여서 음소 추정 구간을 만든다.

　각 음소 추정 구간에서는 '특징 벡터'라는 것을 뽑아낸다. 특징 벡터는 음소를 구별하는 데 필요한 정보를 숫자로 나타낸 것이다. 구간의 길이와 상관없이 각 구간에서 1개씩만 뽑아낸다. 특징 벡터는 음소의 특성을 잘 보여주는 정보를 사용하지만, 사람마다 다른 목소리 특성은 제외한다. 정보의 종류가 많을수록 음소를 더 정확하게 알아낼 수 있지만, 그만큼 계산이 복잡해져서 시간이 오래 걸린다.

　음성을 인식하려면 입력 패턴과 기준 패턴의 특징 벡터를 비교해야 한다. 이때 음소 추정 구간의 개수가 비교하려는 기준 패턴의 음소 개수와 같아야 한다. 각 음소 추정 구간에서 뽑은 특징 벡터를 순서대로 배열해서 입력 패턴을 만든다.

　예를 들어 입력된 음성을 S1, S2, S3 세 개의 단위 구간으로 나눈다고 하자. 비교하려는 기준 패턴의 음소가 3개라면 [S1, S2, S3]으로 3개의 음소 추정 구간을 만든다. 기준 패턴의 음소가 2개라면 [S1, S2\textasciitilde{}S3]이나 [S1\textasciitilde{}S2, S3] 같은 방식으로 2개의 음소 추정 구간을 만든다. 이렇게 여러 개의 입력 패턴이 가능한 경우에는 모든 경우의 '패턴 거리'를 구해서 그중 가장 작은 값을 선택한다. 패턴 거리는 입력 패턴과 기준 패턴의 특징 벡터 차이를 모두 합한 값이다.

　단위 구간을 더 짧게 나누면 음소 추정 구간을 잘못 설정하는 오류를 줄일 수 있다. 하지만 계산량이 많아져서 처리 시간이 길어진다.

　이런 방식으로 컴퓨터에 저장된 모든 기준 패턴과 비교해서 패턴 거리가 가장 작은 기준 패턴을 찾는다. 최종적으로 이 기준 패턴에 해당하는 글자를 입력된 음성에 대한 인식 결과로 출력한다.
}


\kfPassageNote{문장 독해 — 본문 문장}{%
1. 음성 인식 기술은 컴퓨터가 사람이 말하는 소리를 듣고 그것을 글자로 바꾸는 기술이다.

2. 우리가 스마트폰에 "날씨 알려줘"라고 말하면 컴퓨터가 이를 인식해서 날씨 정보를 보여주는 것이 바로 이 기술 덕분이다.

3. 사람의 말은 'ㄱ', 'ㅏ', 'ㄴ' 같은 음소들이 시간 순서대로 이어진 것으로 볼 수 있다.

4. 음소란 말의 가장 작은 소리 단위로, 자음과 모음을 말한다.

5. 컴퓨터는 각 단어의 음소 배열을 '기준 패턴'으로 미리 저장해 두고, 사람이 말한 소리에서 뽑아낸 '입력 패턴'과 비교해서 단어를 알아낸다.

6. 컴퓨터가 음성을 인식하는 과정은 다음과 같다.

7. 먼저 입력된 소리에서 잡음을 제거하고 음성 신호만 골라낸다.

8. 그 다음 음성 신호를 하나의 음소로 판단되는 구간인 '음소 추정 구간'들로 나눈다.

9. 하지만 음성을 정확히 음소 단위로 나누기는 어렵다.

10. 그래서 먼저 음성을 일정한 시간 간격의 '단위 구간'으로 나누고, 이 구간들을 하나씩 또는 여러 개를 이어 붙여서 음소 추정 구간을 만든다.

11. 각 음소 추정 구간에서는 '특징 벡터'라는 것을 뽑아낸다.

12. 특징 벡터는 음소를 구별하는 데 필요한 정보를 숫자로 나타낸 것이다.

13. 구간의 길이와 상관없이 각 구간에서 1개씩만 뽑아낸다.

14. 특징 벡터는 음소의 특성을 잘 보여주는 정보를 사용하지만, 사람마다 다른 목소리 특성은 제외한다.

15. 정보의 종류가 많을수록 음소를 더 정확하게 알아낼 수 있지만, 그만큼 계산이 복잡해져서 시간이 오래 걸린다.

16. 음성을 인식하려면 입력 패턴과 기준 패턴의 특징 벡터를 비교해야 한다.

17. 이때 음소 추정 구간의 개수가 비교하려는 기준 패턴의 음소 개수와 같아야 한다.

18. 각 음소 추정 구간에서 뽑은 특징 벡터를 순서대로 배열해서 입력 패턴을 만든다.

19. 예를 들어 입력된 음성을 S1, S2, S3 세 개의 단위 구간으로 나눈다고 하자.

20. 비교하려는 기준 패턴의 음소가 3개라면 [S1, S2, S3]으로 3개의 음소 추정 구간을 만든다.

21. 기준 패턴의 음소가 2개라면 [S1, S2\textasciitilde{}S3]이나 [S1\textasciitilde{}S2, S3] 같은 방식으로 2개의 음소 추정 구간을 만든다.

22. 이렇게 여러 개의 입력 패턴이 가능한 경우에는 모든 경우의 '패턴 거리'를 구해서 그중 가장 작은 값을 선택한다.

23. 패턴 거리는 입력 패턴과 기준 패턴의 특징 벡터 차이를 모두 합한 값이다.

24. 단위 구간을 더 짧게 나누면 음소 추정 구간을 잘못 설정하는 오류를 줄일 수 있다.

25. 하지만 계산량이 많아져서 처리 시간이 길어진다.

26. 이런 방식으로 컴퓨터에 저장된 모든 기준 패턴과 비교해서 패턴 거리가 가장 작은 기준 패턴을 찾는다.

27. 최종적으로 이 기준 패턴에 해당하는 글자를 입력된 음성에 대한 인식 결과로 출력한다.
}


\kfActivityBg \par\kfMaybeNeedspace{32\baselineskip}\noindent{\kfTipMark\,\,}{\small\color{kfInk} 다음 표의 '의미'를 읽고, 그에 해당하는 '어휘'를 지문에서 찾아 적으시오.}\par\nopagebreak[4]\smallskip\nopagebreak[4]
\begin{kfVocab}
\kfVocabItem{1}{컴퓨터가 사람의 말을 듣고 글자로 바꾸는 기술}
\kfVocabItem{2}{말소리를 구별해 주는 가장 작은 소리의 단위}
\kfVocabItem{3}{컴퓨터에 미리 저장해 둔 각 단어의 음소 배열}
\kfVocabItem{4}{사람이 말한 소리에서 뽑아낸 음소 배열}
\kfVocabItem{5}{소리에서 잡음을 제거하고 남은 순수한 음성 정보}
\kfVocabItem{6}{하나의 음소로 판단되는 구간}
\kfVocabItem{7}{음성 신호를 일정한 시간 간격으로 나눈 구간}
\kfVocabItem{8}{음소의 특성을 구별하기 위해 숫자로 나타낸 정보}
\kfVocabItem{9}{입력 패턴과 기준 패턴의 특징 벡터 차이를 모두 합한 값}
\end{kfVocab}

\kfActivityBg \par\needspace{15\baselineskip}
\kfSecHeader{kfAreaNonlit}{활동}{문장 독해}
\par\kfMaybeNeedspace{32\baselineskip}\noindent{\kfTipMark\,\,}{\small\color{kfInk} 다음은 본문의 번호 매긴 문장입니다. 각 문장을 읽고 물음에 답하시오.}\par\nopagebreak[4]\smallskip\nopagebreak[4]
\kfBeginMC
\begin{kfSentenceUnit}{1}{}
\kfSentQ{1.}{'음성 인식 기술'이란 무엇인가?}
\begin{kfChoices}
\kfChoice 컴퓨터가 사람이 말하는 소리를 듣고 글자로 바꾸는 기술
\kfChoice 사람이 입력한 글자를 컴퓨터가 소리로 바꿔주는 기술
\kfChoice 소리의 파형을 분석하여 사람의 감정 상태를 파악하는 기술
\kfChoice 여러 사람이 동시에 말하는 소리를 각각 분리해 주는 기술
\end{kfChoices}
\end{kfSentenceUnit}

\begin{kfSentenceUnit}{2}{}
\kfSentQ{2.}{이 문장은 '음성 인식 기술'에 대해 어떻게 설명하고 있는가?}
\begin{kfChoices}
\kfChoice 구체적인 예시를 들어 독자의 이해를 돕고 있다.
\kfChoice 전문가들의 견해를 인용하여 기술의 중요성을 강조하고 있다.
\kfChoice 기술의 발전 과정을 시간 순서대로 나열하여 설명하고 있다.
\kfChoice 기술의 핵심 원리를 비유적인 표현을 사용하여 설명하고 있다.
\end{kfChoices}
\end{kfSentenceUnit}

\begin{kfSentenceUnit}{3}{}
\kfSentQ{3.}{이 문장에서 '이 기술'이 가리키는 것은 무엇인가?}
\begin{kfChoices}
\kfChoice 인공지능 기술
\kfChoice 날씨 예측 기술
\kfChoice 음성 인식 기술
\kfChoice 스마트폰 기술
\end{kfChoices}
\end{kfSentenceUnit}

\begin{kfSentenceUnit}{4}{}
\kfSentQ{4.}{사람의 말은 무엇이 시간 순서대로 이어진 것인가?}
\begin{kfChoices}
\kfChoice 음절
\kfChoice 단어
\kfChoice 문장
\kfChoice 음소
\end{kfChoices}
\end{kfSentenceUnit}

\begin{kfSentenceUnit}{5}{}
\kfSentQ{5.}{'음소'의 정의는 무엇인가?}
\begin{kfChoices}
\kfChoice 말의 가장 작은 소리 단위
\kfChoice 글자의 가장 작은 모양 단위
\kfChoice 문장의 가장 작은 의미 단위
\kfChoice 단어의 가장 작은 문법 단위
\end{kfChoices}
\end{kfSentenceUnit}

\begin{kfSentenceUnit}{6}{}
\kfSentQ{6.}{음소는 구체적으로 무엇과 무엇을 말하는가?}
\begin{kfChoices}
\kfChoice 억양과 성조
\kfChoice 자음과 모음
\kfChoice 단어와 문장
\kfChoice 주어와 서술어
\end{kfChoices}
\end{kfSentenceUnit}

\begin{kfSentenceUnit}{7}{}
\kfSentQ{7.}{컴퓨터가 단어를 인식하기 위해 미리 저장해 두는 음소 배열을 무엇이라고 하는가?}
\begin{kfChoices}
\kfChoice 입력 패턴
\kfChoice 인식 패턴
\kfChoice 기준 패턴
\kfChoice 음성 패턴
\end{kfChoices}
\end{kfSentenceUnit}

\begin{kfSentenceUnit}{8}{}
\kfSentQ{8.}{컴퓨터는 무엇과 무엇을 비교해서 단어를 알아내는가?}
\begin{kfChoices}
\kfChoice 사람 목소리와 기계음
\kfChoice 음성 신호와 특징 벡터
\kfChoice 단위 구간과 음소 구간
\kfChoice 기준 패턴과 입력 패턴
\end{kfChoices}
\end{kfSentenceUnit}

\begin{kfSentenceUnit}{9}{}
\kfSentQ{9.}{'입력 패턴'은 어디에서 뽑아낸 것인가?}
\begin{kfChoices}
\kfChoice 사람이 말한 소리
\kfChoice 미리 저장된 데이터
\kfChoice 컴퓨터의 내부 시계
\kfChoice 입력된 텍스트 파일
\end{kfChoices}
\end{kfSentenceUnit}

\begin{kfSentenceUnit}{10}{}
\kfSentQ{10.}{이 문장은 앞으로 어떤 내용이 이어질 것을 알려주는가?}
\begin{kfChoices}
\kfChoice 음성 인식 기술의 한계와 단점
\kfChoice 컴퓨터가 음성을 인식하는 과정
\kfChoice 음성 인식 기술의 역사적 배경
\kfChoice 음성 인식 기술의 다양한 활용
\end{kfChoices}
\end{kfSentenceUnit}

\begin{kfSentenceUnit}{11}{}
\kfSentQ{11.}{음성 인식 과정의 첫 번째 단계에서 제거하는 것은 무엇인가?}
\begin{kfChoices}
\kfChoice 울림
\kfChoice 공백
\kfChoice 잡음
\kfChoice 성조
\end{kfChoices}
\end{kfSentenceUnit}

\begin{kfSentenceUnit}{12}{}
\kfSentQ{12.}{잡음을 제거한 뒤, 컴퓨터가 최종적으로 골라내는 것은 무엇인가?}
\begin{kfChoices}
\kfChoice 문자 정보
\kfChoice 영상 신호
\kfChoice 배경 음악
\kfChoice 음성 신호
\end{kfChoices}
\end{kfSentenceUnit}

\begin{kfSentenceUnit}{13}{}
\kfSentQ{13.}{'음소 추정 구간'이란 무엇인가?}
\begin{kfChoices}
\kfChoice 음성 신호를 나눈, 하나의 음소로 판단되는 구간
\kfChoice 잡음을 제거하기 전의 전체 소리 구간
\kfChoice 단어와 단어 사이의 묵음이 있는 구간
\kfChoice 여러 개의 음소가 뭉쳐 있는 긴 구간
\end{kfChoices}
\end{kfSentenceUnit}

\begin{kfSentenceUnit}{14}{}
\kfSentQ{14.}{음성 신호를 처리할 때 어려운 점은 무엇인가?}
\begin{kfChoices}
\kfChoice 특징 벡터를 숫자로 표현하는 것
\kfChoice 음성 신호에서 잡음을 제거하는 것
\kfChoice 음성을 디지털 신호로 변환하는 것
\kfChoice 음성을 정확히 음소 단위로 나누는 것
\end{kfChoices}
\end{kfSentenceUnit}

\begin{kfSentenceUnit}{15}{}
\kfSentQ{15.}{음성을 음소 단위로 나누기 어려워서 대신 사용하는 방법은 무엇인가?}
\begin{kfChoices}
\kfChoice 음성을 단어 단위로 끊어서 인식하는 것
\kfChoice 음성을 문장 전체로 한 번에 처리하는 것
\kfChoice 음성을 일정한 시간 간격의 '단위 구간'으로 먼저 나누는 것
\kfChoice 음성의 높낮이를 분석하여 구간을 나누는 것
\end{kfChoices}
\end{kfSentenceUnit}

\begin{kfSentenceUnit}{16}{}
\kfSentQ{16.}{'음소 추정 구간'은 '단위 구간'을 어떻게 해서 만들어지는가?}
\begin{kfChoices}
\kfChoice 단위 구간을 무작위로 섞어서 재배열해서
\kfChoice 단위 구간을 더 잘게 쪼개서 평균을 내서
\kfChoice 단위 구간 중 잡음이 있는 부분을 제거해서
\kfChoice 단위 구간들을 하나씩 또는 여러 개를 이어 붙여서
\end{kfChoices}
\end{kfSentenceUnit}

\begin{kfSentenceUnit}{17}{}
\kfSentQ{17.}{'특징 벡터'는 어디에서 뽑아내는가?}
\begin{kfChoices}
\kfChoice 각 음소 추정 구간
\kfChoice 전체 음성 신호
\kfChoice 제거된 잡음 구간
\kfChoice 저장된 기준 패턴
\end{kfChoices}
\end{kfSentenceUnit}

\begin{kfSentenceUnit}{18}{}
\kfSentQ{18.}{'특징 벡터'의 정의는 무엇인가?}
\begin{kfChoices}
\kfChoice 음소의 발음 시간을 초 단위로 잰 것
\kfChoice 음소의 배열 순서를 기호로 표시한 것
\kfChoice 음소의 소리 크기를 그래프로 나타낸 것
\kfChoice 음소를 구별하는 데 필요한 정보를 숫자로 나타낸 것
\end{kfChoices}
\end{kfSentenceUnit}

\begin{kfSentenceUnit}{19}{}
\kfSentQ{19.}{특징 벡터는 각 음소 추정 구간에서 몇 개씩 뽑아내는가?}
\begin{kfChoices}
\kfChoice 2개
\kfChoice 3개
\kfChoice 1개
\kfChoice 4개
\end{kfChoices}
\end{kfSentenceUnit}

\begin{kfSentenceUnit}{20}{}
\kfSentQ{20.}{특징 벡터를 뽑아낼 때, 추출하는 개수에 영향을 주지 \uline{않는} 요인은 무엇인가?}
\begin{kfChoices}
\kfChoice 음소의 종류
\kfChoice 잡음의 크기
\kfChoice 정보의 종류
\kfChoice 구간의 길이
\end{kfChoices}
\end{kfSentenceUnit}

\begin{kfSentenceUnit}{21}{}
\kfSentQ{21.}{특징 벡터를 만들 때 포함하는 정보는 무엇인가?}
\begin{kfChoices}
\kfChoice 음소의 특성을 잘 보여주는 정보
\kfChoice 사람마다 다른 목소리의 높낮이
\kfChoice 말하는 사람의 기분이나 감정
\kfChoice 말하는 속도와 억양의 변화
\end{kfChoices}
\end{kfSentenceUnit}

\begin{kfSentenceUnit}{22}{}
\kfSentQ{22.}{특징 벡터를 만들 때 제외하는 정보는 무엇인가?}
\begin{kfChoices}
\kfChoice 자음과 모음의 구분
\kfChoice 음소의 고유한 특성
\kfChoice 발음의 정확성 여부
\kfChoice 사람마다 다른 목소리 특성
\end{kfChoices}
\end{kfSentenceUnit}

\begin{kfSentenceUnit}{23}{}
\kfSentQ{23.}{특징 벡터에 정보 종류가 많을 때의 장점은 무엇인가?}
\begin{kfChoices}
\kfChoice 계산 횟수가 줄어들어 속도가 빨라진다
\kfChoice 필요한 저장 공간을 절약할 수 있다
\kfChoice 음소를 더 정확하게 알아낼 수 있다
\kfChoice 소음이 있어도 인식률이 변하지 않는다
\end{kfChoices}
\end{kfSentenceUnit}

\begin{kfSentenceUnit}{24}{}
\kfSentQ{24.}{특징 벡터에 정보 종류가 많을 때의 단점은 무엇인가?}
\begin{kfChoices}
\kfChoice 음소 추정 구간을 나누기가 어려워진다
\kfChoice 음소 인식의 정확도가 오히려 떨어진다
\kfChoice 잡음을 제거하는 데 시간이 더 걸린다
\kfChoice 계산이 복잡해져서 시간이 오래 걸린다
\end{kfChoices}
\end{kfSentenceUnit}

\begin{kfSentenceUnit}{25}{}
\kfSentQ{25.}{음성을 인식하기 위해 비교해야 하는 것은 무엇인가?}
\begin{kfChoices}
\kfChoice 입력 패턴과 기준 패턴의 특징 벡터
\kfChoice 입력된 소리와 저장된 소리의 파형
\kfChoice 말하는 사람의 목소리와 표준 발음 목소리
\kfChoice 입력된 음성의 길이와 기준 패턴의 길이
\end{kfChoices}
\end{kfSentenceUnit}

\begin{kfSentenceUnit}{26}{}
\kfSentQ{26.}{입력 패턴과 기준 패턴의 특징 벡터를 비교하는 이유는 무엇인가?}
\begin{kfChoices}
\kfChoice 잡음을 제거하기 위해서
\kfChoice 음성을 인식하기 위해서
\kfChoice 정보 양을 줄이기 위해서
\kfChoice 단위 구간을 나누기 위해서
\end{kfChoices}
\end{kfSentenceUnit}

\begin{kfSentenceUnit}{27}{}
\kfSentQ{27.}{특징 벡터를 비교하는 단계에서 서로 개수가 같아야 하는 두 가지는 무엇인가?}
\begin{kfChoices}
\kfChoice 입력 패턴의 길이와 기준 패턴의 길이
\kfChoice 단위 구간의 개수와 음소 추정 구간의 개수
\kfChoice 음소 추정 구간의 개수와 기준 패턴의 음소 개수
\kfChoice 추출된 특징 벡터의 개수와 음소의 개수
\end{kfChoices}
\end{kfSentenceUnit}

\begin{kfSentenceUnit}{28}{}
\kfSentQ{28.}{'입력 패턴'은 어떻게 만들어지는가?}
\begin{kfChoices}
\kfChoice 음성 신호 전체를 하나의 벡터로 압축해서
\kfChoice 단위 구간을 시간 순서대로 단순히 연결해서
\kfChoice 기준 패턴의 특징 벡터를 복사하여 변형해서
\kfChoice 각 음소 추정 구간에서 뽑은 특징 벡터를 순서대로 배열해서
\end{kfChoices}
\end{kfSentenceUnit}

\begin{kfSentenceUnit}{29}{}
\kfSentQ{29.}{이 문장은 앞으로 무엇에 대해 설명할 것을 알려주는가?}
\begin{kfChoices}
\kfChoice 입력 패턴을 만드는 과정에 대한 구체적인 예시
\kfChoice 특징 벡터를 추출하는 다양한 수학적 방법
\kfChoice 음성 인식 기술이 실패하는 대표적인 사례
\kfChoice 올바른 입력 패턴과 잘못된 입력 패턴의 차이
\end{kfChoices}
\end{kfSentenceUnit}

\begin{kfSentenceUnit}{30}{}
\kfSentQ{30.}{기준 패턴의 음소가 3개일 때, 음소 추정 구간은 몇 개로 만들어야 하는가?}
\begin{kfChoices}
\kfChoice 1개
\kfChoice 3개
\kfChoice 2개
\kfChoice 4개
\end{kfChoices}
\end{kfSentenceUnit}

\begin{kfSentenceUnit}{31}{}
\kfSentQ{31.}{기준 패턴의 음소가 2개일 때, 음소 추정 구간은 몇 개로 만들어야 하는가?}
\begin{kfChoices}
\kfChoice 1개
\kfChoice 3개
\kfChoice 2개
\kfChoice 4개
\end{kfChoices}
\end{kfSentenceUnit}

\begin{kfSentenceUnit}{32}{}
\kfSentQ{32.}{단위 구간이 S1, S2, S3 세 개일 때, 음소 추정 구간을 2개로 만들기 위해 어떤 방법을 사용해야 하는가?}
\begin{kfChoices}
\kfChoice 단위 구간을 더 잘게 쪼개야 한다
\kfChoice 새로운 단위 구간을 추가해야 한다
\kfChoice 단위 구간 중 하나를 생략해야 한다
\kfChoice 일부 단위 구간들을 하나로 합쳐야 한다
\end{kfChoices}
\end{kfSentenceUnit}

\begin{kfSentenceUnit}{33}{}
\kfSentQ{33.}{입력 패턴이 여러 개일 때 어떤 값을 선택해야 하는가?}
\begin{kfChoices}
\kfChoice 패턴 거리가 가장 작은 값
\kfChoice 패턴 거리가 가장 큰 값
\kfChoice 패턴 거리의 평균값
\kfChoice 임의로 선택한 값
\end{kfChoices}
\end{kfSentenceUnit}

\begin{kfSentenceUnit}{34}{}
\kfSentQ{34.}{'패턴 거리'란 무엇인가?}
\begin{kfChoices}
\kfChoice 각 특징 벡터의 크기를 모두 곱하여 얻은 값
\kfChoice 입력 패턴과 기준 패턴의 길이 차이를 계산한 값
\kfChoice 기준 패턴의 특징 벡터 평균값에서 입력 패턴을 뺀 값
\kfChoice 입력 패턴과 기준 패턴의 특징 벡터 차이를 모두 합한 값
\end{kfChoices}
\end{kfSentenceUnit}

\begin{kfSentenceUnit}{35}{}
\kfSentQ{35.}{단위 구간을 더 짧게 나눌 때의 장점은 무엇인가?}
\begin{kfChoices}
\kfChoice 계산량이 줄어들어 처리 속도가 빨라진다
\kfChoice 필요한 저장 용량을 크게 줄일 수 있다
\kfChoice 음소 추정 구간을 잘못 설정하는 오류를 줄일 수 있다
\kfChoice 잡음을 더 효과적으로 제거할 수 있다
\end{kfChoices}
\end{kfSentenceUnit}

\begin{kfSentenceUnit}{36}{}
\kfSentQ{36.}{단위 구간을 더 짧게 나눌 때의 단점은 무엇인가?}
\begin{kfChoices}
\kfChoice 인식 결과의 정확도가 낮아질 수 있다
\kfChoice 기준 패턴을 만드는 과정이 복잡해진다
\kfChoice 음소 추정 구간을 설정하기가 어려워진다
\kfChoice 계산량이 많아져서 처리 시간이 길어진다
\end{kfChoices}
\end{kfSentenceUnit}

\begin{kfSentenceUnit}{37}{}
\kfSentQ{37.}{컴퓨터는 수많은 기준 패턴 중에서 어떤 것을 최종 후보로 찾는가?}
\begin{kfChoices}
\kfChoice 패턴 거리가 가장 작은 기준 패턴
\kfChoice 패턴 거리가 가장 큰 기준 패턴
\kfChoice 가장 최근에 저장된 기준 패턴
\kfChoice 길이가 가장 긴 기준 패턴
\end{kfChoices}
\end{kfSentenceUnit}

\begin{kfSentenceUnit}{38}{}
\kfSentQ{38.}{최종적으로 출력되는 인식 결과는 무엇인가?}
\begin{kfChoices}
\kfChoice 음소 추정 구간의 개수 정보
\kfChoice 추출된 특징 벡터의 숫자 배열
\kfChoice 입력된 음성과 가장 비슷한 소리
\kfChoice 패턴 거리가 가장 작은 기준 패턴에 해당하는 글자
\end{kfChoices}
\end{kfSentenceUnit}

\kfEndMC

\kfActivityBg \par\kfMaybeNeedspace{32\baselineskip}\noindent{\kfTipMark\,\,}{\small\color{kfInk} 글의 전체적인 논리 흐름을 파악하며 빈칸에 알맞은 내용을 채워 보세요.}\par\nopagebreak[4]\smallskip\nopagebreak[4]
\begin{kfStructure}
\kfNode{0}{}{}
\end{kfStructure}

\kfSecHeader{kfAreaNonlit}{문제}{}

\par\medskip
\kfBeginMC
\begin{kfQBlock}{01}{객관식}
\kfQStem{윗글에서 설명하는 음성 인식의 핵심 원리에 대한 설명으로 적절하지 \uline{\underline{않은}} 것은?}
\begin{kfChoices}
\kfChoice 입력된 음성 신호로부터 ‘입력 패턴’을 만들어 낸다.
\kfChoice 미리 저장해 두었던 단어별 ‘기준 패턴’과 비교한다.
\kfChoice 두 패턴 사이의 차이인 ‘패턴 거리’를 측정한다.
\kfChoice 패턴 거리가 가장 작게 나오는 단어를 찾아 인식한다.
\kfChoice 음성의 주파수를 분석해 말하는 사람의 감정을 파악한다.
\end{kfChoices}
\end{kfQBlock}
\begin{kfQBlock}{02}{객관식}
\kfQStem{윗글의 내용과 일치하지 \uline{않는} 것은?}
\begin{kfChoices}
\kfChoice 음성 신호를 정확히 음소 단위로 나누는 것은 어려운 작업이다.
\kfChoice 특징 벡터는 각 음소 추정 구간의 길이와 상관없이 1개씩 추출된다.
\kfChoice 특징 벡터에 담는 정보의 종류가 많아질수록 인식 처리 시간은 짧아진다.
\kfChoice 입력 패턴의 음소 추정 구간 개수는 기준 패턴의 음소 개수와 같아야 한다.
\kfChoice 단위 구간을 짧게 나눌수록 계산량은 많아진다.
\end{kfChoices}
\end{kfQBlock}
\begin{kfQBlock}{03}{객관식}
\kfQStem{음성 인식 과정의 일부를 순서대로 나열한 것으로 가장 적절한 것은?}
\begin{kfChoices}
\kfChoice 특징 벡터 추출 → 단위 구간 나누기 → 잡음 제거
\kfChoice 잡음 제거 → 단위 구간 나누기 → 특징 벡터 추출
\kfChoice 단위 구간 나누기 → 잡음 제거 → 특징 벡터 추출
\kfChoice 잡음 제거 → 특징 벡터 추출 → 단위 구간 나누기
\kfChoice 특징 벡터 추출 → 잡음 제거 → 단위 구간 나누기
\end{kfChoices}
\end{kfQBlock}
\begin{kfQBlock}{04}{객관식}
\kfQStem{윗글을 바탕으로 추론한 내용으로 가장 적절한 것은?}
\begin{kfChoices}
\kfChoice 음성 인식 기술의 정확도는 주변 소음의 크기와 무관하다.
\kfChoice 사람의 목소리가 아닌 다른 소리는 인식 과정에서 고려되지 않는다.
\kfChoice '패턴 거리'가 크다는 것은 두 패턴이 서로 유사하다는 의미이다.
\kfChoice 기준 패턴의 종류가 적을수록 더 많은 단어를 인식할 수 있다.
\kfChoice 음성 인식의 속도를 높이려면 단위 구간의 길이를 늘리는 것이 유리하다.
\end{kfChoices}
\end{kfQBlock}
\begin{kfQBlock}{05}{객관식}
\kfQStem{다음 <보기>와 같은 상황에서 만들어질 수 \uline{\underline{없는}} 입력 패턴은?}
\begin{kfEvidence}입력된 음성을 다섯 개의 단위 구간 [S1, S2, S3, S4, S5]로 나누었다. 이 음성을 음소가 세 개인 기준 패턴과 비교하려고 한다.\end{kfEvidence}
\begin{kfChoices}
\kfChoice S1, S2, S3\textasciitilde{}S5
\kfChoice S1, S2\textasciitilde{}S3, S4\textasciitilde{}S5
\kfChoice S1\textasciitilde{}S2, S3, S4\textasciitilde{}S5
\kfChoice S1\textasciitilde{}S3, S4, S5
\kfChoice S1\textasciitilde{}S2, S3\textasciitilde{}S5
\end{kfChoices}
\end{kfQBlock}
\begin{kfQBlock}{06}{단답형}
\kfQStem{컴퓨터가 사람이 말하는 소리를 듣고 글자로 바꾸는 기술은 무엇인가?}
\kfAnswerLines{1}
\end{kfQBlock}
\begin{kfQBlock}{07}{단답형}
\kfQStem{컴퓨터가 미리 저장해 둔 각 단어의 음소 배열을 무엇이라고 하는가?}
\kfAnswerLines{1}
\end{kfQBlock}
\begin{kfQBlock}{08}{단답형}
\kfQStem{음성 인식 과정에서 가장 먼저 처리하는 작업은 무엇인가?}
\kfAnswerLines{1}
\end{kfQBlock}
\begin{kfQBlock}{09}{단답형}
\kfQStem{음성을 일정한 시간 간격으로 나눈 구간을 무엇이라고 하는가?}
\kfAnswerLines{1}
\end{kfQBlock}
\begin{kfQBlock}{10}{단답형}
\kfQStem{음소를 구별하는 데 필요한 정보를 숫자로 나타낸 것을 무엇이라고 하는가?}
\kfAnswerLines{1}
\end{kfQBlock}
\begin{kfQBlock}{11}{서술형}
\kfQStem{컴퓨터가 여러 가능한 입력 패턴들 중에서 최종 인식 결과를 결정하는 과정을 <조건>에 맞게 서술하시오.}
\kfAnswerNote{답안}{4}
\end{kfQBlock}
\kfEndMC



\end{document}
